% Options for packages loaded elsewhere
\PassOptionsToPackage{unicode}{hyperref}
\PassOptionsToPackage{hyphens}{url}
\documentclass[
  11pt,
]{article}
\usepackage{xcolor}
\usepackage[margin=1in]{geometry}
\usepackage{amsmath,amssymb}
\setcounter{secnumdepth}{-\maxdimen} % remove section numbering
\usepackage{iftex}
\ifPDFTeX
  \usepackage[T1]{fontenc}
  \usepackage[utf8]{inputenc}
  \usepackage{textcomp} % provide euro and other symbols
\else % if luatex or xetex
  \usepackage{unicode-math} % this also loads fontspec
  \defaultfontfeatures{Scale=MatchLowercase}
  \defaultfontfeatures[\rmfamily]{Ligatures=TeX,Scale=1}
\fi
\usepackage{lmodern}
\ifPDFTeX\else
  % xetex/luatex font selection
\fi
% Use upquote if available, for straight quotes in verbatim environments
\IfFileExists{upquote.sty}{\usepackage{upquote}}{}
\IfFileExists{microtype.sty}{% use microtype if available
  \usepackage[]{microtype}
  \UseMicrotypeSet[protrusion]{basicmath} % disable protrusion for tt fonts
}{}
\makeatletter
\@ifundefined{KOMAClassName}{% if non-KOMA class
  \IfFileExists{parskip.sty}{%
    \usepackage{parskip}
  }{% else
    \setlength{\parindent}{0pt}
    \setlength{\parskip}{6pt plus 2pt minus 1pt}}
}{% if KOMA class
  \KOMAoptions{parskip=half}}
\makeatother
\setlength{\emergencystretch}{3em} % prevent overfull lines
\providecommand{\tightlist}{%
  \setlength{\itemsep}{0pt}\setlength{\parskip}{0pt}}
\usepackage{mathtools,mathrsfs}
\emergencystretch=2em
\usepackage{bookmark}
\IfFileExists{xurl.sty}{\usepackage{xurl}}{} % add URL line breaks if available
\urlstyle{same}
\hypersetup{
  hidelinks,
  pdfcreator={LaTeX via pandoc}}

\author{}
\date{}

\begin{document}
\section*{Reader scope}
The complete manuscripts in Parts 4--8 construct and relate the stated Gamma cyclic spectral representation, amplitude windows and configuration observable and operator-algebra maps. That Gamma cyclic representation does not reconstruct local interacting Yang--Mills theory. Retaining the full origin amplitude window adds a zero-energy sector beyond the single vacuum. The original Hamiltonians, hypotheses and limiting scopes are given in full in the manuscripts below. The supplied checker distinguishes exact algebra identities from separately labelled numerical illustrations. The subsequent interaction notes retain the full original electric and Wilson operator coefficients, a controlled fixed-box vacuum expansion, and the six-dimensional first physical gap matrix. These fixed-box corrections do not establish a regulator-uniform interacting continuum or its mass-gap failure.
\clearpage

\part{Physical eigenvalues at fixed volume}


\section{The original finite-box SU(2) Hamiltonian at small positive
coupling}\label{weak-coupling-part-1-the-original-finite-box-su2-hamiltonian-at-small-positive-coupling}

8 September 2026. This note retains the full nonlinear Wilson
Hamiltonian and proves a limit of its actual physical eigenvalues at
each fixed open box. The maximal-tree map, Haar density, kinetic tensor,
residual gauge action, and harmonic comparison are all specified. The
resulting fixed-volume theorem does not identify an interacting
continuum limit.

\subsection{1. The unchanged operator and its physical
space}\label{weak-coupling-part-1-the-unchanged-operator-and-its-physical-space}

Fix L at least 2, set m=2L, and use vertices n in \{-L,\ldots,L\} cubed,
all positive contained nearest-neighbor edges, and all contained
elementary faces. Write V,E,P for these sets, with \[
 |V|=(m+1)^3,\quad |E|=3m(m+1)^2,\quad |P|=3m^2(m+1).
\] On the product Haar probability space of SU(2) link variables, use
\(T_a=-i\sigma_a/2\), left derivatives X\_\{e,a\}, and
\(E_e=-\sum_aX_{e,a}^2\). Reverse traversal uses the inverse of the same
link. The trace W\_p is the fundamental trace of the original oriented
four-edge word. With physical spacing a\textgreater0 and coupling
g\textgreater0, \[
 H_g=\frac{2g^2}{a}\sum_e E_e
       +\frac1{2g^2a}\sum_p(2-W_p).
\tag{1}
\] Thus \(\kappa=2g^2/a\), \(b=1/(2g^2a)\), and \(\xi=1/(4g^4)\) remain
the previous parameters. In the original metric
\(c(A,B)=-\operatorname{tr}(AB)/2\), \(-\Delta_e^c=4E_e\); (1) includes
that factor. The scalar 2b\textbar P\textbar{} is contained in its
displayed potential and is not discarded.

Vertex gauge transformations act by \(U_e\mapsto h_sU_eh_t^{-1}\). Their
Haar average is an orthogonal projection. Its range is the physical
Hilbert space. The operator and form domains are the invariant parts of
H\^{}2 and H\^{}1 on the compact product group. Bounded smooth
multiplication by the potential preserves self-adjointness, compact
resolvent and these domains. The positive unit ground state is unique:
the modulus inequality gives a nonnegative minimizer; elliptic
regularity and the maximum principle make it smooth and positive; and
integration by parts gives \[
 q_{H_g-\mathcal E_g}(\psi_g f)
       =\frac{2g^2}{a}\sum_{e,a}\int\psi_g^2|X_{e,a}f|^2.
\tag{2}
\] Dividing a second ground state by the positive smooth first one in
(2) proves uniqueness. Gauge transformations therefore fix this actual
vacuum. Let \(\mathcal E_0(g)\le\mathcal E_1(g)\le\cdots\) be the
physical eigenvalues, with multiplicity and index zero for the ground
state, and let \(\Delta_L(g,a)=\mathcal E_1(g)-\mathcal E_0(g)\).

\subsection{2. Exact nonlinear maximal-tree coordinates, including
measure}\label{weak-coupling-part-1-exact-nonlinear-maximal-tree-coordinates-including-measure}

Root the graph at o=(-L,-L,-L). For each other vertex choose its parent
by decreasing the first coordinate, in the order 1,2,3, which exceeds
-L. These positively oriented parent edges form a tree \(\mathcal T\):
the sum of the three coordinate offsets strictly decreases toward o, and
every other vertex has one parent. There are \textbar V\textbar-1 tree
edges. Denote the remaining chords by \(\mathcal C\), of cardinality \[
 r=|E|-|V|+1=2m^3+3m^2.
\tag{3}
\] Let t\_v(U) be the ordered holonomy on the tree path from o to v,
with t\_o=I. Define, for every chord c from s to t, \[
 Z_c=t_s(U)U_ct_t(U)^{-1}.
\tag{4}
\] The coordinate map U -\textgreater{} ((U\_e)\emph{\{e in
tree\},(Z\_c)}\{c in chords\}) is a smooth global bijection of product
compact groups. Its inverse retains the specified tree links, computes
their same t\_v, and sets \(U_c=t_s^{-1}Z_ct_t\). For each fixed tree
configuration the chord changes are independent left and right
translations of Haar measure. Fubini therefore proves exact preservation
of the full product Haar probability measure by this bijection.

The gauge transformation h\_v=t\_v makes every tree link I and makes
each chord precisely Z\_c.~Under an arbitrary original gauge
transformation k, tree holonomies become \(t_v\mapsto k_ot_vk_v^{-1}\),
so \(Z_c\mapsto k_oZ_ck_o^{-1}\) simultaneously for all c.~Consequently
\[
 \mathcal U:F\longmapsto F((Z_c(U))_c)
\tag{5}
\] is an isometry from \(L^2(SU(2)^r,dZ)\) onto the subspace invariant
under based gauge transformations (k\_o=I). To establish surjectivity on
L\^{}2, first average a smooth function over the compact based gauge
group. In tree/chord variables its average is independent of the tree
variables: that group acts transitively on each tree-variable fibre and
leaves Z fixed. Smooth functions are dense and Haar averaging is an
orthogonal projection, so this equality extends by continuity to the
entire invariant L\^{}2 subspace. No evaluation of an L\^{}2 class on a
measure-zero slice is needed. The full physical space corresponds
exactly to simultaneous-conjugation-invariant F. No quotient measure or
residual gauge condition is omitted.

The exact transformed potential is
\(W(Z)=\sum_p(2-\operatorname{tr}U_p(Z))\), where tree links in the
original face word are I and chord links are Z\_c or Z\_c\^{}\{-1\} with
their original orientation. For an edge e and generator T\_a, define the
vector field \(\mathcal X_{e,a}\) on the chord manifold by varying that
single original edge through \(e^{tT_a}U_e\), starting from the tree-I
representative, recomputing (4), and differentiating at t=0. This is an
explicit smooth vector field: its coefficients are obtained by the
finite products and inverses (4), not a retained-loop approximation.
Here is its complete expression. Let \(\epsilon(v,e)\) be one when the
tree edge e lies on the root-to-v tree path, and zero otherwise. Write
\(L_{c,\alpha}F(Z)=\left.\frac d{dt}F(\ldots,e^{tT_\alpha}Z_c,\ldots)
\right|_{t=0}\) and
\(R_{c,\alpha}F(Z)=\left.\frac d{dt}F(\ldots,Z_ce^{tT_\alpha},\ldots)
\right|_{t=0}\). Then \[
 \mathcal X_{e,\alpha}=
 \begin{cases}
 L_{e,\alpha},&e\in\mathcal C,\\
 \displaystyle\sum_{c\in\mathcal C}
 [\epsilon(s(c),e)L_{c,\alpha}-\epsilon(t(c),e)R_{c,\alpha}],
       &e\in\mathcal T.
 \end{cases}
\tag{5a}
\] Indeed a varied tree edge changes each downstream t\_v to
\(e^{tT_\alpha}\) at the tree-I representative; differentiating (4)
gives the stated left and right fields with the minus sign from the
inverse t\_t. Each displayed field is divergence-free for product Haar.
Thus the exact transformed kinetic operator on smooth functions is
\(-(2g^2/a)\sum_{e,\alpha}\mathcal X_{e,\alpha}^{,2}\). Gauge invariance
of the original sum of squared derivatives shows that the exact
based-invariant Dirichlet form is \[
 q_g[F]=\frac{2g^2}{a}\int\sum_{e,a}|\mathcal X_{e,a}F|^2dZ
                   +\frac1{2g^2a}\int W(Z)|F|^2dZ.
\tag{6}
\] Indeed differentiation of (5) gives these fields at the
representative; changing tree representative rotates the three generator
components at each edge by an orthogonal adjoint matrix, preserving
their summed squares. The product-Haar identity just proved then
integrates that pointwise identity. Chord-edge variations alone supply a
full tangent frame, since they vary their own Z\_c by left
multiplication while leaving all other chords fixed. Thus the form is
elliptic on the compact chord manifold. Its domain is H\^{}1, its
represented operator has domain H\^{}2, and (5) intertwines both closed
forms and their self-adjoint operators. Restriction to the residual
invariant subspace preserves these statements.

\subsection{3. The unique minimum and its exact linear coordinate
maps}\label{weak-coupling-part-1-the-unique-minimum-and-its-exact-linear-coordinate-maps}

Each summand of W is nonnegative. For SU(2), trace equal to 2 forces the
group element to be I. If W=0, every face holonomy is I. Any monotone
path from the minimum corner o to a vertex consists of positive
coordinate steps; any two such words differ by interchanges of adjacent
steps in different directions. Each interchange is an actual contained
elementary square. Flatness therefore makes their holonomies identical.
For every positive edge s-\textgreater t, append it to a monotone path
to s and compare with a monotone path to t. This gives t\_s U\_e=t\_t,
in particular Z\_c=I for all chords. Conversely those identities make
every face I. The potential has exactly one zero on the chord manifold,
the identity configuration. This proof uses the original open box, not a
torus.

Use real edge and face cochains with their original counting inner
products. Define \[
 (d_0 f)(e)=f(t)-f(s),\qquad
 (d_1 x)(n;i,j)=x_i(n)+x_j(n+e_i)-x_i(n+e_j)-x_j(n).
\tag{7}
\] The target of d\_0 is \(\mathbb R^E\), and that of d\_1 is
\(\mathbb R^P\). Substitution proves d\_1d\_0=0. The same monotone path
interchange argument, now with additive real edge integrals, proves
\(\ker d_1=\operatorname{im}d_0\).

Let p\_v(x) be the additive integral along the specified tree path, and
define \(T:\mathbb R^E\to\mathbb R^{\mathcal C}\) by \[
 (Tx)_c=p_s(x)+x_c-p_t(x).
\tag{8}
\] This is precisely the derivative of (4) in each of the three
Lie-algebra components. Let j insert chord coordinates into the full
edge vector with zero tree entries. Then Tj=I, and
\(\ker T=\operatorname{im}d_0\): one implication follows by telescoping,
and in the other the potentials p\_v give x=d\_0p when all the chord
identities (8) vanish. Put \[
 G=TT^*,\quad R=T^*G^{-1/2},\quad K=\ker d_0^*,\quad C=d_1j.
\tag{9}
\] G is strictly positive because T is surjective. The displayed R is an
isometry onto K: R\^{}*R=I, its image is perpendicular to ker T, and the
dimensions agree. The map \(J=T^*G^{-1}\) is the inverse of T on K.
Since TJ=I=Tj, their difference is a gradient, giving \[
 d_1J=C,\qquad d_1R=CG^{1/2},\qquad
 R^*d_1^*d_1R=G^{1/2}C^*CG^{1/2}.
\tag{10}
\] These identities prove the metric relation between tree coordinates
and transverse edge cochains. In particular neither a tree metric nor a
Euclidean chord metric is silently substituted for the original kinetic
tensor. C is injective by ker d\_1=im d\_0 and the zero tree entries of
j.

In the exponential chart \(Z_c=\exp(y_c^aT_a)\), the principal kinetic
matrix in (6) at zero is \(G\otimes I_3\), by differentiating (4). The
Taylor expansion of each original face word, using
\(\operatorname{tr}(T_aT_b)=-\delta_{ab}/2\), gives \[
 W(y)=\frac14\sum_{a=1}^3\|Cy^a\|^2+O_L(|y|^3).
\tag{11}
\] One can check the coefficient without choosing a logarithm of the
whole word: its first derivative is the signed sum of its edge
generators, its trace derivative is zero, and the second trace
derivative is the trace of that sum squared. The commutator terms have
trace zero. Thus the quadratic term of 2-tr is one quarter the squared
signed sum. All higher word terms remain in W. Their fixed-chart bound
follows from Taylor's integral remainder for the finite smooth word, not
their removal. Injectivity of C proves nondegeneracy of this unique
minimum.

The Haar density in this chart is exactly \[
 dZ=\mathcal J(y)\,dy,\qquad
 \mathcal J(y)=(16\pi^2)^{-r}
       \prod_{c\in\mathcal C}
       \left[\frac{\sin(|y_c|/2)}{|y_c|/2}\right]^2,
       \quad |y_c|<2\pi.
\tag{12}
\] The continuous value at zero is used. Indeed unit-quaternion polar
coordinates have Haar measure \(\sin^2(r/2)dr,d\Omega/(4\pi^2)\);
division by the Cartesian polar Jacobian r\^{}2 gives (12) for one
factor. Its integral over the radius-2pi chart is one. This retains the
full density and its factor 16pi squared for every original chord.

\subsection{4. Exact open-box transverse
frequencies}\label{weak-coupling-part-1-exact-open-box-transverse-frequencies}

Set \(\lambda_j=4\sin^2(\pi j/[2(m+1)])\), j=0,\ldots,m. The complete
positive squared singular values of d\_1 on K are \[
 \lambda_{j_1}+\lambda_{j_2}+\lambda_{j_3},
 \quad j_i\in\{0,\ldots,m\},\quad
 k=\#\{i:j_i>0\}\ge2,
 \quad\hbox{multiplicity }k-1.
\tag{13}
\] We prove the boundary-sensitive formula. On the one-dimensional
interval with vertices l=0,\ldots,m and edges l=0,\ldots,m-1, let
\(Df(l)=f(l+1)-f(l)\). The vectors \[
 v_0(l)=(m+1)^{-1/2},\qquad
 v_j(l)=\sqrt{\frac2{m+1}}\cos\frac{\pi j(l+1/2)}{m+1},
\] \[
 w_j(l)=-\sqrt{\frac2{m+1}}\sin\frac{\pi j(l+1)}{m+1},
       \quad 1\le j\le m,
\tag{14}
\] are orthonormal bases in their respective counting spaces. The
geometric sum formula for roots of unity, followed by the cosine and
sine product identities, proves their inner products; their counts equal
the dimensions. The difference-of-cosines formula gives
\(Dv_j=\sqrt{\lambda_j}w_j\), and taking the adjoint gives
\(D^*w_j=\sqrt{\lambda_j}v_j\). Here l=n\_i+L is an explicit index map,
not a change of physical coordinates: edge midpoints remain
\(a(n+e_i/2)\).

Tensor these bases in the three coordinate directions. At a fixed triple
j, one-cochains have one basis vector in each direction i with
j\_i\textgreater0: use w\_\{j\_i\} in that direction and v\_\{j\_h\} in
the other two. These k vectors are orthonormal. The gradient of the
vertex tensor is the vector with components
\(s_i=\sqrt{\lambda_{j_i}}\). The oriented curl in (7) has components
\(s_i x_h-s_h x_i\). Expanding their squares gives \[
 \sum_{i<h}|s_i x_h-s_h x_i|^2
  =\left(\sum_i s_i^2\right)\sum_i|x_i|^2
                  -\left|\sum_i s_i x_i\right|^2.
\tag{15}
\] Consequently the orthogonal complement of the gradient has dimension
k-1 and squared curl value \(\sum\lambda_{j_i}\). For k=1 it is zero
dimensional. This proves (13) on a complete tensor basis, with all open
boundaries included. Its total multiplicity is \(3m^2+2m^3=r\), agreeing
with (3), not an extra polarization count.

Since sine is strictly increasing on the relevant half interval, the
smallest frequency is \[
 \sigma_{\min}=\sqrt8\sin\frac{\pi}{2(m+1)}.
\tag{16}
\] It occurs exactly at triples (1,1,0),(1,0,1),(0,1,1), each with one
transverse component, hence has multiplicity three.

\subsection{5. The comparison operator and the residual gauge
constraint}\label{weak-coupling-part-1-the-comparison-operator-and-the-residual-gauge-constraint}

Write \(\sigma_1,\ldots,\sigma_r\) for the positive square roots in
(13), with multiplicity. In the original chord coordinates rescaled by
the explicit map y=g x, the comparison form is the operator \[
 H_{\rm osc}=\frac1a\left[-2\sum_{\alpha=1}^3\sum_{c,d\in\mathcal C}
                     G_{cd}\partial_{x_c^\alpha}\partial_{x_d^\alpha}
              +\frac18\sum_{a=1}^3\|Cx^a\|^2\right].
\tag{17}
\] The kinetic term in (17) is written with ordinary coordinate
derivatives; its quadratic form is positive because G is positive
definite. The exact local Hilbert map is \[
 F(\exp(gx))=g^{-3r/2}\mathcal J(gx)^{-1/2}f(x),
\tag{18}
\] for functions supported inside this chart. By (12) this preserves
their norm exactly. In particular the rescaling has retained both its
volume power and its nonconstant Haar factor.

Let A\^{}\{ij\}(y) be the full smooth principal matrix of the kinetic
form in (6), with A(0)=G tensor I\_3. Under (18) the exact kinetic form
is \[
 \frac2a\int A^{ij}(gx)
 \left(\partial_i f-\frac g2(\partial_i\log\mathcal J)(gx)f\right)^*
 \left(\partial_j f-\frac g2(\partial_j\log\mathcal J)(gx)f\right)dx,
\tag{19}
\] and its exact potential is multiplication by \(W(gx)/(2g^2a)\).
Equation (19) follows by differentiating the displayed density in (18)
inside (6). Every metric coefficient and density derivative remains in
it. On bounded x domains these forms converge to (17) by smoothness and
(11).

The linear change x=G\^{}\{1/2\}z in each color component uses the
constant Jacobian \(\det(G)^{3/2}\); multiplication by its square root
is part of the unitary change of L\^{}2 coordinates. Diagonalize the
positive matrix \(G^{1/2}C^*CG^{1/2}\) by an orthogonal matrix. By (10)
its eigenvalues are exactly (13). Thus (17) is unitarily equivalent to
\[
 \sum_{\nu=1}^r\left[-\frac2a\Delta_{z_\nu}
                    +\frac{\sigma_\nu^2}{8a}|z_\nu|^2\right],
                  \quad z_\nu\in\mathbb R^3.
\tag{20}
\] All the linear matrices act on spatial/chord indices only. Residual
SU(2) conjugation therefore remains simultaneous adjoint rotation of the
r three-component vectors, not r independently removed color rotations.

For a single real oscillator in (20), factorization with
\(\partial_z+\sigma z/4\) and its adjoint gives energy
\(\sigma/(2a)+n\sigma/a\), n=0,1,\ldots{} . The Gaussian vacuum and
successive Hermite polynomials give a complete eigenbasis (equivalently,
their generating functions are complete in the Gaussian L\^{}2 space).
Consequently the oscillator ground energy is \[
 \mu_0=\frac3{2a}\sum_{\nu=1}^r\sigma_\nu.
\tag{21}
\] It is a simultaneous-rotation singlet. Every one-quantum state
transforms as an adjoint vector, which has no invariant vector:
invariance under all three coordinate-axis rotations forces its
components to vanish. Every state with at least two quanta has
excitation energy at least \(2\sigma_{\min}/a\). This value is achieved
by the invariant quadratic \(\sum_a a^\dagger_{\nu,a}a^\dagger_{\nu,a}\)
on a lowest mode. There are six independent invariant quadratics among
the three lowest spatial modes: the contractions indexed by
\(\nu\le\eta\). For each ordered pair of vector indices, invariance of a
bilinear tensor under rotations forces it to be a scalar multiple of the
identity (first use the sign-changing half-turns to remove off-diagonal
entries, then quarter-turns to equate diagonal entries). Symmetry of
creation operators therefore gives precisely those six, with no
additional invariants at that energy. Thus the first physical oscillator
excitation has energy \(2\sigma_{\min}/a\) and multiplicity six.

\subsection{6. Limit of actual physical eigenvalues, not a substitution
of
operators}\label{weak-coupling-part-1-limit-of-actual-physical-eigenvalues-not-a-substitution-of-operators}

Let \(\mu_k\) be the sorted eigenvalues of (20) on its simultaneous-
rotation-invariant space. We claim, for every fixed L,a,k, \[
 \lim_{g\downarrow0}\mathcal E_k(g)=\mu_k.
\tag{22}
\] Here is a form/min--max proof retaining the nonconstant coefficients.
All localization functions used below depend only on the chart radius
\(|y|\), so are invariant under the residual simultaneous rotations.

For the upper bound choose the first k+1 invariant oscillator
eigenvectors, multiply them by a smooth cutoff equal to one on
\textbar y\textbar\textless=rho and supported in
\textbar y\textbar\textless2rho, and use (18). The number rho is fixed
inside the exponential chart. In x coordinates the cutoffs tend to one.
The oscillator vectors are polynomials times a decaying Gaussian. On
this fixed y chart, A, the density derivatives, and their first
derivatives are bounded; moreover \(W(gx)/(g^2)\le c_\rho |x|^2\),
because W vanishes to second order. The cutoff derivatives in (19) have
an additional g factor. Dominated convergence therefore gives the
identity matrix as their norm Gram limit and the oscillator energy
matrix as their form Gram limit. These vectors are independent for small
g, remain exactly gauge invariant, and lie in the full original form
domain. The min--max principle gives \(\limsup\mathcal E_k(g)\le\mu_k\).

For the lower bound fix any epsilon\textgreater0. By (11), positive
definiteness and smoothness, choose rho\textgreater0 so that on
\textbar y\textbar\textless2rho the principal matrix and the quadratic
potential differ from their values in (17) by relative form errors at
most epsilon. The density is positive and smooth there. After
multiplying by its square root, the kinetic derivatives are
\(\partial_i-\tfrac12\partial_i\log\mathcal J\). The elementary
inequality \(|u-v|^2\ge(1-\epsilon)|u|^2-(\epsilon^{-1}-1)|v|^2\) shows
that these bounded density derivatives cost at most
\(C_{\rho,\epsilon}g^2\|F\|^2\) in the original form. After reducing rho
or epsilon if needed, the local form is therefore bounded below by
(1-epsilon) times the full quadratic comparison form, minus that
displayed error. The constants are finite suprema of the original
coefficients on this fixed compact chart; they need not be independent
of L and are not treated as such.

Choose smooth invariant chi\_0,chi\_1 with chi\_0 squared plus chi\_1
squared equal to one, chi\_0 supported in
\textbar y\textbar\textless2rho and equal to one on
\textbar y\textbar\textless=rho. For example use the cosine and sine of
a smooth radial function constant at 0 near the inner ball and at pi/2
beyond the outer ball. The exact product-rule localization identity for
(6) is \[
 q_g[F]=q_g[\chi_0F]+q_g[\chi_1F]
       -\frac{2g^2}{a}\int\sum_{e,a,j}|\mathcal X_{e,a}\chi_j|^2|F|^2dZ.
\tag{23}
\] Expand each squared derivative to prove it; the mixed terms vanish by
\(\sum_j\chi_j\mathcal X\chi_j=0\). The error is bounded by
\(C_\rho g^2\|F\|^2\). On the support of chi\_1 the actual W has a
positive lower bound c\_rho, since its only zero is the identity and the
manifold is compact. Thus the exterior form is at least
\(c_\rho/(2g^2a)\|\chi_1F\|^2\).

Map chi\_0F by the chart density and dilation into ordinary Euclidean
L\^{}2, extending by zero. It is a bounded linear map of Hilbert spaces,
preserves its norm, and its image remains in the invariant form space.
The oscillator bound \[
 q_{\rm osc}[u]\ge\mu_k\|u\|^2
       -\sum_{j<k}(\mu_k-\mu_j)|\langle\varphi_j,u\rangle|^2
\tag{24}
\] follows by expansion in its complete invariant eigenbasis. Each of
the k linear functionals in (24), composed with the local map, is
bounded on the original Hilbert space. Combining (23), the local
comparison, and the exterior lower bound gives a lower bound
\((1-\epsilon)\mu_k-C_{\rho,\epsilon}g^2\) on a subspace of codimension
at most k, once g is small enough that the exterior bound exceeds
(1-epsilon)mu\_k. Min--max gives that same lower bound for
\(\mathcal E_k(g)\). First send g to zero, then epsilon to zero. This
proves (22). The density maps, cutoffs, finite-rank comparisons and
domains in this argument are all for the original operator (1).

In particular the actual physical finite-box gap satisfies \[
 \boxed{\lim_{g\downarrow0}\Delta_L(g,a)
      =\frac{4\sqrt2}{a}\sin\frac{\pi}{4L+2}.}
\tag{25}
\] The next six eigenvalues have this same limiting excitation energy.
For each fixed box this is a limit at positive couplings of the full
interacting physical spectrum. It has not been obtained by using an
unprojected one-quantum color-vector state.

\subsection{7. What this does and does not prove about simultaneous
limits}\label{weak-coupling-part-1-what-this-does-and-does-not-prove-about-simultaneous-limits}

Equation (25) supplies a genuine gap-closing sequence of finite
interacting regulators. For example take L\_j=j squared and
a\_j=1/(100j), j\textgreater=2. Its limiting comparison gap is
\(400\sqrt2 j\sin[\pi/(4j^2+2)]\), which tends to zero. For each fixed
j, (25) proves the existence of a number g\_j\textgreater0, which can
also be required to satisfy g\_j\textless1/j, such that the actual gap
differs from this comparison value by less than 1/j. Every Hamiltonian
in this sequence has its full nonlinear potential and a strictly
positive coupling. The actual first excited eigenvectors can be chosen
of unit norm and orthogonal to their own actual vacua; their
finite-regulator spectral measures are then the atoms at those actual
positive excitation energies.

This existence statement supplies neither an explicit admissible running
coupling at every cutoff nor convergence of local observables/vacua to a
nontrivial interacting four-dimensional theory. In particular
convergence of these excited spectral atoms toward zero does not
identify their limits as distinct physical continuum states, or prove
nonzero continuum spectral weight on every interval (0,epsilon). Neither
identification is assumed. The nonuniform L-dependence in the
localization proof is retained, and the next calculation must control
these actual state/observable maps in the simultaneous limit. Equation
(25) is not a Yang--Mills mass-gap counterexample.

\subsection{8. Primary mathematical and lattice
context}\label{weak-coupling-part-1-primary-mathematical-and-lattice-context}

The Hamiltonian framework is J. Kogut and L. Susskind, \emph{Hamiltonian
formulation of Wilson's lattice gauge theories}, Physical Review D 11
(1975), 395--408, https://doi.org/10.1103/PhysRevD.11.395. Maximal-tree
SU(2) coordinates are also treated by C. W. Bauer, I. D'Andrea, M.
Freytsis and D. M. Grabowska, \emph{A new basis for Hamiltonian SU(2)
simulations}, https://arxiv.org/abs/2307.11829; this note derives its
own full maps (4)--(10) without any spin truncation. For clarity, in the
original link representation their displayed maximal-tree Hamiltonian
corresponds to \(H_{\rm B}(g_{\rm B},a)=g_{\rm B}^2\sum_eE_e/(2a)
+\sum_p(2-\operatorname{tr}U_p)/(g_{\rm B}^2a)\). The exact
parameter/energy dictionary for the same edge Casimir and face trace is
\[
 H_g=\sqrt2\,H_{\rm B}(2^{3/4}g,a).
\] Substitution gives respectively the original coefficients \(2g^2/a\)
and \(1/(2g^2a)\) in (1). This is an explicit operator identity on the
common original-link domain; its energy multiplier is not dropped. Each
chosen maximal tree is then related to that link representation by its
own exact Haar unitary (5). The leading small-coupling comparison uses
the localization/min--max method of B. Simon, \emph{Semiclassical
analysis of low lying eigenvalues. I. Non-degenerate minima: asymptotic
expansions}, Annales de l'Institut Henri Poincare A 38 (1983), 295--308,
Sections 2,3,6,
https://www.numdam.org/item/AIHPA\_1983\_\_38\_3\_295\_0/. The
finite-box density, kinetic tensor and residual-gauge implementation
needed here are proved explicitly above. No novelty of the general
harmonic-approximation method is claimed.

\subsection{9. Actual-state
continuation}\label{weak-coupling-part-1-actual-state-continuation}

The companion calculation \emph{Actual weak-coupling vacuum, observable
and spectral maps} proves strong convergence of the actual vacuum under
the full-measure logarithm/Haar/dilation isometry, and operator-norm
convergence of every fixed finite spectral projection. It constructs an
explicit bounded physical radial observable whose centered raw spectral
measure retains all Laguerre weights, rather than inferring state
survival from (25) alone. For a lowest spatial mode and the stated
explicit observable parameter, the first positive limiting weight is
3/(8sqrt(2)), at the energy in (25), and the total centered mass is
1-2\^{}(-3/2). A particular simultaneous very-weak-coupling
macroscopic-mode diagonal instead sends this entire raw measure to a
nonzero atom at zero. These maps and that zero-energy consequence are
proved in the companion; none is identified with the target interacting
continuum. The companion is included in the same readable source root
with this manuscript.

\clearpage
\part{Actual vacuum, observables and spectral maps}
\section{Actual weak-coupling vacuum, observable and spectral
maps}\label{weak-coupling-part-2-actual-weak-coupling-vacuum-observable-and-spectral-maps}

8 September 2026. The finite open box and its physical spacing remain
fixed through Sections 1--6. We prove convergence of the actual vacuum
and finite spectral projections of the nonlinear Hamiltonian, and
compute an exact family of limiting observable spectral measures.
Section 7 records what a particular simultaneous low-energy diagonal
actually does to those measures. It does not identify that diagonal with
interacting four-dimensional Yang--Mills theory.

\subsection{1. Original objects and the full-measure Hilbert
map}\label{weak-coupling-part-2-original-objects-and-the-full-measure-hilbert-map}

Use the open box with vertices in \(\{-L,\ldots,L\}^3\), \(L\ge2\), all
its positive edges E and faces P, and physical spacing \(a>0\). Put
\(m=2L\), \(r=2m^3+3m^2\). The original operator is \[
 H_g=\frac{2g^2}{a}\sum_{e\in E}E_e+
       \frac1{2g^2a}\sum_{p\in P}(2-\operatorname{tr}U_p),
 \qquad E_e=-\sum_{\alpha=1}^3X_{e,\alpha}^2,
 \quad T_\alpha=-i\sigma_\alpha/2.                       \tag{1}
\] In particular \(\kappa=2g^2/a\), \(b=1/(2g^2a)\), \(\xi=1/(4g^4)\),
and the scalar \(2b|P|\) are unchanged. The physical Hilbert space is
the invariant subspace for all vertex gauge transformations. Let
\(\psi_g>0\) be the actual unit ground state and \(\mathcal E_g\) its
actual energy.

Here are the exact coordinate data used below. Root the fixed maximal
tree at \((-L,-L,-L)\), with parent obtained by decreasing the first
coordinate in the order 1,2,3 that exceeds \(-L\). If \(t_v\) is its
ordered tree holonomy, the chord variables are \[
 Z_c=t_{s(c)}U_ct_{t(c)}^{-1},\qquad c\in\mathcal C.       \tag{2}
\] The inverse retains the tree links and sets \(U_c=t_s^{-1}Z_ct_t\).
Independent Haar translations in each chord, followed by Fubini, prove
that this coordinate map preserves product Haar measure. Based gauge
averaging eliminates the tree variables exactly. The residual physical
condition on functions of Z is simultaneous conjugation of every Z\_c by
one element of SU(2), not independent conjugation of each chord.

The full proof of the quotient form, boundary-sensitive frequencies and
finite-box eigenvalue convergence is in the companion manuscript
\emph{The original finite-box SU(2) Hamiltonian at small positive
coupling}. We retain its exact maps, not merely its spectral conclusion.
For real edge cochains, let \(p_v\) be the additive tree-path integral,
and set \[
 (Tx)_c=p_s(x)+x_c-p_t(x),\quad G=TT^*,\quad C=d_1j,
 \quad B=G^{1/2}C^*CG^{1/2}.                              \tag{3}
\] Here j inserts chord cochains with zero tree entries; d\_1 is the
original oriented face curl with the counting inner products. The
identities \(Tj=I\), \(\ker T=\operatorname{im}d_0=\ker d_1\) and
\(d_1T^*G^{-1}=C\) show that G and B are strictly positive. Choose once
an orthogonal matrix O with \[
 O^TBO=\operatorname{diag}(\sigma_1^2,\ldots,\sigma_r^2),
 \quad \sigma_\nu>0,
 \qquad z^\alpha=O^TG^{-1/2}x^\alpha.                     \tag{4}
\] The three Lie-algebra components are denoted by \(\alpha\); they are
not the physical spacing a. O acts on chord indices only.

Except for the Haar-null set where a chord equals \(-I\), every Z has a
unique logarithm \(Z_c=\exp(y_c^\alpha T_\alpha)\), \(|y_c|<2\pi\). The
exact density and dilated chart are \[
 \mathcal J(y)=(16\pi^2)^{-r}
    \prod_c\left(\frac{\sin(|y_c|/2)}{|y_c|/2}\right)^2,
 \qquad \Omega_g=\{x:|x_c|<2\pi/g\text{ for all }c\}.
                                                               \tag{5}
\] For a function F in the based quotient define \[
 (\mathcal B_gF)(x)=
 \begin{cases}
 g^{3r/2}\mathcal J(gx)^{1/2}F(\exp(gx)),&x\in\Omega_g,\\
 0,&x\notin\Omega_g.
 \end{cases}                                                   \tag{6}
\] The change of variables in (5) proves this is an isometry into
\(L^2(\mathbb R^{3r},dx)\), onto the closed subspace supported in
\(\Omega_g\). Its adjoint has the reciprocal density and dilation on
that subspace and vanishes on its orthogonal complement. It is not
claimed onto the full Euclidean space for fixed g. Both \(\Omega_g\) and
the density are invariant under simultaneous color rotations, so (6)
intertwines the exact residual physical projections.

The full transformed local form, on compactly supported chart functions,
is \[
 \frac2a\int A^{ij}(gx)
  \left(\partial_i f-\frac g2(\partial_i\log\mathcal J)(gx)f\right)^*
  \left(\partial_j f-\frac g2(\partial_j\log\mathcal J)(gx)f\right)dx
 +\int\frac{W(gx)}{2g^2a}|f|^2dx,                           \tag{7}
\] where W is the full sum in (1) expressed by the original chord words,
\(A(0)=G\otimes I_3\), and \[
 W(y)=\tfrac14\sum_\alpha\|Cy^\alpha\|^2+O_L(|y|^3).
                                                               \tag{8}
\] All nonconstant coefficients and all word remainders remain in (7).

\subsection{2. Compactness of actual bounded-energy
vectors}\label{weak-coupling-part-2-compactness-of-actual-bounded-energy-vectors}

We first prove the compactness needed for actual eigenvectors. Suppose
\(g_n\downarrow0\), \(\|F_n\|\le K\), and \(q_{g_n}[F_n]\le K\), with
fixed L and a. Write \(f_n=\mathcal B_{g_n}F_n\). These vectors have a
strongly convergent subsequence in Euclidean L\^{}2.

To prove this, choose a fixed small radius \(\rho>0\) within the
logarithm chart. The positivity of \(C^*C\), (8), and Taylor's integral
remainder give a number \(c>0\) with \(W(y)\ge c|y|^2\) for
\(|y|<\rho\). Flatness of every original face on the contractible open
box implies all chords are I, by the ordered square-interchange argument
for monotone tree paths. Thus W has exactly one zero on the compact
chord manifold. On the complement of this radius-\(\rho\) neighborhood
it has a strictly positive minimum \(c_\rho\).

Both terms of the original form are nonnegative. Its potential part
therefore yields, when \(g_n R<\rho\), \[
 \int_{|x|>R}|f_n|^2dx
   \le \frac{2aK}{cR^2}+\frac{2aK g_n^2}{c_\rho}.        \tag{9}
\] Indeed on \(R<|x|<\rho/g_n\) the potential is at least
\(c|x|^2/(2a)\); outside that region it is at least
\(c_\rho/(2g_n^2a)\). These are estimates for the full W, not a
replacement quadratic potential.

On any fixed ball \(|x|\le R\), for sufficiently small g the smooth
principal matrix \(A(gx)\) has a positive lower eigenvalue independent
of g. The logarithmic-density derivative in (7) is bounded there. The
kinetic part of (7), the L\^{}2 bound, and \(|u|^2\le2|u-v|^2+2|v|^2\)
consequently bound \(\|f_n\|_{H^1(B_R)}\). No derivative estimate is
asserted across the boundary of \(\Omega_g\); every fixed ball is
eventually strictly inside. Rellich compactness on a succession of fixed
balls and a diagonal subsequence give strong local L\^{}2 convergence.
The uniform tail estimate (9), first with R large and then n large,
upgrades it to global strong L\^{}2 convergence. It also preserves inner
products for any finite family of such vectors by polarization. Every
constant here may depend on L,a; none has been assumed uniform in
volume.

\subsection{3. Actual vacuum and finite spectral
projections}\label{weak-coupling-part-2-actual-vacuum-and-finite-spectral-projections}

Define \[
 H_{\rm osc}=\frac1a\left[-2\sum_{\alpha,c,d}G_{cd}
      \partial_{x_c^\alpha}\partial_{x_d^\alpha}
                +\frac18\sum_\alpha\|Cx^\alpha\|^2\right].
                                                               \tag{10}
\] Its physical subspace consists of the simultaneous-rotation
invariants. The companion finite-box proof gives actual physical
eigenvalues \(\mathcal E_k(g)\to\mu_k\), where the \(\mu_k\) are the
sorted physical eigenvalues of (10), with multiplicity. That proof uses
upper trial spaces and the full-form localization identity; it does not
choose the oscillator vacuum in place of \(\psi_g\).

We now prove the stronger vector conclusion. The comparison trial bound
makes \(q_g[\psi_g]=\mathcal E_g\) bounded at fixed L,a. Section 2 gives
compactness of \(\mathcal B_g\psi_g\). For a smooth compactly supported
invariant test function h, pull h back by (6). Its support stays in the
logarithm chart, so it is in the exact original form domain. The
eigenvector identity is \[
 q_g(\psi_g,\mathcal B_g^*h)
       =\mathcal E_g\langle\mathcal B_g\psi_g,h\rangle.
                                                               \tag{11}
\] Local H\^{}1 weak compactness from Section 2 and the uniform
coefficient convergence on the support of h let (7) pass to the limit in
(11). Every strongly converging subsequence therefore has a norm-one
limit f satisfying the weak equation \(H_{\rm osc}f=\mu_0f\). The same
energy bounds put f in the oscillator form domain: take increasing balls
in the local lower semicontinuity estimates for derivatives and the
quadratic potential. Compactly supported invariant test functions are a
form core, obtained by invariant smooth radial cutoffs followed by
convolution and rotation averaging. Thus the weak equation identifies
the represented self-adjoint operator, not a different extension.

The oscillator ground eigenspace is one-dimensional. The exact maps (6)
send the positive \(\psi_g\) to nonnegative functions, fixing their
phase. It follows that the entire family converges, not just a
subsequence: \[
 \mathcal B_g\psi_g\longrightarrow\Phi_0
           \quad\hbox{strongly in }L^2(\mathbb R^{3r}).      \tag{12}
\] Put \[
 S=G^{-1/2}B^{1/2}G^{-1/2}.
\] The exact unit Gaussian in chord coordinates is \[
 \Phi_0(x)=(\det S)^{3/4}(4\pi)^{-3r/4}
       \exp\left(-\frac18\sum_\alpha (x^\alpha)^TSx^\alpha\right).
                                                               \tag{13}
\] In (4) it is the product \(\prod_\nu(\sigma_\nu/(4\pi))^{3/4}
e^{-\sigma_\nu|z_\nu|^2/8}\), with the additional constant Jacobian
factor \((\det G)^{-3/4}\) when expressed as an x-wavefunction. Each
real z coordinate has vacuum variance \(2/\sigma_\nu\). The density
factor, this variance, and the ground energy
\(\mu_0=3\sum_\nu\sigma_\nu/(2a)\) are all retained.

Let I be a bounded interval whose endpoints avoid the spectrum of (10),
and let P\_g(I) and P\_0(I) be the physical spectral projections of H\_g
and H\_osc. Then, as operators on the common Euclidean physical
subspace, \[
 \|\mathcal B_gP_g(I)\mathcal B_g^*-P_0(I)\|\longrightarrow0.
                                                               \tag{14}
\] Here and below the adjoint is zero on the complement of \(\Omega_g\).
For proof, eigenvalue convergence gives the same finite rank on both
sides for small g. Choose an orthonormal eigenbasis of the range on the
left. Each vector has bounded absolute energy, so Section 2 gives a
simultaneously strongly convergent subsequence. The argument (11)
identifies the limits as eigenvectors in the indicated oscillator
interval. Their inner products persist. Equality of ranks makes them a
basis of that entire interval. Finite sums of their rank-one projections
converge in operator norm. Every subsequence has the same projection
limit, proving (14) for the original family. Applying the same reasoning
to intervals shifted by the convergent \(\mathcal E_g\) gives (14) for
excitation operators \(H_g-\mathcal E_g\) and \(H_{\rm osc}-\mu_0\).

\subsection{4. Transport of bounded observables and raw spectral
measures}\label{weak-coupling-part-2-transport-of-bounded-observables-and-raw-spectral-measures}

Suppose \(B_g\) is multiplication by a bounded physical function on the
original chord manifold, with \(\sup_g\|B_g\|<\infty\), and its
coordinate multiplier \(B_g(\exp(gx))\) converges at each x to a bounded
invariant multiplier B\_0(x). This statement defines the class of maps
considered in this section; Section 5 constructs such a map exactly.
Equations (12), the multiplier bound, and dominated convergence against
\(|\Phi_0|^2\) give \[
 \mathcal B_gB_g\psi_g\to B_0\Phi_0,\qquad
 c_g:=\langle\psi_g,B_g\psi_g\rangle
                  \to c_0:=\langle\Phi_0,B_0\Phi_0\rangle.
                                                               \tag{15}
\] Consequently the actual vacuum-orthogonal vectors \[
 v_g=(B_g-c_g)\psi_g
       \quad\hbox{satisfy}\quad
 \mathcal B_gv_g\to v_0=(B_0-c_0)\Phi_0.                  \tag{16}
\] No unit-vector rescaling or division by a small norm is used here.

Define their raw positive excitation measures by \[
 \nu_g(J)=\langle v_g,\mathbf1_J(H_g-\mathcal E_g)v_g\rangle,
 \qquad
 \nu_0(J)=\langle v_0,\mathbf1_J(H_{\rm osc}-\mu_0)v_0\rangle.
                                                               \tag{17}
\] They converge weakly against every bounded continuous function on
\([0,\infty)\), with their total masses retained. To verify the tail
claim needed here, choose an oscillator energy cutoff R not at an
eigenvalue. From (14) and (16), the masses in {[}0,R{]} and the total
masses converge. Completeness of the oscillator eigenbasis makes
\(\nu_0((R,\infty))\to0\) as \(R\to\infty\); hence the family \(\nu_g\)
is tight as g tends to zero. On {[}0,R{]}, isolate the finitely many
eigenvalue clusters and use eigenvalue convergence, (14), and uniform
continuity of the test function. The bounded test-function tail is then
controlled by tightness. This proves the stated convergence without
asserting a moment bound for an unspecified unbounded observable.

In particular for every fixed physical \(t\ge0\), \[
 \langle v_g,e^{-t(H_g-\mathcal E_g)}v_g\rangle
       \to\langle v_0,e^{-t(H_{\rm osc}-\mu_0)}v_0\rangle.
                                                               \tag{18}
\] The same finite-projection argument applies to any actual vectors
whose images under (6) converge strongly, including limits of bounded
averaged translation operators. Multiplication is not a necessary
hypothesis for that last vector statement.

\subsection{5. An explicit gauge-invariant field observable with nonzero
mass}\label{weak-coupling-part-2-an-explicit-gauge-invariant-field-observable-with-nonzero-mass}

Choose one spatial mode \(\nu\) in (4), and a real number \(\theta\).
Let \(\chi(y)\) be a smooth simultaneous-rotation-invariant function
equal to one near zero and supported strictly inside the product
logarithm chart. Set on that chart \[
 B_{g,\theta}(Z)=\exp\left(i\theta\chi(y)
              |(O^TG^{-1/2}(y/g))_\nu|^2\right),          \tag{19}
\] and set it equal to 1 outside the chart. This is a globally smooth
physical multiplication operator for each g\textgreater0. The support
property makes the extensions agree on a full neighborhood of the chart
boundary. It is unitary, and its pointwise multiplier limit is
\(B_{0,\theta}(x)=e^{i\theta|z_\nu|^2}\). Thus (15)--(18) apply to the
actual nonlinear model.

Write \(\sigma=\sigma_\nu\), \(k=3/2\), and \(\beta=4\theta/\sigma\).
Under \(|\Phi_0|^2\), the exact radial variable \(q=\sigma|z_\nu|^2/4\)
has density \(q^{k-1}e^{-q}/\Gamma(k)\) on \([0,\infty)\). Integrating
the Gaussian with its original variance yields \[
 c_0=(1-i\beta)^{-3/2},\qquad
 \|v_0\|^2=1-(1+\beta^2)^{-3/2}.                          \tag{20}
\] The branch of the power is fixed by continuity from \(\beta=0\). In
particular \(\theta\ne0\) gives positive limiting centered mass. An
unchanged bounded continuous physical function B(Z), without the
explicit coordinate map (19), instead has limit B(I) and zero centered
mass by (15). Both statements follow from the same map; neither
establishes that the underlying observables are unrelated.

\subsection{6. Complete spectral coefficients of the explicit
observable}\label{weak-coupling-part-2-complete-spectral-coefficients-of-the-explicit-observable}

For clarity we compute all coefficients, not only a first energy moment.
Use the generalized Laguerre polynomials \(L_n^{k-1}(q)\), fixed by \[
 \sum_{n=0}^\infty L_n^{k-1}(q)z^n
       =(1-z)^{-k}\exp\left(-\frac{qz}{1-z}\right),\quad |z|<1.
                                                               \tag{21}
\] Equivalently their coefficients are
\(\sum_{j=0}^n(-1)^j\Gamma(n+k)q^j/
[\Gamma(k+j)(n-j)!j!]\). The Rodrigues expression
\(q^{1-k}e^q(d/dq)^n(e^{-q}q^{n+k-1})/n!\), obtained from Leibniz's
rule, proves by n integrations by parts that their orthogonal squared
norms in the Gamma probability measure are \((k)_n/n!\). Boundary terms
vanish at infinity exponentially and at zero by the powers appearing
before the nth integration. The equation \(qL_n''+(k-q)L_n'+nL_n=0\)
follows directly from their coefficients. Conjugating the radial
oscillator by its ground Gaussian gives \[
 (H_{\rm osc}-\mu_0)(L(q)\Phi_0)
     =\frac{2\sigma}{a}[-qL''-(k-q)L']\Phi_0.             \tag{22}
\] Thus \(\sqrt{n!/(k)_n}L_n^{k-1}(q)\Phi_0\) has excitation energy
\(2n\sigma/a\). These vectors are complete in the radial subspace: if a
radial L\^{}2 function is orthogonal to every polynomial, its
Gamma-weighted Laplace transform is analytic on \(\Re t>-1/2\) by
Cauchy--Schwarz and has every derivative zero at t=0. Analyticity and
uniqueness of the Fourier transform of its integrable weighted measure
then make that function zero. Other modes remain in their vacuum
factors.

Integrate (21) against \(e^{i\beta q}\) times the Gamma density. For z
near zero the integral is absolutely convergent, and it equals
\((1-i\beta+i\beta z)^{-k}\). Comparing coefficients gives \[
 \langle \sqrt{n!/(k)_n}L_n^{k-1}\Phi_0,
                   e^{i\theta|z_\nu|^2}\Phi_0\rangle
  =\sqrt{(k)_n/n!}\frac{(-i\beta)^n}{(1-i\beta)^{n+k}}.
                                                               \tag{23}
\] The vacuum coefficient is precisely (20); centering removes exactly
n=0. Therefore the complete limiting raw measure from the original
actual state is \[
 \boxed{\nu_0=
   \sum_{n=1}^\infty
    \frac{(3/2)_n}{n!}
    \frac{\beta^{2n}}{(1+\beta^2)^{n+3/2}}
          \delta_{2n\sigma/a}.}                         \tag{24}
\] Every weight is nonnegative, and the binomial generating series sums
them to \(1-(1+\beta^2)^{-3/2}\), matching the actual norm limit. The
full physical-time correlation is consequently \[
 \boxed{\lim_{g\downarrow0}
   \langle v_{g,\theta},e^{-t(H_g-\mathcal E_g)}v_{g,\theta}\rangle
 = [1+\beta^2(1-e^{-2\sigma t/a})]^{-3/2}
                   -(1+\beta^2)^{-3/2},\quad t\ge0.}     \tag{25}
\] This identity describes a proved fixed-box limit of full interacting
matrix elements. It does not assert that the original finite-g measure
equals the right side before taking the limit.

For a lowest mode, \(\sigma=\sqrt8\sin(\pi/(4L+2))\). Choosing the
explicit \(\theta=\sigma/4\) gives \(\beta=1\), centered mass
\(d=1-2^{-3/2}\) and first positive atom of weight \(3/(8\sqrt2)\) at
energy \(4\sqrt2\sin(\pi/(4L+2))/a\). Its positive raw mass is a state
calculation, not a claim based only on the eigenvalue location.

\subsection{7. Exact fate of this particular macroscopic low-mode
diagonal}\label{weak-coupling-part-2-exact-fate-of-this-particular-macroscopic-low-mode-diagonal}

There is a useful test of what the gap-closing diagonal in the companion
paper actually retains. Take \(L_j=j^2\), \(a_j=1/(100j)\),
j\textgreater=2, choose a lowest mode at each j, and set
\(\theta_j=\sigma_j/4\). Let \(\nu_{g,j}\) be its actual centered raw
measure (17),(19). The proved finite-box weak convergence permits a
choice of strictly positive \(g_j<1/j\) such that \[
 d_{\rm BL}(\nu_{g_j,j},\nu_{0,j})<1/j,
 \qquad |\|v_{g_j,j}\|^2-d|<1/j.                         \tag{26}
\] Here \(d_{\rm BL}\) is the supremum over real test functions with
supremum norm and Lipschitz constant at most one. On finite measures on
\([0,\infty)\), weak convergence with tightness and mass convergence
implies (26): restrict tests to a compact interval containing all but a
chosen tail mass; their restrictions form a totally bounded family by
Arzela--Ascoli, so a finite uniform net reduces the assertion to
finitely many continuous test functions. No uniform-in-j semiclassical
error was inserted in this choice.

The comparison measures have j-independent weights in (24) with
\(\beta=1\), and their nth atom has energy \(2n\sigma_j/a_j\to0\) for
every fixed n.~Their summable tail weights are also j-independent.
Dominated convergence in this series, followed by (26), proves \[
 \nu_{g_j,j}\ \Longrightarrow\ d\,\delta_0,
 \qquad
 \langle v_{g_j,j},e^{-t(H_{g_j,j}-\mathcal E_{g_j,j})}v_{g_j,j}\rangle
       \longrightarrow d\quad(t\ge0).                   \tag{27}
\] Every finite vector remains exactly vacuum-orthogonal; every finite
Hamiltonian retains g\_j\textgreater0 and every original plaquette.
Nevertheless (27) is an atom at zero, not positive limiting spectral
weight throughout arbitrarily small intervals away from zero.

More precisely, suppose one tried to represent these limiting vectors
and correlations by a Hilbert vacuum \(\Omega\), a nonnegative
self-adjoint Hamiltonian K with \(K\Omega=0\), and a vector
\(v\perp\Omega\) of squared norm d.~Equality
\(\langle v,e^{-tK}v\rangle=d\) for t\textgreater0 implies
\(\int(1-e^{-t\omega})d\nu_v(\omega)=0\). Positivity makes \(\nu_v\)
supported at zero, hence Kv=0. Since d\textgreater0, that representation
has at least two independent zero-energy vectors. It cannot at the same
time have a unique vacuum. This is an exact consequence for this
retained observable/state limit, not a claim that all spatial limits of
(1) behave this way. Different states and local-observable
identifications require their own maps and spectral calculations.

The result pinpoints a concrete issue: finite gaps can close while a
retained macroscopic mode becomes a stationary zero-energy sector. To
test the desired interacting theory one must transport the actual local
gauge-invariant fields and their nonzero finite-energy spectral weights,
not infer them from that single vanishing gap. The original nonlinear
Hamiltonians, exact chart maps and the failure of uniformity are
explicit in this calculation; no interacting continuum has been replaced
by a claimed free-theory counterexample.

\subsection{Primary context}\label{weak-coupling-part-2-primary-context}

The exact finite-box quotient construction and eigenvalue comparison are
included in the companion source cited in Section 1. The semiclassical
localization method is B. Simon, \emph{Semiclassical analysis of low
lying eigenvalues. I. Non-degenerate minima: asymptotic expansions},
Annales de l'Institut Henri Poincare A 38 (1983), 295--308, Sections
2,3,6, https://www.numdam.org/item/AIHPA\_1983\_\_38\_3\_295\_0/. For
maximal-tree lattice coordinates see C. W. Bauer, I. D'Andrea, M.
Freytsis and D. M. Grabowska, \emph{A new basis for Hamiltonian SU(2)
simulations}, https://arxiv.org/abs/2307.11829. Their Hamiltonian
coefficients are not substituted for (1). The target requirement of a
nontrivial four-dimensional theory and its vacuum/spectral axioms is A.
Jaffe and E. Witten, \emph{Quantum Yang--Mills Theory}, Sections 3--4,
https://www.claymath.org/wp-content/uploads/2022/06/yangmills.pdf. The
oscillator and Gamma identities used here are proved in Sections 3,6; no
novelty is asserted for those classical formulas.

\clearpage
\part{Potential moments and the electric-state graph limit}
\section{Full-vacuum potential moments and the electric-state graph
limit}\label{weak-coupling-part-3-full-vacuum-potential-moments-and-the-electric-state-graph-limit}

8 September 2026. All statements through Section 6 concern a fixed
original open box L\textgreater=2 and fixed spacing a\textgreater0, at
small positive g. The nonlinear Wilson operator, actual vacuum, full
gauge action, raw state norms and all physical factors are retained. We
derive the missing derivative-state map from the vacuum equation, not
from strong vacuum convergence alone.

\subsection{1. Exact original operator and potential-gradient
inequality}\label{weak-coupling-part-3-exact-original-operator-and-potential-gradient-inequality}

Use the finite-box definitions and exact coordinate maps of the
companion manuscripts \emph{The original finite-box SU(2) Hamiltonian at
small positive coupling} and \emph{Actual weak-coupling vacuum,
observable and spectral maps}. In particular \[
 H=\kappa H_0+bW,\quad H_0=\sum_eE_e,
 \quad \kappa=2g^2/a,\quad b=1/(2g^2a),
 \quad W=\sum_p w_p,\quad w_p=2-\operatorname{tr}U_p,
 \quad E_e=-\sum_{\alpha=1}^3X_{e,\alpha}^2.                 \tag{1}
\] The anti-Hermitian generators are \(T_\alpha=-i\sigma_\alpha/2\). All
positive contained links and all contained faces are included. The
strictly positive unit vacuum is \(\psi\), with actual energy
\(\mathcal E\), so \(H\psi=\mathcal E\psi\). The scalar \(2b|P|\) is
part of bW and is not removed in this equation.

For any face p and any of its four links e, \[
 \sum_\alpha|X_{e,\alpha}w_p|^2=w_p-w_p^2/4\le w_p.       \tag{2}
\] To check every factor, hold the other links fixed. Its oriented face
word is a Haar translate of U\_e or U\_e\^{}\{-1\}; conjugations rotate
the three generators orthogonally. Write the resulting SU(2) element as
\(u_0I+i\boldsymbol u\cdot\boldsymbol\sigma\), with
\(u_0^2+|\boldsymbol u|^2=1\). Its trace is 2u\_0, and
\(\operatorname{tr}(T_\alpha U)=u_\alpha\). The three squared
derivatives sum to \(|\boldsymbol u|^2\). Since \(w_p=2-2u_0\), this is
exactly (2). The inverse-link case has the same sum by inversion and
bi-invariance, not a changed orientation.

Let \(r_e\) be the actual number of faces incident to e; on this open
box \(r_e\le4\), including every boundary edge. Cauchy--Schwarz in that
finite incident set, followed by summation in e, gives \[
 \sum_{e,\alpha}|X_{e,\alpha}W|^2
  \le\sum_e r_e\sum_{p\ni e}\sum_\alpha|X_{e,\alpha}w_p|^2
  \le16W.                                                   \tag{3}
\] Thus the constant 16 is independent of the box; W itself still
contains every face. In particular this is not a local replacement
potential.

\subsection{2. All potential moments from the actual vacuum
equation}\label{weak-coupling-part-3-all-potential-moments-from-the-actual-vacuum-equation}

Write \(M_k=\int W^k\psi^2\,dU\), with \(M_0=1\). For every integer
k\textgreater=1, \[
 \boxed{bM_{k+1}\le\mathcal E M_k+4\kappa k^2M_{k-1},
 \qquad M_1\le\mathcal E/b.}                              \tag{4}
\] Here is a proof that retains the derivative terms. The exact
ground-state identity for a real smooth f is \[
 q_H[f\psi]-\mathcal E\|f\psi\|^2
   =\kappa\int\psi^2\sum_{e,\alpha}|X_{e,\alpha}f|^2dU.
                                                               \tag{5}
\] It follows by differentiating the product twice in (1), integrating
with the original Haar measure, and substituting the full
eigen-equation. The kinetic term of \(q_H[f\psi]\) is nonnegative. For
k\textgreater=2 take \(f=W^{k/2}\), or its smooth positive
regularization if needed. Equation (3) bounds the right side by
\(4\kappa k^2M_{k-1}\); the left side is at least
\(bM_{k+1}-\mathcal E M_k\), proving (4). For k=1 use
\(f=(W+\epsilon)^{1/2}\). Its squared gradient is at most
\(4W/(W+\epsilon)\le4\). In (5) the lower bound is
\(b(M_2+\epsilon M_1)-\mathcal E(M_1+\epsilon)\). Letting epsilon
decrease to zero proves the same formula for k=1. For higher odd k the
identical regularization and dominated convergence on the compact group
justify all powers. The bound for M1 follows directly from \(bW\le H\)
in quadratic forms on \(\psi\).

The companion finite-box eigenvalue theorem gives a finite K such that
\(\mathcal E\le K\) for sufficiently small positive g, with L,a fixed.
No uniformity in L is assumed. Substituting the original coefficients
into (4) gives the exact scaled recursion \[
 \frac{M_{k+1}}{g^{2k+2}}
 \le 2aK\frac{M_k}{g^{2k}}+
                 16k^2\frac{M_{k-1}}{g^{2k-2}}.             \tag{6}
\] Define explicitly \(B_0=1\), \(B_1=2aK\), and
\(B_{k+1}=2aK B_k+16k^2 B_{k-1}\). Induction proves
\(M_k\le B_k g^{2k}\) for every fixed k. The recursion keeps the full
vacuum energy estimate and every original g,a factor. For example,
\(B_2=(2aK)^2+16\) and \(B_3=(2aK)^3+80(2aK)\). Only fixed-box
boundedness of K is used.

\subsection{3. Potential-weighted strong convergence and the total
electric
graph}\label{weak-coupling-part-3-potential-weighted-strong-convergence-and-the-total-electric-graph}

Retain the exact full-measure map \(\mathcal B_g\), matrices T,G,C, and
Gaussian \(\Phi_0\) from the state companion. Put \[
 V_2(x)=\tfrac14\sum_\alpha\|Cx^\alpha\|^2,
 \qquad D=-\sum_{\alpha,c,d}G_{cd}
                         \partial_{x_c^\alpha}\partial_{x_d^\alpha}.
                                                               \tag{7}
\] The companion proves \(\mathcal B_g\psi\to\Phi_0\) strongly. We
strengthen it to \[
 \mathcal B_g\bigl((W/g^2)\psi\bigr)\to V_2\Phi_0,
 \qquad
 \mathcal B_g(g^2H_0\psi)\to D\Phi_0
                \quad\hbox{strongly in }L^2.                \tag{8}
\] The first convergence is local by the exact Taylor limit W(gx)/g²→V2
and the strong vacuum convergence. To prove that no weighted tail is
lost, choose the same fixed radius rho as in the companion:
W(y)\textgreater=c\textbar y\textbar² inside it and W\textgreater=c\_rho
outside it. On the subset \textbar x\textbar\textgreater R with
\textbar gx\textbar\textless rho, the variable t=W/g² is at least cR².
Outside the rho neighborhood it is at least c\_rho/g². Therefore (6) at
k=3 gives \[
 \int_{|x|>R}|\mathcal B_g((W/g^2)\psi)|^2dx
  \le \frac{B_3}{cR^2}+\frac{g^2B_3}{c_\rho}.             \tag{9}
\] The limiting polynomial times Gaussian has an integrable squared
tail. Let R tend to infinity after the local convergence to obtain the
first part of (8).

For the second, the actual eigen-equation gives the exact identity \[
 g^2H_0\psi=\frac{a\mathcal E}{2}\psi-\frac{W}{4g^2}\psi.
                                                               \tag{10}
\] The first limit in (8) and \(\mathcal E\to\mu_0\) show that its image
tends to \((a\mu_0/2-V_2/4)\Phi_0\). The exact oscillator equation
\([(2/a)D+V_2/(2a)]\Phi_0=\mu_0\Phi_0\) identifies this as DΦ0. This
proves (8) from full potential moments, not merely from convergence in
L2 of the vacuum.

\subsection{4. Every retained electric weight and its graph
limit}\label{weak-coupling-part-3-every-retained-electric-weight-and-its-graph-limit}

For any fixed real edge weights w define \[
 \Gamma_w=\sum_e w_eE_e,\qquad
 G_w=T\operatorname{diag}(w)T^*,\qquad
 D_w=-\sum_{\alpha,c,d}(G_w)_{cd}
                        \partial_{x_c^\alpha}\partial_{x_d^\alpha}.
                                                               \tag{11}
\] No positivity of w is needed. The electric operators on distinct
links strongly commute by their product-group spectral decomposition.
Their eigenvalues are nonnegative, so on \(\operatorname{Dom}H_0\) \[
 \|\Gamma_w F\|\le\|w\|_\infty\|H_0F\|.                \tag{12}
\] Their individual Casimir spectra and all signs of w remain in this
inequality. Restriction to the physical space is valid because every
E\_e commutes with the full vertex gauge projection.

We claim the actual graph limit \[
 \boxed{\mathcal B_g(g^2\Gamma_w\psi)\longrightarrow D_w\Phi_0.}
                                                               \tag{13}
\] For a smooth invariant compactly supported u(x), its reciprocal map
\(\mathcal B_g^*u\) is a smooth physical function supported in the
logarithm chart for all sufficiently small g. Differentiating the exact
tree fields and density gives \[
 \mathcal B_g(g^2H_0\mathcal B_g^*u)\to Du,
 \qquad
 \mathcal B_g(g^2\Gamma_w\mathcal B_g^*u)\to D_wu          \tag{14}
\] strongly, since every coefficient and derivative converges uniformly
on the fixed compact support. In detail the principal coefficients of
gX\_e are its derivative columns of T at y=0; derivatives of these
coefficients and of the Haar density each carry an extra factor g.
Summing their exact squared fields gives G and Gw. The lower-order terms
tend to zero on this compact support, not by deletion from the original
operator.

Choose smooth invariant radial cutoffs u\_R=χ\_RΦ0. Their Gaussian decay
gives \(u_R\to\Phi_0\) in the graph norms of both constant-coefficient
operators D and Dw. The first limit in (14), combined with (8), shows \[
 \lim_{g\to0}\|g^2H_0(\psi-\mathcal B_g^*u_R)\|
        =\|D(\Phi_0-u_R)\|.                                \tag{15}
\] Apply (12) to the difference, then use the second limit in (14). The
limsup error in (13) is at most \(\|w\|_\infty\|D(\Phi_0-u_R)\|+
\|D_w(u_R-\Phi_0)\|\), which tends to zero as R grows. This proves (13).
Every limit here fixes L,a,w before g tends to zero.

Inner products with the strongly converging unit vacua now prove \[
 g^2\langle\Gamma_w\rangle_\psi\to
          \langle\Phi_0,D_w\Phi_0\rangle,
\quad
 \mathcal B_g\bigl(g^2(\Gamma_w-\langle\Gamma_w\rangle)\psi\bigr)
     \to V_w:=(D_w-\langle D_w\rangle)\Phi_0.              \tag{16}
\] This is the raw actual covariance-state map whose derivative control
was not supplied by (12) of the state companion alone.

\subsection{5. Explicit complete limiting covariance spectral
measure}\label{weak-coupling-part-3-explicit-complete-limiting-covariance-spectral-measure}

In the mode coordinates (4) of the state companion write \[
 \Sigma=\operatorname{diag}(\sigma_1,\ldots,\sigma_r),
 \quad A_w=O^TG^{-1/2}G_wG^{-1/2}O,
 \quad J_w=\Sigma^{1/2}A_w\Sigma^{1/2}.                    \tag{17}
\] These are exact finite real matrices with the original weights. Let
\(a_{\nu,\alpha}^{\dagger}\) be the oscillator creation operators fixed
by \(z_{\nu,\alpha}=\sqrt{2/\sigma_\nu}
(a_{\nu,\alpha}+a_{\nu,\alpha}^{\dagger})\), with
\([a_{\nu,\alpha},a_{\eta,\beta}^{\dagger}]
=\delta_{\nu\eta}\delta_{\alpha\beta}\). In particular each real
coordinate has vacuum variance 2/σν. Differentiating the full Gaussian
of the companion gives \[
 \langle D_w\rangle=\frac38\operatorname{tr}(A_w\Sigma),
 \qquad
 V_w=-\frac18\sum_{\nu,\eta}(J_w)_{\nu\eta}
                \sum_{\alpha=1}^3
                   a_{\nu,\alpha}^{\dagger}a_{\eta,\alpha}^{\dagger}\Phi_0.
                                                               \tag{18}
\] For a direct check before passing to creation operators, put
\(S=G^{-1/2}B^{1/2}G^{-1/2}\). Ordinary differentiation gives \[
 D_w\Phi_0=
 \left[\frac34\operatorname{tr}(G_wS)
  -\frac1{16}\sum_\alpha(x^\alpha)^TSG_wSx^\alpha\right]\Phi_0.
                                                               \tag{19}
\] Gaussian covariance 2S\^{}-1 subtracts exactly half of its constant
term. Then (17) and the stated coordinate/creation relation give (18),
with no omitted constant or color multiplicity.

The diagonal two-quantum vector
\(\sum_\alpha(a_{\nu,\alpha}^{\dagger})^2\Phi_0\) has squared norm 6;
the off-diagonal vector
\(\sum_\alpha a_{\nu,\alpha}^{\dagger}a_{\eta,\alpha}^{\dagger}\Phi_0\)
for ν\textless η has squared norm 3. These follow by commuting the two
annihilators through the two creators; different unordered mode pairs
are orthogonal. The exact limiting raw excitation measure is
consequently \[
 \boxed{\nu_{w,0}=
 \frac3{32}\sum_\nu (J_w)_{\nu\nu}^2\delta_{2\sigma_\nu/a}
 +\frac3{16}\sum_{\nu<\eta}(J_w)_{\nu\eta}^2
                          \delta_{(\sigma_\nu+\sigma_\eta)/a}.}
                                                               \tag{20}
\] Repeated energies add their positive masses. The total is
\(3\operatorname{tr}(J_w^2)/32\); its first moment is
\(3\operatorname{tr}(\Sigma J_w^2)/(16a)\). These are moments of the
displayed limiting finite measure; convergence of an unbounded
first-moment test for finite-g measures is not inferred merely from weak
convergence.

For the actual full-model measures the precise conclusion is \[
 g^4\langle v_w,\mathbf1_{\bullet}(H-\mathcal E)v_w\rangle
       \Longrightarrow\nu_{w,0},\quad
 v_w=(\Gamma_w-\langle\Gamma_w\rangle)\psi,
 \qquad
 g^4\|v_w\|^2\to\frac3{32}\operatorname{tr}(J_w^2).        \tag{21}
\] This follows from the proved strong vector map (16), actual finite
spectral projections of the state companion, and its tightness argument
for raw measures. It holds against bounded continuous tests, with total
mass convergence. The factor g4 is explicitly retained; the original
norm and original positive-g spectrum are not declared equal to the
limit.

\subsection{6. Native weights and the actual state, without a fictitious
vacuum}\label{weak-coupling-part-3-native-weights-and-the-actual-state-without-a-fictitious-vacuum}

For the original native electric weight, \(w_{(n,1)}=n_2^2\) and all
other w\_e=0. Then \[
 G_{\rm nat}=\sum_{e=(n,1)}n_2^2\,t_et_e^T,
 \qquad t_e\text{ is column e of }T.                       \tag{22}
\] It is positive semidefinite and nonzero. For example take a
direction-1 edge at transverse coordinate n2=L and choose an adjacent
contained face. That edge participates in a chord loop, and its column
of T is nonzero: otherwise d1=CT would have zero column there,
contradicting the nonzero coefficient of that face in the original
oriented curl. Its weight L² is positive. Thus (22) is nonzero, and all
invertible congruences in (17) preserve that property. Hence \[
 \lim_{g\downarrow0}g^4\|v_{\rm nat}\|^2
              =\frac3{32}\operatorname{tr}(J_{\rm nat}^2)>0.\tag{23}
\] This resolves the fixed-box derivative/norm question for the actual
native Γ state. Equation (20) supplies all its limiting mode-pair
masses, including their original physical energies. It does not
establish a volume-uniform error in (21) or spectral tightness along
g=g\_j,L=L\_j. The exact finite matrices (22),(17) provide the next
explicit spatial objects to compute, rather than an assumption that the
missing joint limit has already been proved.

Equations (13)--(23) concern the actual electric derivative state. They
do not identify a finite-angle translated state with that derivative
state. In particular no small-angle or simultaneous volume/coupling
remainder has been imported into (21) without its own estimate.

\subsection{7. Finite-energy mass of the native comparison measure as
the box
grows}\label{weak-coupling-part-3-finite-energy-mass-of-the-native-comparison-measure-as-the-box-grows}

The matrices in (17),(22) also permit a boundary-sensitive estimate of
the limiting native measure, without diagonalizing a large box
numerically. Let \(R=T^*G^{-1/2}\), an isometry onto \(\ker d_0^*\), and
\(P_T=RR^*\). Then \[
 A_{\rm nat}=O^TR^*\operatorname{diag}(w)RO,\quad
 0\le A_{\rm nat}\le L^2I,\quad
 \operatorname{tr}A_{\rm nat}=\sum_e w_e(P_T)_{ee}.         \tag{24}
\] Every original edge belongs to a contained elementary face, including
boundary edges. The oriented circulation around that face, divided by 2,
is a unit vector in \(\ker d_0^*\) with edge component of absolute value
1/2. The variational formula for the norm of an orthogonal projection
therefore gives \((P_T)_{ee}\ge1/4\). For the original native weights
the exact sum is \[
 S_1:=\sum_e w_e
   =(2L)(2L+1)\sum_{n_2=-L}^L n_2^2
   =\frac23L^2(L+1)(2L+1)^2.                               \tag{25}
\] The least frequency obeys \(\sigma_{\min}=\sqrt8\sin(\pi/(4L+2))
\ge\sqrt8/(2L+1)\), using \(\sin x\ge2x/\pi\) on \([0,\pi/2]\).
Positivity and the trace Cauchy--Schwarz inequality, with
\(r=4L^2(4L+3)\), imply \[
 \begin{split}
 \operatorname{tr}(J_{\rm nat}^2)
 &\ge\frac{(\operatorname{tr}J_{\rm nat})^2}{r}
 \ge\frac{\sigma_{\min}^2 S_1^2}{16r}\\
 &\ge\frac{L^2(L+1)^2(2L+1)^2}{18(4L+3)}
 \ge\frac{L^5}{27}.                                      \tag{26}
 \end{split}
\] The last, deliberately explicit lower estimate uses \(4L+3\le6L\) for
L\textgreater=2; neither the earlier exact expression nor its boundary
factors is replaced in the definition of the state.

Fix a physical energy ceiling \(E_*>0\) and put \(s=aE_*\). Only mode
pairs with both \(\sigma_\nu\le s\) and \(\sigma_\eta\le s\) can occur
in (20) below that ceiling. Let P\_s be the diagonal mode projection
onto these frequencies. Equations (17),(24) give \[
 \|P_sJ_{\rm nat}P_s\|\le sL^2,\qquad
 \nu_{{\rm nat},0}([0,E_*])
       \le\tfrac3{32}\operatorname{rank}(P_s)s^2L^4.       \tag{27}
\] The exact open-box frequency list is
\(\sigma^2=\sum_{i=1}^3 4\sin^2(\pi j_i/[2(2L+1)])\), \(0\le j_i\le2L\),
with multiplicity k-1 when exactly k\textgreater=2 indices are positive.
The same sine inequality implies \(j_i\le(2L+1)s/2\) for every counted
triple. Each multiplicity is at most 2. Even including the extra
excluded zero/one-positive triples therefore gives the valid bound \[
 \operatorname{rank}(P_s)
   \le2\left(1+\left\lfloor\frac{(2L+1)s}{2}\right\rfloor\right)^3.
                                                               \tag{28}
\] Combining (26)--(28) with the exact total raw mass in (21) yields \[
 \boxed{\nu_{{\rm nat},0}([0,E_*])
 \le \min\!\left\{1,
 \frac{54(aE_*)^2}{L}
       \left(1+\frac{(2L+1)aE_*}{2}\right)^3\right\}
                          \nu_{{\rm nat},0}([0,\infty)).} \tag{29}
\] This retains the original raw mass; it is not an unrecorded change of
state norm. On the original spatial diagonal L\_j=j²,a\_j=1/(100j), its
explicit relative bound is \[
 \frac{54E_*^2}{10^4 j^4}
          \left(1+\frac{(2j^2+1)E_*}{200j}\right)^3
                   \longrightarrow0.                     \tag{30}
\] Thus the weak-coupling native \emph{comparison} measures send a
vanishing fraction of their nonzero raw mass to every fixed finite
physical-energy interval along this box/spacing diagonal. This statement
is fully proved for the explicit matrix measures; it has not assumed a
uniform remainder for the actual finite-g measures in (21).

There is a corresponding precise existence statement for actual
positive-g regulators. At each fixed j, divide both sides of (21) by
their explicitly retained nonzero total mass to compare probability
measures. Their weak convergence follows from (21). For each integer
k\textgreater=1 choose once a number R\_k in (k,k+1) outside the
countable union, over all j\textgreater=2, of the finite comparison
spectra. At stage j, fixed-box convergence gives a positive
g\_j\textless1/j for which the actual and comparison fractions below
each R\_k, k\textless=j, differ by less than 1/j. Only finitely many
inequalities are imposed at that stage. For fixed E\_* choose k with
R\_k\textgreater E\_\emph{. Eventually j\textgreater=k; the bound (30)
at the fixed ceiling R\_k and the error 1/j show that the actual
fraction below E\_} tends to zero. This proves existence of an actual
all-positive-g native electric-state sequence escaping every fixed
finite energy interval. It gives no specified logarithmic/fixed-kappa
running-coupling estimate.

The macroscopic radial observable of the state companion and this native
electric observable therefore have different, explicitly computed
spectral limits on sufficiently weak chosen diagonals: a zero-energy
atom for the former, finite-energy escape for the latter. Their common
coordinate maps, actual vacuum and full Hamiltonian are retained. This
does not prove they are unrelated, nor rule out a further state family
or spectral selection with a useful interacting continuum limit.

\subsection{8. Exact Hamiltonian filtering, including the zero-atom
distinction}\label{weak-coupling-part-3-exact-hamiltonian-filtering-including-the-zero-atom-distinction}

For any finite regulator and smooth vacuum-orthogonal \(v\), put
\(A=H-\mathcal E\) and \[
 u_t=e^{-tA/2}v,\quad
 C(t)=\|u_t\|^2=\int_{[0,\infty)}e^{-t\omega}\,d\nu_v(\omega).
\tag{31}
\] This is the raw spectral measure, with no normalization by a possibly
vanishing norm. Since the vectors used above lie in the operator domain,
spectral differentiation for \(t>0\) gives \[
 -C'(t)=\int\omega e^{-t\omega}d\nu_v(\omega),\quad
 m(t):=-C'(t)/C(t),\quad
 m'(t)=-\operatorname{Var}_{\nu_{v,t}}(\omega)\le0,
\tag{32}
\] where \(d\nu_{v,t}=e^{-t\omega}d\nu_v/C(t)\). Euclidean-time
filtering therefore lowers the exact spectral mean while retaining the
raw amplitude \(C(t)\).

For a varying regulator, Markov's inequality gives the exact criterion
\[
 \nu_{v_j}([0,R])\ge C_j(t_j)\left(1-\frac{m_j(t_j)}R\right),
 \qquad R>0.\tag{33}
\] Equation (33) is a sufficient estimate on raw low-energy mass, not a
necessary criterion for gaplessness. A previous version incorrectly
called a positive lower bound on \(C_j(t_j)\) necessary. The following
exact calculation corrects that statement.

Take \(k=3/2\), \(c=100\sqrt2\pi\), and the finite measure \[
 d\gamma_c(\omega)=
 \frac{\omega^{k-1}e^{-\omega/c}}{\Gamma(k)c^k}\,d\omega.
\] On \(\mathbb C\Omega\oplus L^2(\gamma_c)\), let
\(K=0\oplus M_\omega\), where multiplication by \(\omega\) has domain
\(\{f:\int\omega^2|f|^2d\gamma_c<\infty\}\). This real multiplication
operator is self-adjoint: its resolvent at nonreal \(z\) is
multiplication by \((\omega-z)^{-1}\), and its adjoint has the same
maximal domain. Absolute continuity gives \(\ker K=\mathbb C\Omega\).
For each \(\varepsilon>0\), \((0,\mathbf1_{(0,\varepsilon)})\) is a
nonzero vacuum-orthogonal vector of energy quotient at most
\(\varepsilon\). Consequently this operator has no positive gap. For its
fixed raw vector \(v=(0,1)\), direct Gamma integration gives \[
 C(t)=(1+ct)^{-k},\qquad
 m(t)=\frac{kc}{1+ct},\qquad
 m'(t)=-\frac{kc^2}{(1+ct)^2}.
\] Thus both \(C(t)\to0\) and \(m(t)\to0\). No positive lower bound on
the filtered norm is necessary.

In fact, suppose finite raw measures \(\nu_j\) converge weakly to a
finite \(\nu\), \(\liminf C_j(t_j)\ge\eta>0\), and \(m_j(t_j)\to0\). For
every \(R>0\), (33) gives \(\liminf\nu_j([0,R])\ge\eta\). Since
\([0,R]\) is closed, the portmanteau inequality then gives
\(\nu([0,R])\ge\eta\). Continuity from above as \(R\downarrow0\) proves
\(\nu(\{0\})\ge\eta\). Under finite weak convergence the proposed pair
of bounds therefore forces a positive zero atom, rather than detecting
only positive energies arbitrarily close to zero.

For a fixed positive Hamiltonian on a fixed Hilbert space, a family of
nonzero vectors orthogonal to its entire kernel with Rayleigh quotients
tending to zero does contradict a positive spectral gap, regardless of
the raw norms: a gap \(\Delta>0\) would imply
\(\langle v,Kv\rangle\ge\Delta\|v\|^2\) for every such vector. For
varying regulators, the additional issue is to construct the actual
common state and observable limit identifying the continuum theory. It
is not a defect in that variational inequality.

The fixed-parameter radial family in Section 7 of the companion
\emph{Actual weak-coupling vacuum, observable and spectral maps} has
\(C_j(t)\to1-2^{-3/2}\) and \(m_j(t)\to0\) at fixed physical \(t>0\), so
it retains an extra zero sector. Differentiation of its limit at such
positive times is justified because \(\omega e^{-t\omega}\) is bounded.
The native electric family here has the vanishing finite-energy fraction
(29)--(30). The growing-parameter radial family in the separate Gamma
note instead has the atomless measure displayed above. These are three
explicit spectral outcomes of specified states, not a classification of
all possible state maps.

\subsection{Reproducibility and
references}\label{weak-coupling-part-3-reproducibility-and-references}

Every operator and Hilbert-space map used here is given explicitly or in
the two complete companion proofs included with this document. The
gradient constant, moment recursion, graph limit and all spectral
coefficients are derived above. The standard ground-state transform and
localization context are cited in the companions (Kogut--Susskind,
Bauer--D'Andrea--Freytsis-- Grabowska, and Simon). No interacting
continuum construction or mass-gap counterexample is asserted. In
particular this is not a proof that the original S6 completion or any
claimed external counterexample exists.

\clearpage
\part{Radial states and the full Gamma spectral limit}
\section{A radial state with a non-atomic low-energy scaling
limit}\label{weak-coupling-part-4-a-radial-state-with-a-non-atomic-low-energy-scaling-limit}

8 September 2026. This note constructs an exact finite-regulator state
sequence for the unchanged open-box SU(2) Hamiltonian. Unlike the
fixed-\(\beta\) lowest-mode diagonal, whose mass collapses to a zero
atom, the radial parameter grows with the box. The raw spectral measures
converge to an explicit Gamma density supported down to zero. Every
finite regulator has positive coupling, the full Wilson potential, the
actual positive vacuum, and the exact physical projection.

\subsection{1. The finite-regulator
state}\label{weak-coupling-part-4-the-finite-regulator-state}

Use the operator, maximal-tree quotient, physical Hilbert space, actual
ground state \(\psi_{g,L,a}\), and chart observable of the companion
proofs. The original operator, with no Wilson term removed, is \[
 H_{g,L,a}=\frac{2g^2}{a}\sum_{e\in E}
 \left(-\sum_{\alpha=1}^3X_{e,\alpha}^2\right)
 +\frac1{2g^2a}\sum_{p\in P}(2-\operatorname{tr}U_p),
 \qquad T_\alpha=-i\sigma_\alpha/2.
\] Here \(E\) and \(P\) contain all positive edges and contained faces
in the open box with vertices \(\{-L,\ldots,L\}^3\), and the Hilbert
measure is product Haar restricted to all vertex-gauge-invariant
functions. Its unique positive unit vacuum and excitation energy are
denoted by \(\psi_{g,L,a}\) and \(\mathcal E_{g,L,a}\). The additive
tree quotient retains \(G=TT^*>0\), chord insertion \(\iota\), and face
curl \(C=d_1\iota\). Choose \(O\) with
\(O^{\mathsf T}G^{1/2}C^*CG^{1/2}O=\operatorname{diag}(\sigma_\nu^2)\).
These are the actual kinetic and curl tensors, not a substituted metric.
For integers \(j\ge2\), \(L_j=j^2\), \(a_j=1/(100j)\), choose a lowest
physical oscillator mode with \[
 \sigma_j=\sqrt8\sin\frac{\pi}{4j^2+2},\qquad
 \beta_j=\sqrt j,\qquad \theta_j=\frac{\beta_j\sigma_j}{4}.
 \tag{1}
\] For each fixed \(j\), use the companion's globally smooth
chart-cutoff multiplier \[
 B_{g,j}(Z)=
 \exp\!\left(i\theta_j\chi(y)
 \left|(O^{\mathsf T}G^{-1/2}(y/g))_{\nu_j}\right|^2\right)
 \tag{2}
\] in the logarithm chart and \(1\) outside it. It is invariant under
the residual simultaneous colour action. Define the actual centered
vector and its raw positive spectral measure: \[
 v_{g,j}=(B_{g,j}-c_{g,j})\psi_{g,j},\quad
 c_{g,j}=\langle\psi_{g,j},B_{g,j}\psi_{g,j}\rangle,
\] \[
 \nu_{g,j}(J)=
 \left\langle v_{g,j},
 \mathbf1_J(H_{g,j}-\mathcal E_{g,j})v_{g,j}\right\rangle .
 \tag{3}
\] The cutoff \(\chi\) is real, smooth, invariant under simultaneous
colour rotations, equal to one near zero and supported strictly inside
the product logarithm chart. Thus (2) extends smoothly by one. Since
\(O,G\) act only on spatial chord indices, (2) is exactly gauge
invariant. It has modulus one, so the raw squared norm in (3) is
\(1-|c_{g,j}|^2\). The multiplier is nonconstant on an open set of the
selected three coordinates. Positivity of \(\psi_{g,j}\) and equality
conditions in Cauchy--Schwarz give \(|c_{g,j}|<1\); hence the vector is
nonzero for every \(g>0\). Subtracting \(c_{g,j}\) removes its entire
vacuum component. Its finite-regulator domain is smooth functions on the
compact quotient, contained in the domain of the full Hamiltonian.

The fixed-box theorem permits a strictly positive \(g_j<1/j\) matching
the actual raw measure to its comparison measure. Section 3 gives the
exact selection, without replacing a state by a unit-norm surrogate.

\subsection{2. Exact comparison measure}\label{weak-coupling-part-4-exact-comparison-measure}

The companion Laguerre calculation applies with \(k=3/2\) and
\(\beta=\beta_j\). Its raw comparison measure is \[
 \nu_{0,j}=
 \sum_{n=1}^{\infty}
 \frac{(k)_n}{n!}
 \frac{\beta_j^{2n}}{(1+\beta_j^2)^{n+k}}\,
 \delta_{\,2n\sigma_j/a_j}.
 \tag{4}
\] Its exact total mass is \[
 d_j=1-(1+\beta_j^2)^{-k}
     =1-(j+1)^{-3/2}\longrightarrow1.                 \tag{5}
\] The \(n=0\) term is absent because the state is vacuum-orthogonal; it
is not silently folded into the raw measure.

For \(t\ge0\), summing the negative-binomial series gives \[
 L_j(t):=\int e^{-t\omega}\,d\nu_{0,j}(\omega)
 =
 \left[1+j\left(1-e^{-t\lambda_j}\right)\right]^{-k}
 -(j+1)^{-k},
 \quad \lambda_j=\frac{2\sigma_j}{a_j}.                 \tag{6}
\] The physical spacing and frequency give \[
 j\lambda_j
 =400\sqrt2\,j^2\sin\frac{\pi}{4j^2+2}
 \longrightarrow c:=100\sqrt2\,\pi .                  \tag{7}
\] Hence, for each fixed \(t\), \[
 L_j(t)\longrightarrow(1+ct)^{-3/2}.                  \tag{8}
\] This is the Laplace transform of \[
 d\gamma_c(\omega)=
 \frac{\omega^{1/2}e^{-\omega/c}}
 {\Gamma(3/2)c^{3/2}}\,d\omega,\qquad \omega\ge0.       \tag{9}
\] Here is a direct quantitative proof of raw-measure convergence. Let
\(Q\) have density \(q^{k-1}e^{-q}/\Gamma(k)\) on \(q>0\), and, for each
\(j\), let \(N_j\), conditional on \(Q=q\), be Poisson with mean \(jq\).
Integration, with the original parameter \(j\) retained, gives \[
 \begin{split}
 \mathbb P(N_j=n)
 &=\frac{j^n}{n!\Gamma(k)}
     \int_0^\infty q^{n+k-1}e^{-(j+1)q}\,dq\\
 &=\frac{(k)_n}{n!}\frac{j^n}{(j+1)^{n+k}} .
 \end{split}
\] Consequently the law of \(\lambda_jN_j\) is exactly
\(\nu_{0,j}+(j+1)^{-k}\delta_0\). This is an auxiliary probability
coupling; the original state measure still excludes the zero
coefficient. Put \(c_j=j\lambda_j\). Conditional expectation makes the
mixed term in the following square vanish, and the Poisson conditional
variance is \(jQ\). Since \(\mathbb EQ=k\), \(\mathbb EQ^2=k(k+1)\), \[
 \begin{split}
 \mathbb E|\lambda_jN_j-cQ|^2
 &=\mathbb E|c_j(N_j/j-Q)+(c_j-c)Q|^2\\
 &=\frac{kc_j^2}{j}+k(k+1)(c_j-c)^2 .
 \end{split}
\] For \(x_j=\pi/(4j^2+2)\), the elementary inequalities
\(0\le x-\sin x\le x^3/6\) give the explicit retained-scale estimate \[
 0<c-c_j\le
 \frac{c}{2j^2+1}
 +\frac{400\sqrt2\,\pi^3j^2}{6(4j^2+2)^3}.
\] Indeed \(c-400\sqrt2j^2x_j=c/(2j^2+1)\). The sine inequality follows
by integrating \(1-\cos s\le s^2/2\), which itself follows by
integrating \(\sin u\le u\). Thus no asymptotic constant is left
unspecified.

Define \(d_{\rm BL}\) on finite measures using all tests \(f\) with
\(\|f\|_\infty\le1\) and Lipschitz constant at most one. The coupling,
Cauchy--Schwarz, and the removed zero coefficient now prove \[
 d_{\rm BL}(\nu_{0,j},\gamma_c)
 \le
 \left[\frac{kc_j^2}{j}+k(k+1)(c_j-c)^2\right]^{1/2}
 +(j+1)^{-k}\longrightarrow0 .
\] The law of \(cQ\) is (9), by the explicit change of variable
\(\omega=cq\). This proves \[
 \nu_{0,j}\Longrightarrow\gamma_c.                    \tag{10}
\] In particular \(\gamma_c((0,\varepsilon))>0\) for every
\(\varepsilon>0\), while \(\gamma_c(\{0\})=0\).

\subsection{3. Transport to the actual nonlinear
Hamiltonian}\label{weak-coupling-part-4-transport-to-the-actual-nonlinear-hamiltonian}

For each fixed \(j\), the actual-vacuum and finite-spectral-projection
theorem for the full nonlinear Hamiltonian gives weak convergence of the
actual raw measure to (4) as \(g\downarrow0\), with total-mass
convergence. Choose \(0<g_j<1/j\) so that \[
 d_{\mathrm{BL}}(\nu_{g_j,j},\nu_{0,j})<\frac1j,\qquad
 \left|\|v_{g_j,j}\|^2-d_j\right|<\frac1j.               \tag{11}
\] For a definite selection let \(n_j\) be the least positive integer
such that \(g_j=2^{-n_j}\) satisfies these inequalities. Existence
follows because both inequalities hold at every sufficiently small
positive coupling for that fixed \(j\). This uses no volume-uniform
estimate and provides no explicit renormalized running-coupling formula.
All coefficients remain \(\kappa_j=200jg_j^2\), \(b_j=50j/g_j^2\),
\(\xi_j=1/(4g_j^4)\), including the scalar \(2b_j|P_j|\). Combining
(10), (11) and the explicit coupling estimate proves \[
 \nu_{g_j,j}\Longrightarrow\gamma_c,\qquad
 \|v_{g_j,j}\|^2\longrightarrow1.                       \tag{12}
\] The raw norm and raw amplitudes remain explicit; no state rescaling
is used in (11)--(13). Since the Gamma law has no endpoint atoms,
portmanteau gives \[
 \lim_{j\to\infty}
 \nu_{g_j,j}((0,\varepsilon))
 =\gamma_c((0,\varepsilon))>0
 \qquad(\varepsilon>0).                                \tag{13}
\] Every \(v_{g_j,j}\) is a nonzero smooth gauge-invariant vector
orthogonal to the actual finite-volume vacuum, and every \(g_j\) is
strictly positive. Thus (12)--(13) concern the full nonlinear
Hamiltonian at every selected positive coupling. The limiting comparison
calculation does not assert that interactions survive as nontrivial
continuum dynamics.

For \(t>0\) and any nonnegative integer \(r\), the test
\(\omega^r e^{-t\omega}\) is bounded and continuous. Consequently the
actual correlations and all their derivatives at positive times
converge: \[
 \int\omega^r e^{-t\omega}\,d\nu_{g_j,j}
 \longrightarrow
 (k)_r c^r(1+ct)^{-k-r}.
\] This does not assert convergence of unbounded energy moments at
\(t=0\).

\subsection{4. The cyclic Hilbert-space map determined by these
states}\label{weak-coupling-part-4-the-cyclic-hilbert-space-map-determined-by-these-states}

Let \(\mathcal H_\gamma=L^2((0,\infty),\gamma_c)\), \(v=1\), and
\(M_\omega f=\omega f\) on \(\{f:\int\omega^2|f|^2d\gamma_c<\infty\}\).
Real multiplication on this maximal domain is self-adjoint: its nonreal
resolvent is the bounded multiplier \((\omega-z)^{-1}\); testing the
adjoint on compactly supported functions gives the same maximal domain.
Its spectral projections are multiplication by indicators. In particular
its vector \(v\) has exactly the raw measure (9).

For finitely many coefficients \(a_\ell\) and times \(t_\ell\ge0\), put
\(A_j=H_{g_j,j}-\mathcal E_{g_j,j}\). Expanding the scalar products
yields \[
 \left\|\sum_\ell a_\ell e^{-t_\ell A_j}v_{g_j,j}\right\|^2
 \longrightarrow
 \int\left|\sum_\ell a_\ell e^{-t_\ell\omega}\right|^2d\gamma_c .
\] This specifies an exact isometry from the limiting time-orbit Gram
space to \(\mathcal H_\gamma\). The image is dense: if \(h\) is
orthogonal to all exponentials, then the finite measure
\(\overline h\,d\gamma_c\) has zero Laplace transform. Push it by
\(x=e^{-\omega}\) to \((0,1]\) and extend it to \([0,1]\) with zero mass
at zero. Testing integer times shows all its polynomial moments vanish.
Polynomials are uniformly dense in continuous functions on \([0,1]\), so
the measure is zero and \(h=0\). The semigroup is intertwined with
multiplication by \(e^{-t\omega}\); the time and energy scales have not
been changed.

Adding \(\mathbb C\Omega\) with energy zero retains the finite-vacuum
Gram vectors orthogonal to this cyclic sector. The resulting particular
spectral representation has a unique vacuum and positive spectrum down
to zero. It is not yet a representation of the continuum local gauge
observable algebra.

\subsection{5. Exact scope}\label{weak-coupling-part-4-exact-scope}

This proves a non-atomic low-energy scaling measure for one explicitly
selected sequence of finite nonlinear lattice Hamiltonians. It does not
yet prove that the varying Hilbert spaces reconstruct the prescribed
nontrivial Yang--Mills theory on \(\mathbb R^4\), or that this selected
bare coupling sequence has its renormalized physical scale. Those are
additional state-identification, reconstruction, and nontriviality maps.
The result supplied here is the exact state, raw measure, Gamma limit,
and positive mass in every interval \((0,\varepsilon)\); no free-field
replacement is used.

The subsequent full amplitude/time proof constructs the precise larger
representation containing these states. A nonzero carrier
\(b_0+\alpha/\sqrt j\) retains the kernel
\([1+cb_0^2t+i(\alpha-\delta)]^{-3/2}\), but also applying the common
modulation to the vacuum retains an entire additional zero-energy origin
fibre. The configuration-bridge proof further shows that on a compatible
selected diagonal the unchanged bounded local cylinder algebra acts by
scalars on the complete Gamma time orbit. Thus the unique-vacuum
statement for the restricted cyclic representation of Section 4 must not
be read as a unique-vacuum claim for the full retained modulation
algebra or for a local field theory.

\subsection{6. Reproducibility and
literature}\label{weak-coupling-part-4-reproducibility-and-literature}

The exact series coefficients and moment identities are algebraic
checks; the Gamma integral, sine estimate and convergence are proved
analytically above. The supplied checker separates exact symbolic
identities from floating-point illustrations; neither certifies the
nonlinear companion theorem. That theorem is proved in the included
finite-box and state-map companions, which are necessary proof inputs.

For the continuum target, A. Jaffe and E. Witten, \emph{Quantum
Yang--Mills Theory}, Section 2, specifies a nontrivial four-dimensional
quantum field theory, not only an excitation operator and a cyclic
vector
(\href{https://www.claymath.org/wp-content/uploads/2022/06/yangmills.pdf}{Clay
problem text}). S. Chatterjee, \emph{Yang--Mills for probabilists},
Sections 3 and 6--7, explains the continuum loop-observable construction
and its scaling
(\href{https://arxiv.org/abs/1803.01950}{arXiv:1803.01950}). Those works
identify which additional mathematical structures must be constructed;
they are not invoked as premises establishing this limit.

\clearpage
\part{Exact amplitude-time kernels and carrier windows}
\section{Exact radial amplitude/time kernel and carrier-window
representation}\label{weak-coupling-part-5-exact-radial-amplitudetime-kernel-and-carrier-window-representation}

8 September 2026. This note computes the complete complex amplitude/time
kernel of the radial spectral family. It proves the continuous
modulation limit about a nonzero carrier, its exact Hamiltonian, and its
relation to the globally discontinuous amplitude family and the origin
zero sector.

\subsection{Original objects and
conventions}\label{weak-coupling-part-5-original-objects-and-conventions}

All inner products are conjugate linear in their first argument. Retain

\[
 L_j=j^2,\quad a_j=\frac1{100j},\quad
 \sigma_j=\sqrt8\sin\frac{\pi}{4j^2+2},\quad
 \lambda_j=\frac{2\sigma_j}{a_j}
 =400\sqrt2\,j\sin\frac{\pi}{4j^2+2},\quad j\ge2.
\]

Thus \(j\lambda_j\to c=100\sqrt2\pi\), with physical time t and every
spacing/frequency factor unchanged. Put \(k=3/2\), \(s_j=\sqrt j\), and
\(q=\sigma_j|z|^2/4\). In the exact oscillator radial Hilbert map, the
vacuum is the constant function one in

\[
 \mathcal H_\rho=L^2([0,\infty),\rho),\qquad
 d\rho(q)=\frac{q^{k-1}e^{-q}}{\Gamma(k)}dq.
\]

The normalized Laguerre vectors
\(\ell_n(q)=\sqrt{n!/(k)_n}L_n^{k-1}(q)\) form an orthonormal basis; the
excitation operator is \(\lambda_j N\), where \(N\ell_n=n\ell_n\). This
is the exact radial spectral-coordinate map of the comparison operator,
not replacement of the finite-g operator. The actual nonlinear
Hamiltonian, Haar isometry, physical projection, and fixed-box transport
are those proved in the included companions \emph{The original
finite-box SU(2) Hamiltonian at small positive coupling} and
\emph{Actual weak-coupling vacuum, observable and spectral maps}.

Define uncentered and centered radial vectors

\[
 \phi_j(b)=e^{ib s_jq},\quad
 m_j(b)=\langle1,\phi_j(b)\rangle=(1-ib s_j)^{-k},\quad
 v_j(b)=\phi_j(b)-m_j(b),\qquad b\in\mathbb R.
\]

Their raw centered norm is \(1-(1+b^2j)^{-k}\). In particular v\_j(0)=0
exactly. No division by this norm is used in the kernel calculation.

\subsection{1. The finite-j kernel, including complex powers and
subtraction}\label{weak-coupling-part-5-the-finite-j-kernel-including-complex-powers-and-subtraction}

The exact coefficients are

\[
 a_n(b)=\langle\ell_n,\phi_j(b)\rangle
 =\sqrt{\frac{(k)_n}{n!}}
       \frac{(-ib s_j)^n}{(1-ib s_j)^{n+k}}.
\]

Let

\[
 P_j(b,d)=(1+ib s_j)(1-id s_j)
          =1+bdj+i(b-d)s_j.
\]

The factors 1+ib s\_j and 1-id s\_j lie in the open right half-plane.
Their principal arguments add to
\(\arctan(bs_j)-\arctan(ds_j)\in(-\pi,\pi)\). Consequently their
principal logarithms add to the principal logarithm of P\_j, including
when bd\textless0. Therefore

\[
 \overline{a_n(b)}a_n(d)
 =\frac{(k)_n}{n!}\frac{(bdj)^n}{P_j(b,d)^{n+k}},
 \quad
 \overline{m_j(b)}m_j(d)=P_j(b,d)^{-k}.
\]

For t\textgreater=0 the binomial-series argument has modulus strictly
less than one:

\[
 \left|\frac{bdj\,e^{-t\lambda_j}}{P_j(b,d)}\right|
 \le\frac{|bd|j}{\sqrt{(1+b^2j)(1+d^2j)}}<1.
\]

If b or d is zero the argument is zero and the same formulas apply.
Summing first from n=0 gives the uncentered kernel

\[
 \begin{split}
 F_j(b,d;t)
 &:=\langle\phi_j(b),e^{-t\lambda_jN}\phi_j(d)\rangle\\
 &=\left[1+bdj(1-e^{-t\lambda_j})+i(b-d)\sqrt j\right]^{-k}.
 \tag{1}
 \end{split}
\]

The power in (1) is the principal one. To verify that no phase was lost
when combining the factors, let z run from 0 to e\^{}\{-t lambda\_j\}
and put E(z)=P\_j-bdjz. When b=d, both P\_j and E(z) are positive real.
When b!=d, E(z) stays in one open half-plane because its nonzero
imaginary part is constant. It never crosses the principal logarithm
cut. Also 1-bdjz/P\_j lies in the open right half-plane by the strict
modulus bound. At z=0 the logarithmic factorization is exact.
Differentiating the two logarithms along this path gives the same
derivative -bdj/E(z), so the factorization remains exact throughout.
This proves the principal-power identity even for negative bd.

Centering removes precisely the n=0 term. Thus the complete raw kernel
is

\[
 \boxed{
 K_j(b,d;t)=
 \left[1+bdj(1-e^{-t\lambda_j})+i(b-d)\sqrt j\right]^{-k}
 -\left[(1+ib\sqrt j)(1-id\sqrt j)\right]^{-k}.
 }
 \tag{2}
\]

Equivalently it is the absolutely convergent series
\(\sum_{n\ge1}\overline{a_n(b)}a_n(d)e^{-t\lambda_jn}\). Absolute
convergence also follows directly from Cauchy--Schwarz and the unit
norms of the uncentered vectors. In particular
\(K_j(d,b;t)=\overline{K_j(b,d;t)}\), as required for a Gram kernel.
When either amplitude is zero, the two terms in (2) coincide and the
centered kernel is zero exactly.

For any finite list (b\_l,u\_l), u\_l\textgreater=0, the time-vector
Gram matrix is

\[
 \left\langle e^{-u_l\lambda_jN}v_j(b_l),
                e^{-u_m\lambda_jN}v_j(b_m)\right\rangle
 =K_j(b_l,b_m;u_l+u_m).
 \tag{3}
\]

Every such matrix is positive semidefinite, directly from its
definition. No real-part operation has been applied to its off-diagonal
entries.

\subsection{2. Fixed global amplitudes and the discontinuous limiting
family}\label{weak-coupling-part-5-fixed-global-amplitudes-and-the-discontinuous-limiting-family}

For each fixed t\textgreater=0,

\[
 j(1-e^{-t\lambda_j})\longrightarrow ct.
 \tag{4}
\]

For b=d!=0, equation (2) consequently gives

\[
 K_j(b,b;t)\longrightarrow(1+cb^2t)^{-k}.
 \tag{5}
\]

The subtracted vacuum mass is
(1+b\textsuperscript{2j)}\{-k\}-\textgreater0. The limit in (5) is the
Laplace transform of Gamma(k,scale cb\^{}2), with raw total mass one.
Its physical scale contains the full b\^{}2 dependence.

For fixed b!=d, the imaginary part of the first denominator in (2) is
(b-d)sqrt(j), while its real part converges to 1+cbdt. Its inverse power
therefore tends to zero. When bd!=0 the second denominator has magnitude
of order j, so its inverse power also tends to zero. When bd=0, the
centered kernel is already zero. Thus

\[
 K_\infty(b,d;t)=
 \begin{cases}
 (1+cb^2t)^{-k},&b=d\ne0,\\
 0,&\text{otherwise}.
 \end{cases}
 \tag{6}
\]

For distinct nonzero b,d the leading off-diagonal phase is retained by
the stronger asymptotic

\[
 j^{k/2}K_j(b,d;t)\longrightarrow
 [i(b-d)]^{-k}
 =|b-d|^{-k}
       e^{-i\,\operatorname{sgn}(b-d)k\pi/2}.
 \tag{7}
\]

The vacuum subtraction is of smaller order j\^{}\{-k\} in this case.
Formula (7) does not apply when one amplitude is zero, because then its
two terms cancel identically. For bd\textless0 the vacuum term itself
has asymptotic phase \(e^{-i\,\operatorname{sgn}(b-d)k\pi}\); replacing
its denominator by an unsigned real magnitude would give an incorrect
finite-j kernel.

The uncentered limit additionally has F\_infinity(0,0;t)=1; all its
other entries agree with (6). At equal time it is the Kronecker kernel
on the set of real amplitudes:

\[
 F_\infty(b,d;0)=\mathbf1_{\{b=d\}}.
 \tag{8}
\]

Its amplitude-only cyclic Hilbert space is therefore
\(\ell^2(\mathbb R_{\mathrm{discrete}})\), with unit basis e\_b and
vacuum e\_0. The finite-j identity \(B_j(a)B_j(b)=B_j(a+b)\) gives the
exact limiting action \(U(a)e_b=e_{a+b}\). This is a unitary
representation of the additive group of real numbers, but is not
continuous in its usual topology:

\[
 \langle e_0,U(a)e_0\rangle=\mathbf1_{\{a=0\}},\qquad
 \|U(a)e_0-e_0\|=\sqrt2\quad(a\ne0).
\]

In particular this amplitude group has no self-adjoint generator P with
U(a)=exp(iaP). Such exponentials are strongly continuous by dominated
convergence in the spectral measure of every vector. The centered family
has the related discontinuity v(0)=0 and norm(v(b))=1 for every b!=0.

The time-cyclic completion has more vectors than this amplitude-only
subspace. An exact representation of all fixed-global-amplitude/time
limits is

\[
 \mathcal H_{\mathrm{global}}=
 \mathbb C\Omega\ \oplus\!
 \bigoplus_{b\in\mathbb R\setminus\{0\}}\mathcal H_\rho,
 \qquad H_b=cb^2Q,\quad Qf(q)=qf(q),
 \tag{9}
\]

with H Omega=0 and phi(b)=1 in its b summand. The operator Q and its
domain are proved in Section 5. Formula (9) gives (6) exactly. Each
fibre is cyclic under its semigroup: if f is orthogonal to all exp(-t
cb\^{}2 q), then the Laplace transform of the finite density f rho is
zero for t\textgreater=0. That transform is analytic for positive real
part by dominated differentiation on smaller half-planes. The identity
theorem makes it zero on that half-plane. Its values on s+iR,
s\textgreater0, are the Fourier transform of the integrable density
e\^{}\{-sq\}f rho, so Fourier uniqueness makes f=0.

Different nonzero amplitude sectors in (9) are orthogonal, but their
relation is explicit: all have the same q Hilbert space and generators
differing by the physical scalar cb\^{}2. Equivalently the map
f(omega)-\textgreater f(cb\^{}2q) is a unitary from L2(Gamma(k,scale
cb\^{}2)) to H\_rho, intertwining multiplication by omega with cb\^{}2Q.
In particular the b and -b fibres have the same energy operator despite
their orthogonality in this limiting family. Orthogonality does not
establish unrelatedness.

\subsection{3. Exact positive-g diagonal transport of the joint
kernels}\label{weak-coupling-part-5-exact-positive-g-diagonal-transport-of-the-joint-kernels}

For the actual full nonlinear Hamiltonian at fixed j use the globally
smooth physical multiplier

\[
 B_{g,j}(b;Z)=
 \exp\left(ib\sqrt j\,\chi(y)
 \frac{\sigma_j}{4}
 \left|(O^TG^{-1/2}(y/g))_{\nu_j}\right|^2\right)
 \tag{10}
\]

on the chart and one outside it. The same chi is used for all
amplitudes; it is one near zero and supported inside the chart. All
group products \(B_{g,j}(a)B_{g,j}(b)=B_{g,j}(a+b)\) are exact. The
original nonlinear Hamiltonian, all original plaquettes, the actual
positive unit vacuum psi\_gj, and the exact gauge projection are
unchanged. Define

\[
 m_{g,j}(b)=\langle\psi_{g,j},B_{g,j}(b)\psi_{g,j}\rangle,
 \quad v_{g,j}(b)=(B_{g,j}(b)-m_{g,j}(b))\psi_{g,j},
 \quad A_{g,j}=H_{g,j}-\mathcal E_{g,j}.
\]

At every fixed j, the companion's bounded-observable theorem transports
each of these vectors strongly to its exact comparison vector. Its
finite spectral projections and tightness argument then give

\[
 \langle v_{g,j}(b),e^{-tA_{g,j}}v_{g,j}(d)\rangle
       \longrightarrow K_j(b,d;t),\qquad t\ge0.
 \tag{11}
\]

Cross terms follow either directly from finite projections or by complex
polarization of strongly transported finite sums of these vectors. This
retains both their real and imaginary parts.

In particular, at stage j one may prescribe any finite list of amplitude
values b\_\{j,l\} and physical times t\_\{j,r\}, and choose
0\textless g\_j\textless1/j so that all their coefficient errors and all
their pairwise errors in (11) are less than 1/j. Each j is fixed before
applying its convergence theorem; the finite lists may contain
amplitudes of the form b\_0+alpha/sqrt(j). The minimum of finitely many
positive small-coupling thresholds is still positive. This proves the
finite-set diagonal statement.

There is also a useful justified strengthening to one diagonal covering
all bounded amplitude/time sets. Here are the details, so no uniformity
in volume is assumed. For each fixed j, write f\_g=mathcal B\_g psi\_g
and extend the transformed multiplier in (10) by any unit-modulus value
off Omega\_g. It equals exp(ib sqrt(j) q(x)) on a ball whose radius
tends to infinity as g decreases to zero. Therefore

\[
 \sup_{b\in\mathbb R}
 \|\mathcal B_g B_{g,j}(b)\psi_g-e^{ib\sqrt j q}\Phi_0\|
 \le \|f_g-\Phi_0\|
       +2\|\mathbf1_{\{|x|>R_g\}}\Phi_0\|
 \longrightarrow0,
 \tag{12}
\]

for a suitable R\_g-\textgreater infinity at that fixed j. Inner
products with f\_g and Phi\_0 imply uniform convergence of m\_gj(b) to
m\_j(b), and hence uniform strong convergence of centered vectors as
well.

For a compact amplitude interval {[}-M,M{]}, the comparison vectors form
a norm-compact set: continuity in b follows by dominated convergence
against the Gaussian measure. Finite oscillator spectral cutoffs tend
strongly to the identity, so their tail errors are uniform on this
compact set. To see this explicitly, approximate the compact set by a
finite epsilon-net, control every net vector by one cutoff, and use that
cutoff and its complement have norm at most one. Convergence in operator
norm of the transported finite spectral projections, together with (12),
gives the same uniform tail bound for the actual vectors when g is
small.

On a fixed finite spectral range, separate its finitely many eigenvalue
clusters. Their projections converge in operator norm, and their actual
excitation energies converge. Uniformly for 0\textless=t\textless=T, the
difference of the corresponding exponential scalars is bounded by T
times the energy difference, since all excitation energies are
nonnegative. Thus the finite-range kernel convergence is uniform for
\textbar b\textbar,\textbar d\textbar\textless=M and t in {[}0,T{]}. The
omitted same-projection tails contribute at most the product of the two
vector tail norms, by the contraction property of the semigroup.
Consequently (11) is uniform on each such compact parameter set.

At stage j choose a positive dyadic g\_j\textless1/j for which this
supremum error on \textbar b\textbar,\textbar d\textbar\textless=j,
0\textless=t\textless=j and the coefficient supremum error are both
\textless1/j. Such a dyadic exists by the fixed-j result just proved.
This yields a single actual positive-coupling diagonal for every fixed
real b,d,t and for all carrier-window sequences with fixed finite
b\_0,alpha,delta,t. It also gives the raw norms by t=0. No estimate
uniform in j was asserted for how small that dyadic must be, and no
prescribed renormalized running coupling is obtained.

\subsection{4. Rescaling around zero: continuous amplitude and zero
energy}\label{weak-coupling-part-5-rescaling-around-zero-continuous-amplitude-and-zero-energy}

Set b=alpha/sqrt(j), d=delta/sqrt(j), with alpha,delta fixed. The
original parameter in the multiplier becomes beta=alpha, exactly; q, its
original variance convention, the chart map, and a\_j are unchanged.
Equation (2) is now

\[
 \begin{split}
 K_j^{(0)}(\alpha,\delta;t)
 &=[1+\alpha\delta(1-e^{-t\lambda_j})+i(\alpha-\delta)]^{-k}\\
 &\quad-[(1+i\alpha)(1-i\delta)]^{-k}.
 \end{split}
\]

Since lambda\_j-\textgreater0, its limit is the continuous kernel

\[
 K^{(0)}(\alpha,\delta)=
 [1+i(\alpha-\delta)]^{-k}
 -[(1+i\alpha)(1-i\delta)]^{-k},
 \tag{13}
\]

independent of physical time. It is represented exactly on H\_rho by

\[
 v^{(0)}(\alpha)=e^{i\alpha q}-(1-i\alpha)^{-k},\qquad
 \Omega=1,
 \quad H^{(0)}=0.
 \tag{14}
\]

The raw norm in (14) is \(1-(1+\alpha^2)^{-k}\), strictly positive for
alpha!=0 and tending to zero continuously as alpha tends to zero. The
complete raw positive spectral measure is
\([1-(1+\alpha^2)^{-k}]\delta_0\). One may verify the energy collapse
directly in the exact Laguerre coordinate map: the coefficients of exp(i
alpha q) no longer depend on j, and exp(-t lambda\_j N) converges
strongly to the identity by dominated convergence of the squared
Laguerre coefficients. The first raw comparison energy moment is
lambda\_j k alpha\^{}2-\textgreater0 as an additional calculation, not
as an inference from weak convergence.

The unitary group W(alpha)f=exp(i alpha q)f is strongly continuous on
H\_rho, with self-adjoint generator Q on its maximal domain. It provides
a continuous amplitude representation, but every vector in (14) remains
a zero-energy vector. Retaining this entire origin family therefore
enlarges the zero eigenspace beyond the single vacuum.

At finite j the rate of loss of continuity in the unscaled parameter is
also explicit. Differentiation under the Gamma integral gives

\[
 \|\partial_b v_j(b)\|^2
 =j\left[k(k+1)-k^2(1+b^2j)^{-(k+1)}\right].
\]

Indeed the uncentered derivative is i sqrt(j) q phi\_j(b), of squared
norm j E(q\^{}2)=jk(k+1), and subtracting its vacuum component removes
\(|m_j'(b)|^2=jk^2(1+b^2j)^{-(k+1)}\). At b=0 this norm squared is jk.
Passing from b to alpha=b sqrt(j) removes precisely that factor j; it
does not alter the already-retained Hamiltonian spacing lambda\_j.

\subsection{5. Nonzero carrier windows: continuous modulation and Gamma
energy}\label{weak-coupling-part-5-nonzero-carrier-windows-continuous-modulation-and-gamma-energy}

Fix a nonzero real carrier b\_0. Define

\[
 b_j(\alpha)=b_0+\frac{\alpha}{\sqrt j},\qquad
 d_j(\delta)=b_0+\frac{\delta}{\sqrt j}.
\]

The exact multiplier identity is

\[
 B_j(b_j(\alpha))=B_j(b_0)e^{i\alpha q}.
 \tag{15}
\]

At positive g the same identity holds with q replaced by the exact
chart-cutoff radial function in (10), so these states share an actual
finite-regulator multiplier relation.

Now \((b_j-d_j)\sqrt j=\alpha-\delta\) exactly, \(b_jd_j\to b_0^2\), and
(4) holds. The first denominator in (2) therefore tends to
\(1+cb_0^2t+i(\alpha-\delta)\). The subtracted vacuum product tends to
zero because both beta parameters b\_0 sqrt(j)+alpha and b\_0
sqrt(j)+delta have magnitude tending to infinity. Thus

\[
 \boxed{
 K^{(b_0)}(\alpha,\delta;t)=
 [1+cb_0^2t+i(\alpha-\delta)]^{-k},\qquad b_0\ne0.
 }
 \tag{16}
\]

This is a continuous, nonconstant amplitude/time kernel. In particular
its equal-amplitude raw spectral measure is Gamma(k,scale cb\_0\^{}2),
independent of the finite offset alpha. Its off-diagonal entries retain
the full relative modulation phase. Equations (13) and (16) differ at
b\_0=0 precisely because the vacuum subtraction then survives.

Put h\_0=cb\_0\^{}2\textgreater0. On H\_rho define

\[
 \phi^{(b_0)}(\alpha)(q)=e^{i\alpha q},\qquad
 Qf(q)=qf(q),\quad
 \operatorname{Dom}Q=\left\{f:\int q^2|f(q)|^2d\rho(q)<\infty\right\},
 \qquad H^{(b_0)}=h_0Q.
 \tag{17}
\]

Then the Gamma integral proves exactly

\[
 \begin{split}
 \langle\phi^{(b_0)}(\alpha),
             e^{-tH^{(b_0)}}\phi^{(b_0)}(\delta)\rangle
 &=\int_0^\infty e^{-[h_0t+i(\alpha-\delta)]q}d\rho(q)\\
 &=[1+h_0t+i(\alpha-\delta)]^{-k}.
 \end{split}
 \tag{18}
\]

The real part of the integrating coefficient 1+h\_0t is positive, fixing
the same principal branch as in the finite-j formula.

Here are the operator and density details. Multiplication by q on the
maximal domain in (17) is densely defined and symmetric. For nonreal z,
the inverse multiplier (q-z)\^{}\{-1\} is bounded and maps H\_rho into
that domain, since q/(q-z) is bounded. It supplies both nonreal
resolvents; equivalently the adjoint's domain is exactly the same
maximal weighted L2 domain. Thus Q is self-adjoint. The unitary group
W(alpha)=exp(i alpha Q) is multiplication by exp(i alpha q). For every f
in H\_rho, dominated convergence with bound 4\textbar f\textbar\^{}2
proves its strong continuity. On Dom Q the difference quotient tends in
norm to iQf, using \textbar(exp(ihq)-1)/h\textbar\textless=q.
Conversely, existence of a norm derivative forces the difference
quotients to be bounded in L2; Fatou then gives qf in L2. Hence this is
exactly the generator domain, not a smaller unspecified core. Since
h\_0\textgreater0, H\^{}(b\_0) has the same domain.

The span of exp(i alpha q), alpha real, is dense in H\_rho. If f is
orthogonal to all these vectors, Cauchy--Schwarz makes the density
f(q)rho(q) integrable on the real line after extending it by zero to
negative q, and its Fourier transform vanishes at every real frequency.
For an explicit uniqueness argument, convolve that L1 density with a
Gaussian of any positive variance. The Gaussian Fourier integral and
Fubini express the convolution as the integral of its vanishing Fourier
transform times the Gaussian Fourier multiplier, so the convolution is
zero. Gaussian convolutions converge to the original L1 density as the
variance decreases to zero; this follows first for continuous compactly
supported functions by uniform continuity and Gaussian tails and then
for general L1 functions by density and contraction. Thus the density,
and therefore f, is zero. This proves the asserted dense span.

Consequently (18) determines the entire carrier-window cyclic Hilbert
space, rather than a proper unidentified subspace. It is the b\_0 fibre
already present in (9): the fixed carrier corresponds to alpha=0, and
the window supplies its vectors exp(i alpha q), which lie in that same
time-cyclic fibre. This is the exact map from the globally discontinuous
amplitude labels to their continuous local modulations.

The carrier Hamiltonian can also be checked by an exact differential
conjugation, before taking any limit. On the radial coordinate space,

\[
 Nf=-qf''-(k-q)f',\qquad T_\beta f=e^{i\beta q}f.
\]

For every \(f\in C_c^\infty((0,\infty))\), both f and T\_beta f belong
to the radial oscillator operator domain. Product differentiation with
all signs retained gives

\[
 T_\beta^*(\lambda_jN)T_\beta f
 =\lambda_jNf-2i\lambda_j\beta qf'
   +\lambda_j\beta^2qf-i\lambda_j\beta(k-q)f.
 \tag{18a}
\]

For beta=b\_0 sqrt(j), the coefficients satisfy lambda\_j
beta\textsuperscript{2-\textgreater cb\_0}2, lambda\_j
beta-\textgreater0, and lambda\_j-\textgreater0. Every function
multiplying them in (18a) is in L2(rho), since f is smooth with compact
support away from zero. Therefore

\[
 T_{b_0\sqrt j}^*(\lambda_jN)T_{b_0\sqrt j}f
       \longrightarrow cb_0^2Qf
       \quad\text{strongly in }\mathcal H_\rho.
 \tag{18b}
\]

The displayed test functions are an operator core for Q: truncate any f
in its maximal domain to {[}1/R,R{]}, which converges in the norm
\(\int(1+q^2)|f|^2d\rho\), and then approximate each truncated function
by smooth compactly supported functions in a slightly larger interval,
where the weighted density is bounded above and below by positive
constants. This proves a graph limit on a stated core with the exact
phase and physical-time scaling. Equation (18) supplies the complete
semigroup identification; no removal of the differential terms from the
original finite-j operator is asserted in (18a).

The spectral projections of H\^{}(b\_0) multiply by indicators of h\_0
q. Its constant vector has the raw probability law Gamma(k,scale h\_0),
and its kernel at zero is trivial because rho has no atom at q=0. Its
spectrum is {[}0,infinity): outside this interval the resolvent
multiplier is bounded, while normalized indicators of arbitrarily short
positive rho-measure intervals about q=E/h\_0 are approximate
eigenvectors for every E\textgreater=0. In particular the spectrum
reaches zero continuously.

A separate vacuum can be retained by adjoining C Omega with energy zero,
orthogonal to this carrier fibre. The kernel on that direct sum
represents the vacuum and all these carrier states; it has a unique zero
eigenvector within those retained sectors. Formula W(alpha) on the
carrier fibre does not by itself prescribe how the actual modulation
acts on Omega. The actual origin-window vectors in Section 4 give that
additional action if they too are retained, and then add the extra zero
sector described there. No arbitrary scalar action on Omega is inferred.

\subsection{6. The complete relation between different
windows}\label{weak-coupling-part-5-the-complete-relation-between-different-windows}

For two distinct fixed carriers b\_0!=d\_0, and finite offsets
alpha,delta, the term (b\_j-d\_j)sqrt(j) has magnitude tending to
infinity. The first term of (2) therefore vanishes. Each vacuum product
also vanishes unless both carriers are zero, which is excluded in this
case; if one carrier is zero its bounded coefficient multiplies the
other vanishing coefficient. Thus all cross kernels between different
carrier windows tend to zero, including cross kernels between a nonzero
carrier window and the origin window.

All uncentered window kernels together have the exact representation

\[
 \mathcal H_{\mathrm{windows}}
   =\bigoplus_{b_0\in\mathbb R}\mathcal H_\rho,
 \qquad
 \phi(b_0,\alpha)=e^{i\alpha q}\text{ in fibre }b_0,
 \qquad \Omega=\phi(0,0),
 \tag{19}
\]

with H acting as cb\_0\^{}2 Q in fibre b\_0 and as the zero operator on
the entire origin fibre. Its domain is the set of families f satisfying

\[
 \sum_{b_0\ne0}c^2b_0^4\int q^2|f_{b_0}(q)|^2d\rho(q)<\infty.
\]

This direct sum is self-adjoint by its fibre resolvents, whose norms are
bounded by 1/\textbar Im z\textbar{} uniformly in b\_0 for nonreal z.
Every vector has at most countably many nonzero components: the set of
components of norm at least 1/n is finite for every n.~Thus all sums and
norm arguments have their usual Hilbert direct-sum meaning.

The actual common modulation and carrier-shift identities have exact
counterparts on (19):

\[
 (W(\eta)f)_{b_0}(q)=e^{i\eta q}f_{b_0}(q),\qquad
 (U(a)f)_{b_0}(q)=f_{b_0-a}(q).
 \tag{20}
\]

W is strongly continuous by dominated convergence in the sum of the
fibre norms. U is unitary but nonregular, already on Omega. The two
groups commute. On their common dense window vectors they implement
alpha-\textgreater alpha+eta and b\_0-\textgreater b\_0+a respectively,
exactly matching the finite-regulator multiplier identities. Their
extension from that dense span is therefore determined. H has zero
kernel equal to the entire origin fibre, while each nonzero fibre has
trivial zero kernel.

The representation (9) embeds isometrically into (19) by retaining just
the constant vector in the origin fibre and every vector in every
nonzero fibre. That subspace is invariant under H and its semigroup. It
need not be invariant under W: W(eta)Omega is a nonconstant origin
vector for eta!=0. This identifies exactly which extra states appear
when the modulation is also applied to the vacuum. Neither construction
identifies a spatial local gauge algebra.

\subsection{7. Configuration density and retained
phases}\label{weak-coupling-part-5-configuration-density-and-retained-phases}

At every finite g,j and every real amplitude,
\textbar B\_gj(b)\textbar=1 pointwise. Thus for every bounded
multiplication observable F on the original configuration space,

\[
 \langle B\psi,F B\psi\rangle=\langle\psi,F\psi\rangle.
 \tag{21}
\]

This is an exact identity of uncentered configuration densities, before
any limit. It applies in particular to bounded physical multiplication
observables, with the actual vacuum and Haar measure retained.

For m=\textless psi,B psi\textgreater{} and v=(B-m)psi, expansion gives

\[
 \begin{split}
 \langle v,Fv\rangle-\langle\psi,F\psi\rangle
 &=-\overline m\langle\psi,FB\psi\rangle
   -m\langle\psi,FB^*\psi\rangle
   +|m|^2\langle\psi,F\psi\rangle,
 \end{split}
\]

and therefore the raw bound

\[
 \left|\langle v,Fv\rangle-\langle\psi,F\psi\rangle\right|
 \le\|F\|(2|m|+|m|^2).
 \tag{22}
\]

It also holds for complex bounded F. For the carrier-window amplitudes,
m\_j=(1-i(b\_0 sqrt(j)+alpha))\^{}\{-k\}-\textgreater0, and the
positive-g diagonal in Section 3 transports this coefficient.
Consequently (22) tends to zero for every uniformly bounded sequence of
configuration multipliers F\_j. The retained raw centered norm is
1-\textbar m\_j\textbar\^{}2-\textgreater1. If one additionally divides
by that explicitly retained norm, the corresponding error is bounded by
\(\|F\|(2|m|+2|m|^2)/(1-|m|^2)\), which also tends to zero.

These equal configuration statistics coexist with the nontrivial Gamma
energy law proved in (16)--(18). Their compatibility already holds at
finite regulators: multiplication by a phase preserves configuration
density, while electric derivatives differentiate that phase. For the
actual Hamiltonian's ground-state quadratic form and a smooth unit
multiplier B=exp(i h), the exact identity is

\[
 q_{H-\mathcal E}[B\psi]
   =\frac{2g^2}{a}\int\psi^2
      \sum_{e,\alpha}|X_{e,\alpha}B|^2dU
   =\frac{2g^2}{a}\int\psi^2
      \sum_{e,\alpha}|X_{e,\alpha}h|^2dU.
 \tag{23}
\]

It follows from the companion's exact ground-state transform, and
\textbar X exp(ih)\textbar{}\textsuperscript{2=\textbar Xh\textbar{}}2.
Subtracting m psi does not change this excitation form because
(H-mathcal E)psi=0. Equation (23) explains, with the original
coefficient 2g\^{}2/a, how a nonconstant smooth phase can alter energy
without altering configuration probabilities. It does not infer limits
of unspecified electric observables or unbounded moments from the
bounded-configuration identity.

\subsection{8. Proven scope}\label{weak-coupling-part-5-proven-scope}

The exact complex Laguerre sum yields a globally nonregular amplitude
family, a continuous origin rescaling whose energies collapse to zero,
and a continuous modulation window about every nonzero global radial
carrier with generator H=cb\_0\^{}2Q and a Gamma spectral law. The full
phase, vacuum subtraction, raw mass, physical time, and parameter
dictionary are preserved. The positive-g construction transports all
finite Gram matrices, and the fixed-j compactness argument supplies one
diagonal for the stated parameter families.

The carrier-window modulation is about a selected macroscopic radial
mode. Its exact representation, its inclusion into the global spectral
fibres, and its relation to the origin zero sector have been proved.
They do not specify spatial local fields, their complete
gauge-observable algebra, or the renormalized interacting
four-dimensional theory. Failure of the unscaled global amplitude to be
regular does not prove that every other bridge fails; the continuous
nonzero-carrier map above is an explicit further relation.

\clearpage
\part{The Gamma sector and bounded local holonomies}
\section{The Gamma sector and the exact action of bounded local
holonomies}\label{weak-coupling-part-6-the-gamma-sector-and-the-exact-action-of-bounded-local-holonomies}

8 September 2026. This note computes the configuration-observable map on
the radial spectral construction. It retains the original finite
nonlinear Hamiltonians and gives an explicit limit in which the Gamma
spectral sector survives while the unrescaled continuous gauge-invariant
cylinder algebra acts by scalars. No assertion about rescaled curvature
fields or all possible continuum constructions follows from that
particular map.

\subsection{1. Original states, measures, and an exact phase-density
identity}\label{weak-coupling-part-6-original-states-measures-and-an-exact-phase-density-identity}

Use the open-box SU(2) Hamiltonian, product Haar space, full vertex
gauge projection, and positive unit vacuum of the included fixed-box
proofs: \[
 H_j=\kappa_j\sum_eE_e+b_j\sum_p(2-\operatorname{tr}U_p),
 \quad \kappa_j=\frac{2g_j^2}{a_j},\quad
 b_j=\frac1{2g_j^2a_j},\quad L_j=j^2,\quad a_j=\frac1{100j}.
\] The excitation operator is \(A_j=H_j-\mathcal E_j\). Write
\(d\mu_j=|\psi_j|^2dU\) for the actual vacuum configuration measure. The
radial multiplier \(B_j\) of the Gamma note is globally smooth, gauge
invariant and has modulus one. Put \[
 m_j=\int B_j\,d\mu_j,\qquad
 v_j=(B_j-m_j)\psi_j .
\] For any bounded configuration multiplier \(F_j\), direct
multiplication on the original Haar space gives \[
 \begin{split}
 \langle v_j,F_jv_j\rangle
 &=\int F_j\bigl(1-\overline m_j B_j-m_j\overline B_j+|m_j|^2\bigr)d\mu_j,\\
 \left|\langle v_j,F_jv_j\rangle-\mu_j(F_j)\right|
 &\le\|F_j\|_\infty(2|m_j|+|m_j|^2).
 \tag{1}
 \end{split}
\] There is no division by \(\|v_j\|\) in (1). For the uncentered vector
\(B_j\psi_j\), its configuration measure is exactly \(\mu_j\) at every
regulator, not only asymptotically.

The fixed-box comparison mean is \(m_{0,j}=(1-i\sqrt j)^{-3/2}\). The
positive-coupling selection can impose \(|m_j-m_{0,j}|<1/j\) together
with the Gamma raw-measure estimates, because these are finitely many
convergent quantities at each fixed box. Thus \(|m_j|\to0\), so (1)
identifies the limiting bounded configuration expectations of the
centered state with those of the vacuum. The statement does not identify
their electric or general noncommuting observable expectations.

\subsection{2. Positive coupling with a retained local flatness
estimate}\label{weak-coupling-part-6-positive-coupling-with-a-retained-local-flatness-estimate}

For each integer \(j\ge2\), choose a dyadic \(g_j=2^{-n_j}>0\) small
enough to satisfy all the Gamma note's raw spectral and centering
estimates and \(g_j<j^{-5}\). Let \(n_j\) be the least positive integer
satisfying the finite list. The fixed-box limits prove existence.
Consequently \[
 g_j^2j^8<j^{-2},\quad
 \kappa_j=200jg_j^2,\quad b_j=50j/g_j^2,\quad
 \xi_j=\frac1{4g_j^4}.
 \tag{2}
\] Every finite Hamiltonian retains its full electric term, magnetic
term, and scalar Wilson contribution. Along this selected diagonal the
raw spectral measure of \(v_j\) still tends to the Gamma law with shape
\(3/2\) and scale \(c=100\sqrt2\pi\).

Here are the exact spatial maps needed for the local observables. Fix a
coarse positive integer \(k\ge2\). When \(k\) divides \(j\), replace
each coarse edge by its ordered product of \(j/k\) collinear fine links,
with coarse vertex \(n\) placed at fine vertex \((j/k)n\). Physical
positions agree since \(a_j(j/k)n=a_kn\). Denote this continuous map by
\(p_{k,j}\). The small coarse box is contained in the fine box. It
intertwines the vertex gauge actions, with all intermediate endpoint
factors cancelling in their actual order. No reordering of noncommuting
links is used. For a common nested sequence one may restrict to
\(j=n!\); then every fixed \(k\) divides all sufficiently late indices.
The Gamma formulas and their constants are unchanged on this cofinal
subsequence.

Section 13 of the included \emph{Spatial continuum} proof establishes \[
 \int \|I-U_C\|_{\rm HS}^2d\mu_j
 \le 12m^2g_j^2\sqrt{N_jM_j}
 \tag{3}
\] for an \(m\)-by-\(m\) fine square. Its complete proof retains the
ordered non-Abelian filling, the actual vacuum variational bound and all
boundary counts. In particular \(N_j=6j^2(2j^2+1)^2\) is the number of
fine edges, \(M_j=12j^4(2j^2+1)\) the number of fine faces, and \[
 \sqrt{N_jM_j}\le18\sqrt5\,j^6.
\] For clarity the algebraic filling step is not a commuting
approximation: if a disk boundary is \(AqB\) and a newly attached face
replaces its shared segment \(q\) by \(r\), the changed holonomy times
the old inverse is \(A(rq^{-1})A^{-1}\). Iteration gives one conjugated
oriented plaquette per face. Unitary invariance and Cauchy--Schwarz give
\(\|I-U_C\|_{\rm HS}^2\le2m^2\sum_{\text{faces}}(2-\operatorname{tr}U_p)\).
The vacuum bound for the full face sum then gives (3).

Taking \(m=j/k\) gives for every coarse elementary face \[
 \int\|I-U_p\|_{\rm HS}^2d(p_{k,j})_*\mu_j
 \le\frac{216\sqrt5}{k^2}g_j^2j^8
 \le\frac{216\sqrt5}{k^2j^2}.
 \tag{4}
\] The finite number of coarse faces permits summation of (4). The
following identifies its full limiting measure, not just individual
plaquette means. On the contractible coarse box, flat link assignments
are exactly \[
 U_e=q_{s(e)}^{-1}q_{t(e)},\qquad q_o=I,
 \tag{5}
\] where \(o\) is the coarse root. To prove (5), construct \(q_v\) as
the ordered canonical path from \(o\) to \(v\). Each elementary flatness
identity moves one adjacent coordinate step past another, so all paths
have the same holonomy; hence \(q_sU_e=q_t\). Conversely (5) makes each
face telescope to identity. The inverse map is the canonical path
product. The flat configurations therefore form one compact vertex-gauge
orbit. Independent Haar \(q_v\), \(v\ne o\), give its invariant
probability \(\sigma_k\). It is unique: averaging a continuous function
over the compact gauge group is constant on this transitive orbit, so
integration against any invariant probability gives the same result.

Compactness supplies subsequential limits of the pushed vacuum measures.
Equation (4) puts each limit on the flat orbit; gauge invariance passes
to the limit. Uniqueness of \(\sigma_k\) identifies every subsequential
limit, hence the full sequence. Thus for every continuous
gauge-invariant function \(F\) on the fixed coarse link space, with
\(F_j=F\circ p_{k,j}\), \[
 \|(F_j-F(I))\psi_j\|^2
 =\int|F-F(I)|^2d(p_{k,j})_*\mu_j\longrightarrow0 .
 \tag{6}
\] Gauge invariance is essential in (6): a general noninvariant function
need not be constant on the flat orbit.

\subsection{3. Exact semigroup domination on the full nonlinear
model}\label{weak-coupling-part-6-exact-semigroup-domination-on-the-full-nonlinear-model}

The full finite-volume heat semigroup of \(H_j\) is positivity
preserving. One proof uses its original product-group electric heat
kernels and the positive multiplication operator \(\exp(-tV_j)\), where
\(V_j=b_j\sum_p(2-\operatorname{tr}U_p)\) is bounded and real. The
Trotter product converges strongly and preserves the positive cone. For
complex input, each positive integral kernel satisfies \(|Kf|\le K|f|\);
the same domination passes to the limit in \(L^2\) and almost everywhere
along a subsequence. Multiplication by the scalar
\(\exp(t\mathcal E_j)\) preserves it for \(\exp(-tA_j)\). Moreover
\(\exp(-tA_j)\psi_j=\psi_j\). Therefore for any bounded \(h_j\) and
\(t\ge0\), \[
 \big|\exp(-tA_j)(h_j\psi_j)\big|
 \le\exp(-tA_j)(|h_j|\psi_j)
 \le\|h_j\|_\infty\psi_j .
 \tag{7}
\] The physical subspace is invariant, so (7) applies to all the radial
physical states without replacing the Hamiltonian or gauge projection.
Since \(|B_j-m_j|\le2\), (7) yields \[
 |\exp(-tA_j)v_j|\le2\psi_j .
 \tag{8}
\] For any fixed finite coefficients \(a_\ell\) and nonnegative times
\(t_\ell\), let \(w_j=\sum_\ell a_\ell\exp(-t_\ell A_j)v_j\). Equations
(6)--(8) prove the explicit estimate \[
 \|(F_j-F(I))w_j\|
 \le2\left(\sum_\ell|a_\ell|\right)
       \|(F_j-F(I))\psi_j\|\longrightarrow0 .
 \tag{9}
\] The same result holds when finitely many bounded-phase radial
families are included, since each centered multiplier still has supremum
at most two. No estimate on an unbounded electric operator has been
inserted.

\subsection{4. The actual limiting observable
representation}\label{weak-coupling-part-6-the-actual-limiting-observable-representation}

The Gamma note proves convergence of every finite time-orbit Gram matrix
to the one on \(L^2(\gamma_c)\), and that these orbit vectors are dense
there. Equation (9) therefore defines the limit of each fixed continuous
gauge-invariant configuration cylinder on this dense space as \[
 \pi_{\rm cyl}(F)=F(I)\,I_{\mathcal H_\gamma}.
 \tag{10}
\] This definition respects sums, products, conjugation and the unit,
since evaluation at the identity configuration has those properties.
Also \(\|\pi_{\rm cyl}(F)\|\le\|F\|_\infty\), so it extends by
continuity to the uniform closure of this cylinder algebra. The same
scalar acts on the separately retained vacuum. Formula (9) proves the
operator map on actual approximating vectors, not merely an equality of
expectations.

The energy operator on the Gamma sector is nonetheless \(M_\omega\),
with strictly positive measure on every \((0,\varepsilon)\). These two
facts coexist: on this particular selected diagonal, the unchanged
bounded local holonomies cannot generate that nontrivial spectral sector
from the vacuum. Their representation (10) is explicit. It neither
proves that a rescaled curvature or electric algebra has the same limit
nor eliminates another running-coupling path.

The next mathematical map must retain a nonconstant observable beyond
(10). Candidates include the actual plaquette curvature divided by its
specified area and coupling, or gauge-invariant electric fluctuations.
They are unbounded along the regulator sequence, so the bounded factor
estimate (9) does not determine their limit. Their norms, moments,
domains, products and dynamics must be calculated with those factors
present; assigning them the scalar value (10) would be unjustified.

\subsection{Proof inputs and
literature}\label{weak-coupling-part-6-proof-inputs-and-literature}

The full operator, ground-state and fixed-box convergence proofs are
included as \emph{The original finite-box SU(2) Hamiltonian at small
positive coupling} and \emph{Actual weak-coupling vacuum, observable and
spectral maps}. The nonlinear vacuum estimate used in (3) is proved in
Sections 12--13 of the included \emph{Spatial continuum} manuscript.
These are integral mathematical inputs, not assumptions replacing the
estimates.

The distinction between raw loop observables and their continuum scaling
is discussed in S. Chatterjee, \emph{Yang--Mills for probabilists},
Sections 3 and 6--7
(\href{https://arxiv.org/abs/1803.01950}{arXiv:1803.01950}). The
nontrivial local-field target is specified by A. Jaffe and E. Witten,
\emph{Quantum Yang--Mills Theory}, Section 2
(\href{https://www.claymath.org/wp-content/uploads/2022/06/yangmills.pdf}{Clay
problem text}). Neither source is cited as proving the new operator
limit (9)--(10); that calculation is given above.

\clearpage
\part{Full operator-word limit of the radial carrier algebra}
\section{Full operator-word limit of the radial carrier
algebra}\label{weak-coupling-part-7-full-operator-word-limit-of-the-radial-carrier-algebra}

8 September 2026. This proof extends the complete amplitude/time kernels
to products of the actual bounded operators, including real-time
evolution. It identifies the resulting operator algebra and its centre.
The original spatial lattice, gauge projection, physical time and
coupling coefficients are retained. This is the algebra of the specified
global radial observables; it is not an identification with the
continuum local Yang--Mills algebra.

\subsection{1. Actual finite operators and their previously proved
kernels}\label{weak-coupling-part-7-actual-finite-operators-and-their-previously-proved-kernels}

Keep \(j\ge2\), \(L_j=j^2\), \(a_j=1/(100j)\), and \[
 H_j=\frac{2g_j^2}{a_j}\sum_eE_e+
 \frac1{2g_j^2a_j}\sum_p(2-\operatorname{tr}U_p),\qquad
 A_j=H_j-\mathcal E_j\ge0 .
 \tag{1}
\] Every edge and contained face of the original open box occurs in (1).
The space is the full vertex-gauge-invariant product-Haar space. The
actual positive unit vacuum is \(\psi_j\), and \(A_j\psi_j=0\). Its
uniqueness is proved in the included finite-box companion. In particular
the entire kernel at each regulator is \(\mathbb C\psi_j\).

The globally smooth physical chart function and multipliers are \[
 Q_j=g_j^{-2}\chi(y)\frac{\sigma_j}{4}
       |(O^{\mathsf T}G^{-1/2}y)_{\nu_j}|^2,\quad
 B_j(\beta)=e^{i\beta Q_j},\quad
 \sigma_j=\sqrt8\sin\frac{\pi}{4j^2+2}.
 \tag{2}
\] The cutoff extends the phase by zero outside the logarithm chart,
hence the exponential by one. The same real cutoff is used for every
amplitude. The original quotient kinetic tensor is \(G=TT^*\); \(O\)
diagonalizes \(G^{1/2}C^*CG^{1/2}\). Every \(B_j(\beta)\) is unitary and
exactly physical. Its product law is
\(B_j(\beta)B_j(\gamma)=B_j(\beta+\gamma)\).

Put \(k=3/2\), \(c=100\sqrt2\pi\) and \[
 d\rho(q)=\frac{q^{k-1}e^{-q}}{\Gamma(k)}\,dq,\qquad q>0.
 \tag{3}
\] The included amplitude/time proof constructs one positive dyadic
coupling sequence on which all compact-parameter finite-box kernel tests
converge. For the uncentered vectors \[
 \phi_j(b,\alpha)=B_j(b\sqrt j+\alpha)\psi_j,
 \qquad b,\alpha\in\mathbb R,
 \tag{4}
\] its exact limit is \[
 \lim_j\langle\phi_j(b,\alpha),e^{-tA_j}\phi_j(d,\delta)\rangle
 =\mathbf1_{\{b=d\}}
  [1+cb^2t+i(\alpha-\delta)]^{-k},\qquad t\ge0 .
 \tag{5}
\] This includes \(b=d=0\); it does not discard the origin modulations.
The original finite comparison denominator is \[
 1+(b\sqrt j+\alpha)(d\sqrt j+\delta)
       (1-e^{-t\lambda_j})
       +i\{(b-d)\sqrt j+\alpha-\delta\},
 \quad
 \lambda_j=400\sqrt2j\sin\frac{\pi}{4j^2+2}.
 \tag{6}
\] Its inverse \(k\)-th power, with the branch proved in that companion,
gives (5), because \(j\lambda_j\to c\). The actual positive-\(g\)
estimates, including compact-uniform spectral cutoffs, are part of that
proof; (6) is not substituted for an actual finite-coupling kernel.

\subsection{2. Common dense vectors and the exact limiting
operators}\label{weak-coupling-part-7-common-dense-vectors-and-the-exact-limiting-operators}

Let \[
 \mathcal H=\bigoplus_{b\in\mathbb R}L^2(\rho),\qquad
 \phi(b,\alpha)_d(q)=\mathbf1_{\{d=b\}}e^{i\alpha q}.
 \tag{7}
\] Its vacuum vector is \(\Omega=\phi(0,0)\). The finite span
\(\mathcal D\) of these vectors is dense: finite label support is dense
in a Hilbert direct sum, and exponentials are total in \(L^2(\rho)\).
For the latter assertion, orthogonality gives a finite density with
vanishing Fourier transform. Convolution with Gaussians, followed by
their \(L^1\) approximation to the identity, forces that density to
vanish. This establishes density without assuming a spatial
interpretation for either label.

Define \(J_j:\mathcal D\to\mathcal H_j\) by
\(J_j\phi(b,\alpha)=\phi_j(b,\alpha)\), extended linearly. This
definition is unambiguous. For each fixed \(b\), finitely many distinct
exponentials are linearly independent on \(q>0\): a vanishing finite
exponential sum is analytic and all its derivatives at any point vanish;
the resulting Vandermonde matrix has nonzero determinant. Different
\(b\) fibres are independent by the direct sum. Equation (5) at zero
time proves \[
 \langle J_jv,J_jw\rangle\longrightarrow\langle v,w\rangle
 \qquad(v,w\in\mathcal D).
 \tag{8}
\] No unit-norm adjustment is made to \(J_jv\).

The limit Hamiltonian is maximal multiplication \[
 (Hf)_b(q)=cb^2qf_b(q),\quad
 \operatorname{Dom}H=
 \left\{f:\sum_b c^2b^4\int q^2|f_b|^2d\rho<\infty\right\}.
 \tag{9}
\] Its nonreal resolvents are multiplication by \((cb^2q-z)^{-1}\),
bounded uniformly by \(1/|\operatorname{Im}z|\). The adjoint has the
same maximal domain, so \(H\) is self-adjoint and nonnegative. Every
vector has at most countably many nonzero fibres, since for every
positive integer \(n\) only finitely many have squared norm at least
\(1/n\). Thus sums and dominated-convergence arguments below apply to
each vector even though the full label set is uncountable.

Define the bounded operators \[
 \begin{split}
 (U(a)f)_b(q)&=f_{b-a}(q),&
 U_j(a)&=B_j(a\sqrt j),\\
 (V(\eta)f)_b(q)&=e^{i\eta q}f_b(q),&
 V_j(\eta)&=B_j(\eta),\\
 S(t)&=e^{-tH},&
 S_j(t)&=e^{-tA_j}\quad(t\ge0).
 \end{split}
 \tag{10}
\] The two multiplier correspondences hold exactly on the representing
vectors: \[
 U_j(a)J_jv=J_jU(a)v,\qquad
 V_j(\eta)J_jv=J_jV(\eta)v,\qquad v\in\mathcal D.
 \tag{11}
\] All \(U,V\) and their actual versions are unitary; all \(S,S_j\) are
contractions. Also \(U\) and \(V\) commute, as do their finite versions.
This does not assert commutation with \(S\).

\subsection{3. Strong convergence across the different Hilbert
spaces}\label{weak-coupling-part-7-strong-convergence-across-the-different-hilbert-spaces}

Here is the precise notion used for products. Say \(\xi_j\to\xi\)
strongly in the \(J\) sense if for every \(\varepsilon>0\) there is
\(v\in\mathcal D\) with \[
 \|\xi-v\|<\varepsilon,\qquad
 \limsup_j\|\xi_j-J_jv\|<\varepsilon .
 \tag{12}
\] Allowing a fixed larger multiple of \(\varepsilon\) gives the same
definition by replacing \(\varepsilon\). This convergence implies
convergence of norms and of pairwise scalar products, by (8) and
approximation with the same finite vectors.

A useful sufficient test is \[
 \|\xi_j\|\to\|\xi\|,\qquad
 \langle J_jv,\xi_j\rangle\to\langle v,\xi\rangle
 \quad\hbox{for every }v\in\mathcal D.
 \tag{13}
\] Indeed expansion of the squared difference gives
\(\|\xi_j-J_jv\|^2\to\|\xi-v\|^2\); density then proves (12).

For \(w\in\mathcal D\), (5) and finite linearity give \[
 \langle J_jv,S_j(t)J_jw\rangle\to\langle v,S(t)w\rangle .
 \tag{14}
\] The norm is supplied by the same kernel at twice the time: \[
 \|S_j(t)J_jw\|^2
 =\langle J_jw,S_j(2t)J_jw\rangle
 \longrightarrow
 \langle w,S(2t)w\rangle=\|S(t)w\|^2 .
 \tag{15}
\] Equations (13)--(15) prove strong \(J\)-convergence of \(S_j(t)J_jw\)
to \(S(t)w\). This output-norm step is essential; scalar two-point
convergence alone would not justify products.

More generally suppose \(\|T_j\|,\|T\|\le C\) and \(T_jJ_jv\to Tv\)
strongly in this sense for all \(v\in\mathcal D\). If \(\xi_j\to\xi\),
then \(T_j\xi_j\to T\xi\). To prove it, choose \(v\) as in (12). The
difference from \(T_jJ_jv\) has limsup norm at most \(C\varepsilon\);
the latter sequence converges strongly to \(Tv\), and
\(\|T\xi-Tv\|\le C\varepsilon\). Approximate \(Tv\) by a member of
\(\mathcal D\) in the definition, then let \(\varepsilon\downarrow0\).
This is an actual bounded-operator continuity argument, not an
assumption that arbitrary regulator products converge.

Apply this lemma to (11), (14)--(15). Induction proves for any fixed
finite word \(T_j^{(m)}\cdots T_j^{(1)}\) in the operators (10) and any
\(w\in\mathcal D\), \[
 T_j^{(m)}\cdots T_j^{(1)}J_jw
 \longrightarrow
 T^{(m)}\cdots T^{(1)}w
 \quad\hbox{strongly in the }J\hbox{ sense}.
 \tag{16}
\] All coefficients, amplitudes and nonnegative times in the word are
retained. Pairwise inner products of any two such words therefore
converge as well, including words applied to \(\Omega\).

\subsection{4. The same sequence also retains real
time}\label{weak-coupling-part-7-the-same-sequence-also-retains-real-time}

For \(w\in\mathcal D\), let \(\mu_j^w\) be its actual raw spectral
measure for \(A_j\) on \(J_jw\), and let \(\mu^w\) be the spectral
measure for \(H\) on \(w\). Their masses converge by (8), and their
Laplace transforms converge by (14).

These facts imply tightness; it is not inferred from a moment bound. For
\(t,R>0\), \[
 \mu_j^w([R,\infty))
 \le
 \frac{\mu_j^w([0,\infty))-\int e^{-t\omega}d\mu_j^w}
      {1-e^{-tR}} .
 \tag{17}
\] The limiting numerator tends to zero as \(t\downarrow0\), since
\(S(t)w\to w\) by spectral dominated convergence. Choose such a small
fixed \(t\), then \(R\) large so the denominator is at least \(1/2\).
This controls the tail uniformly for all sufficiently large \(j\). The
finitely many omitted measures have their own vanishing tails, proving
tightness of the sequence.

Every weak subsequential limit has the stated Laplace transform.
Uniqueness can be seen by the change \(x=e^{-\omega}\): equality at
integer times gives equality of all polynomial moments on \([0,1]\),
after assigning zero mass at the added point zero. Polynomial density
determines the measure. Thus \(\mu_j^w\Rightarrow\mu^w\) on
\([0,\infty)\). Complex polarization supplies the same conclusion for
the mixed measures of any \(v,w\in\mathcal D\).

In particular the bounded continuous tests \(e^{-is\omega}\) give
convergence of all real-time matrix elements, for every fixed
\(s\in\mathbb R\). The actual and limiting real-time groups are unitary,
so their output norms also converge. Criterion (13) proves \[
 e^{-isA_j}J_jw\longrightarrow e^{-isH}w .
 \tag{18}
\] The bounded-product argument now proves (16) for finite words also
containing real-time evolution. Neither (17) nor (18) asserts
convergence of unbounded energy moments or uniformity in time intervals
whose lengths diverge with the regulator.

\subsection{5. Domains, shifts, and the retained energy
relation}\label{weak-coupling-part-7-domains-shifts-and-the-retained-energy-relation}

Let \(P\) be multiplication by the label \(b\), with domain
\(\sum_b b^2\|f_b\|^2<\infty\), and let \(Q\) be multiplication by
\(q\), with domain \(\sum_b\int q^2|f_b|^2d\rho<\infty\). Their spectral
projections commute. Formula (9) is the maximal joint
functional-calculus operator \(cQP^2\); it is not restricted by first
demanding both \(P^2f\) and \(Qf\) separately. For example the entire
origin fibre belongs to \(\operatorname{Dom}H\) even if its vector is
not in \(\operatorname{Dom}Q\).

The shifts preserve \(\operatorname{Dom}P\) and \(\operatorname{Dom}Q\),
and direct index substitution proves \[
 U(a)^*PU(a)=P+aI,\qquad U(a)^*QU(a)=Q .
 \tag{19}
\] For the energy they give the exact identity \[
 U(a)^*HU(a)=cQ(P+aI)^2
 \tag{20}
\] on the shifted maximal domain \[
 \left\{f:\sum_b c^2(b+a)^4\int q^2|f_b|^2d\rho<\infty\right\}.
\] A shift need not preserve \(\operatorname{Dom}H\): it can take an
arbitrary origin-fibre vector to a nonzero fibre, where a \(q^2\) moment
is required. On finite-label smooth compactly supported functions of
\(q>0\), all products and polynomial differences implicit in (19)--(20)
are well defined with their displayed signs.

The variable \(Q\) is central for the bounded generators: its bounded
spectral functions commute with \(U,V\) and both time evolutions. This
centrality is a statement about the calculated radial algebra, not about
the centre of all Yang--Mills observables.

\subsection{6. The complete centre of the generated von Neumann
algebra}\label{weak-coupling-part-7-the-complete-centre-of-the-generated-von-neumann-algebra}

Let \(\mathcal M\) be the von Neumann algebra generated by all the
bounded operators (10). Since \(q>0\) almost everywhere, spectral
dominated convergence in each countable-support vector proves \[
 E_0=\operatorname*{s-lim}_{t\to\infty}S(t),
 \tag{21}
\] where \(E_0\) is the projection onto the whole origin fibre.
Consequently \(E_b=U(b)E_0U(-b)\) belongs to \(\mathcal M\) for every
\(b\). The operator \[
 E_bU(b-d)E_d
 \tag{22}
\] identifies fibre \(d\) with fibre \(b\) by the identity on
\(L^2(\rho)\) and vanishes elsewhere. Thus all label matrix units occur
with their full infinite-dimensional \(q\) fibre.

Suppose \(T\) commutes with \(\mathcal M\). Commutation with every
\(E_b\) makes it fibre diagonal, \(T=\bigoplus_b T_b\). Commutation with
(22) forces \(T_b=T_d\) under the displayed fibre identity. Write this
common operator as \(T_*\). Commutation with every \(V(\eta)\) forces
\(T_*\) to commute with all bounded spectral functions of \(Q\). For
completeness the latter functions are generated by the exponentials: the
spectral theorem, or uniqueness of their finite Fourier measures applied
to matrix elements, shows that an operator commuting with every
\(e^{i\eta Q}\) commutes with its spectral projections. On the scalar
space \(L^2(\rho)\), an operator commuting with all multiplication
projections is itself a multiplier. Indeed \(h=T_*1\) satisfies
\(T_*\mathbf1_E=h\mathbf1_E\). The norm bound on every indicator gives
\(|h|\le\|T_*\|\) almost everywhere, and density of simple functions
proves \(T_*=M_h\).

Conversely any common multiplier \(I\otimes M_h\),
\(h\in L^\infty(\rho)\), commutes with every generator. It also belongs
to \(\mathcal M\) as a bounded spectral function of \(Q\). We have
proved \[
 \mathcal M'=\{I\otimes M_h:h\in L^\infty(\rho)\}
 \subset\mathcal M,\qquad
 Z(\mathcal M)=\mathcal M'
 \cong L^\infty(\rho).
 \tag{23}
\] Here tensor notation is the exact unitary direct-sum map proved in
the polymer-representation companion; it does not require a countability
assumption on all possible labels.

Thus the Gamma coordinate survives as the complete centre of this radial
algebra. The proof supplies a relation to known nonregular polymer
kinematics rather than declaring that kinematics new. The novel claim of
scope here is the particular operator-word map from the specified
nonlinear regulators, with its retained spectral measure and physical
scale.

\subsection{7. Infinite-time closure does not commute with the regulator
limit}\label{weak-coupling-part-7-infinite-time-closure-does-not-commute-with-the-regulator-limit}

Equation (16) concerns fixed finite words. It does not imply convergence
of every operator in the von Neumann closure from a naively
corresponding regulator operator. There is an exact example already in
these sources.

Fix \(\alpha\ne0\) and form the actual centered origin vector \[
 z_j=(B_j(\alpha)-\langle\psi_j,B_j(\alpha)\psi_j\rangle)\psi_j .
 \tag{24}
\] It tends strongly in the \(J\) sense to \[
 z=\phi(0,\alpha)-(1-i\alpha)^{-k}\Omega,\qquad
 \|z\|^2=1-(1+\alpha^2)^{-k}>0 .
\] At each regulator, the infinite-time limit of \(S_j(t)\) is the
rank-one vacuum projection. It kills \(z_j\) exactly. In the limiting
representation \(z\) lies in the origin fibre, so \(S(t)z=z\) for every
\(t\), and \(E_0z=z\). Consequently \[
 \lim_j\lim_{t\to\infty}\|S_j(t)z_j\|^2=0,\qquad
 \lim_{t\to\infty}\lim_j\|S_j(t)z_j\|^2
       =1-(1+\alpha^2)^{-k}.
 \tag{25}
\] Both orders are explicitly calculated; they cannot be interchanged.
This is also why the centre argument in (21)--(23), which concerns the
completed limiting algebra, must not be advertised as convergence of the
individual finite-volume vacuum projections.

\subsection{8. Relation to local observables and the research
target}\label{weak-coupling-part-7-relation-to-local-observables-and-the-research-target}

The bounded spatial cylinder conclusion extends to all the operator
words already proved, including real-time conjugates. Here is the
additional argument. Use the common positive-coupling refinement
containing both the kernel tests and the vacuum flatness tests, and the
divisible coarse/fine sequence in the configuration companion. For any
fixed bounded continuous gauge-invariant coarse cylinder \(F\), write
\(F_j\) for its actual pullback. That companion proves \[
 \|(F_j-F(I))\psi_j\|\longrightarrow0,\qquad
 \|F_j\|\le\|F\|_\infty.
 \tag{26}
\] For
\(v=\sum_{\ell=1}^n z_\ell\phi(b_\ell,\alpha_\ell)\in\mathcal D\), the
actual unit-modulus phases give the pointwise bound \[
 |J_jv|\le\left(\sum_{\ell=1}^n|z_\ell|\right)\psi_j.
\] Thus, without an amplitude or time bound on these fixed parameters,
\[
 \|(F_j-F(I))J_jv\|
 \le\left(\sum_{\ell=1}^n|z_\ell|\right)
          \|(F_j-F(I))\psi_j\|\longrightarrow0.
 \tag{27}
\] The bounded-operator lemma of Section 3 applies with limiting
operator \(F(I)I\). It follows that all fixed finite mixed words
containing these cylinders, carrier shifts, modulations, and both time
evolutions converge strongly in the \(J\) sense. In particular, for
every fixed real \(s\) and \(v\in\mathcal D\), \[
 e^{isA_j}F_j e^{-isA_j}J_jv
       \longrightarrow F(I)v.
 \tag{28}
\] This proves the stated temporal action on the existing spatial
cylinders; it is not necessary to assume it as an additional limit.
Finite collections of coarse cylinders can be placed on a common
divisible grid by the least common multiple of their subdivision
indices. The factorial cofinal sequence in the companion then treats
every fixed such cylinder. No assertion for a continuously moving
regulator-dependent spatial support is hidden in this argument.

The configuration companion proves that on a compatible smaller
positive-coupling diagonal, unchanged bounded continuous local
gauge-invariant holonomies act by \(F(I)I\) on the radial time orbits.
The spatial companion separately proves that the actual phase
double-commutator observable \(jR_g\) tends to \(cQ\) on weighted
carrier time vectors, after explicitly adding weighted tests to its
dyadic selection. This realizes a nonconstant electric observable in the
centre (23), with full tail estimates; it does not make that global
selected-mode observable spatially local.

There remains no proved map from the full spatial local interacting
Yang--Mills algebra onto this radial algebra preserving the required
vacuum, field relations and spacetime dynamics. The positive spectrum
down to zero in the nonzero carrier fibres therefore does not establish
a counterexample for that target theory. What is established is the
complete bounded-word representation, its real-time extension, its exact
energy domains, its centre, and its additional zero-energy fibre. Those
results specify concrete structures any proposed extension must
calculate rather than hide in a change of coordinates.

\clearpage
\part{The radial carrier representation and polymer kinematics}
\section{Exact relation of the radial carrier representation to polymer
kinematics}\label{weak-coupling-part-8-exact-relation-of-the-radial-carrier-representation-to-polymer-kinematics}

8 September 2026. This is a representation-theoretic identification of
the already constructed global radial limit. It is not a claim that
Yang--Mills has become quantum gravity, nor a replacement of the
finite-regulator construction by a chosen representation.

\subsection{1. Source and conventions}\label{weak-coupling-part-8-source-and-conventions}

Ashtekar, Fairhurst and Willis construct a nonregular Weyl
representation with an orthonormal basis indexed by the real numbers in
Section III.B of \emph{Quantum gravity, shadow states, and quantum
mechanics}, Classical and Quantum Gravity 20 (2003), 1031--1062
(\href{https://arxiv.org/abs/gr-qc/0207106v3}{arXiv:gr-qc/0207106v3}).
Their equations labelled innerprod, uvops and xop give a Hilbert space
with basis \(|x\rangle\), \(x\in\mathbb R\), and actions \[
 \widehat U(\ell)|x\rangle=e^{i\ell x}|x\rangle,\qquad
 \widehat V(\mu)|x\rangle=|x-\mu\rangle,\qquad
 \widehat x|x\rangle=x|x\rangle.
 \tag{1}
\] In particular the sign of the shift in their convention is negative.
The label, its physical interpretation, and the Hamiltonian must not be
inferred from equality of these kinematic formulas alone.

\subsection{2. The unitary map with all labels and signs
retained}\label{weak-coupling-part-8-the-unitary-map-with-all-labels-and-signs-retained}

The full radial window space proved in the amplitude note is \[
 \mathcal H_{\rm win}
 =\bigoplus_{b\in\mathbb R}L^2((0,\infty),\rho),\qquad
 d\rho(q)=\frac{q^{1/2}e^{-q}}{\Gamma(3/2)}dq .
 \tag{2}
\] Let \(\mathcal H_{\rm pol}\) be the Hilbert space of (1). For
finite-support families \(f=(f_b)\), define \[
 \mathcal I f=\sum_b |b\rangle\otimes f_b
       \quad\hbox{in }\mathcal H_{\rm pol}\otimes L^2(\rho).
 \tag{3}
\] The original raw norm is preserved exactly:
\(\|\mathcal I f\|^2=\sum_b\|f_b\|^2\). Finite-support families are
dense in the direct sum, and simple tensors with finite basis support
are dense in the tensor product. Thus (3) extends to a unitary onto the
whole tensor product. Its inverse sends \(|b\rangle\otimes f\) to \(f\)
in fibre \(b\), zero elsewhere. No measure or amplitude has been
rescaled.

Write \(P f_b=b f_b\) on the maximal label-moment domain, and
\(Qf_b(q)=qf_b(q)\) on its maximal moment domain. Let \[
 (U(a)f)_b=f_{b-a},\qquad
 (Z(\ell)f)_b=e^{i\ell b}f_b,\qquad
 (W(\alpha)f)_b(q)=e^{i\alpha q}f_b(q).
\] Testing these formulas on the dense finite-support space gives \[
 \begin{split}
 \mathcal I U(a)\mathcal I^{-1}
   &=\widehat V(-a)\otimes I,\\
 \mathcal I Z(\ell)\mathcal I^{-1}
   &=\widehat U(\ell)\otimes I,\\
 \mathcal I W(\alpha)\mathcal I^{-1}
   &=I\otimes e^{i\alpha Q},\\
 \mathcal I P\mathcal I^{-1}
   &=\widehat x\otimes I .
 \end{split}
 \tag{4}
\] The first three identities extend by boundedness. The last extends by
the equality of the maximal squared label moments. The exact Weyl phase
is \[
 Z(\ell)U(a)=e^{i\ell a}U(a)Z(\ell),
 \tag{5}
\] as follows by evaluating both sides in fibre \(b\). The central
modulation \(W\) commutes with both. The group \(Z\) is strongly
continuous by dominated convergence in the sum of squared fibre norms;
\(U\) is not, since a nonzero shift sends a single-fibre vector to an
orthogonal fibre. These properties are preserved by (3).

\subsection{3. Dynamics and what the comparison does not
identify}\label{weak-coupling-part-8-dynamics-and-what-the-comparison-does-not-identify}

The actual radial limit has physical Hamiltonian \[
 Hf_b(q)=c\,b^2q f_b(q),\qquad c=100\sqrt2\pi,
\] on \[
 \operatorname{Dom}H=
 \left\{f:\sum_b c^2b^4\int q^2|f_b(q)|^2d\rho(q)<\infty\right\}.
 \tag{6}
\] Equations (3)--(4) therefore identify it with the maximal joint
multiplication operator \(c\,\widehat x^2\otimes Q\). This notation
means the maximal joint operator in (6), not an unclosed product with an
artificially smaller domain. Its nonreal resolvents are fibrewise
multiplication by \((cb^2q-z)^{-1}\); their norm is at most
\(1/|\operatorname{Im}z|\).

The source's label \(x\) thus corresponds exactly to our global carrier
label \(b\), not automatically to a physical spatial coordinate of the
original gauge field. The additional Gamma variable \(q\), its measure,
and the coefficient \(c\) arise from the specific radial state and
physical lattice scaling. They are not parameters supplied by the
kinematic identification (1). The source paper's particle Hamiltonian is
not asserted equal to (6).

In particular the representation is not novel merely because its carrier
basis is uncountable or its shifts are nonregular. The specific results
established here are the map from the actual nonlinear regulators to
that representation, the retained Gamma factor and Hamiltonian, and the
explicit relation of its amplitudes and local-cylinder observables. The
remaining physical question is which spatially local gauge observables,
with all their scale factors retained, survive on it. Neither the
unitary map (3) nor the established polymer terminology answers that
question.

\clearpage
\part{Earlier Gamma state and local-observable continuation}
\section{Exact Gamma state and the next local-observable
calculation}\label{weak-coupling-part-9-exact-gamma-state-and-the-next-local-observable-calculation}

Later checkpoint, 8 September 2026: the full local-spectrum task below
has now been completed. Use
WEB\_PRO\_YM\_LOCAL\_SPECTRUM\_AND\_INTERACTION\_FEEDBACK\_20260908.md
for the new exact results and the next bounded calculation. The original
feedback is retained below as task provenance.

8 September 2026. The material below supersedes the earlier radial-state
stopping point. The target remains nontrivial interacting
four-dimensional Yang--Mills, not a free comparison or an isolated
spectral model.

\subsection{The retained finite
system}\label{weak-coupling-part-9-the-retained-finite-system}

\[
 H_j=\frac{2g_j^2}{a_j}\sum_eE_e+
       \frac1{2g_j^2a_j}\sum_p(2-\operatorname{tr}U_p),\quad
 L_j=j^2,\quad a_j=\frac1{100j},\quad j\ge2.
\] All vertex gauge constraints, plaquettes and the actual positive unit
vacuum are retained. The physical radial phase has comparison form
\(B_j(b)=e^{ib\sqrt j q}\), where \[
 q=\frac{\sigma_j|z|^2}{4},\quad
 \sigma_j=\sqrt8\sin\frac{\pi}{4j^2+2},\quad
 \lambda_j=\frac{2\sigma_j}{a_j},\quad k=\frac32,\quad
 j\lambda_j\longrightarrow c=100\sqrt2\pi.
\] Actual centering is always with the actual finite-coupling vacuum.

For \(b=1\), the raw spectral comparison measure is \[
 \sum_{n\ge1}\frac{(k)_n}{n!}\frac{j^n}{(j+1)^{n+k}}
                    \delta_{n\lambda_j}.
\] It converges to \(\omega^{k-1}e^{-\omega/c}d\omega/(\Gamma(k)c^k)\),
with limiting mass one and no zero atom. A Poisson--Gamma coupling gives
the exact bound \[
 d_{\rm BL}\le
 [kc_j^2/j+k(k+1)(c_j-c)^2]^{1/2}+(j+1)^{-k},
 \qquad c_j=j\lambda_j .
\] The proved fixed-box nonlinear convergence selects a strictly
positive dyadic \(g_j\) matching raw measures and norms. No
uniform-in-volume remainder or prescribed running coupling is asserted.

\subsection{The full phase and time
kernel}\label{weak-coupling-part-9-the-full-phase-and-time-kernel}

The exact centered mixed kernel is \[
 K_j(b,d;t)=
 [1+bdj(1-e^{-t\lambda_j})+i(b-d)\sqrt j]^{-k}
 -[(1+ib\sqrt j)(1-id\sqrt j)]^{-k}.
\] Both principal branches and the vacuum subtraction are proved in
full. Around a nonzero carrier, set \(b=b_0+\alpha/\sqrt j\) and
\(d=b_0+\delta/\sqrt j\). Its limit is \[
 [1+cb_0^2t+i(\alpha-\delta)]^{-k}.
\] It is represented on \(L^2(\rho)\), with
\(d\rho=q^{k-1}e^{-q}dq/\Gamma(k)\), by \(e^{i\alpha q}\) and
\(H_{b_0}=cb_0^2Q\), \(Qf=qf\) on its maximal domain. Continuous
modulations of the carrier survive. The full family is
\(\bigoplus_{b_0\in\mathbb R}L^2(\rho)\); global carrier shifts are
nonregular, whereas modulation within fibres is strongly continuous. The
origin fibre has zero Hamiltonian on every vector. It contains
additional zero-energy modulation states beyond the original vacuum.
Thus keeping only a nonzero carrier plus one vacuum would omit some
states generated by the common modulation algebra.

\subsection{An operator-level local
map}\label{weak-coupling-part-9-an-operator-level-local-map}

The actual positive-coupling diagonal may also impose \(g_j<j^{-5}\).
Ordered coarsening uses \(k\mid j\); a factorial cofinal subsequence
retains every fixed coarse box without fractional lattice paths. The
actual vacuum coarse-face defect is bounded by
\(216\sqrt5\,g_j^2j^8/k^2\). Hence for every fixed continuous
gauge-invariant coarse cylinder \(F_j\), \(\|(F_j-F(I))\psi_j\|\to0\).

For every centered bounded phase state \(v_j\), positivity of the full
nonlinear heat semigroup and its fixed positive vacuum prove \[
 |e^{-t(H_j-\mathcal E_j)}v_j|\le2\psi_j,\qquad t\ge0.
\] Consequently \[
 \left\|(F_j-F(I))\sum_\ell a_\ell e^{-t_\ell(H_j-\mathcal E_j)}v_j\right\|
 \le2\sum_\ell|a_\ell|\,\|(F_j-F(I))\psi_j\|\longrightarrow0 .
\] Thus the unchanged bounded cylinder algebra acts as \(F(I)I\) on the
whole limiting Gamma time-orbit space. This does not determine rescaled
curvature or electric observables, whose norms are not uniformly
bounded.

\subsection{Bounded demanding
continuation}\label{weak-coupling-part-9-bounded-demanding-continuation}

Start with the full weighted local energy for one fixed smooth compactly
supported spatial test function, including electric and magnetic terms,
actual vacuum centering, and every original spacing/coupling factor. At
a fixed positive physical Euclidean filtering time, derive its complete
two-mode spectral kernel from the original edge and face weights and
calculate its spatial continuum limit with raw norms and remainder
bounds. Then calculate matrix elements with the carrier states above.

Determine whether this observable detects a nonconstant operator on
their limiting sector; if a component disappears, identify its exact map
and retain the next nonzero component. Do not assume a volume-uniform
finite-coupling estimate or equate an oscillator calculation with the
interacting continuum. The newest full local-energy/tensor sidebar
provides the discrete electric and magnetic pair matrices for this task.

Complete current proof sources in the canonical workbench:

\begin{itemize}
\tightlist
\item
  weak\_coupling\_radial\_gamma\_continuum\_candidate.md
\item
  amplitude\_time\_kernel\_independent\_audit.md (public mathematical
  proof)
\item
  radial\_gamma\_configuration\_bridge.md
\item
  the full finite-box, state-map and electric-graph companions.
\end{itemize}

These are now Parts 4--6 of the existing weak-coupling draft, following
the corrected Parts 1--3. The exact-check file contains 23 finite
algebra identities and 7 separately labelled numerical illustrations,
not a certificate of an interacting continuum or a Millennium
conclusion.

\clearpage
\part{Compact local energy and the full-mode radial bridge}
\section{Compact local energy, radial electric observables, and a
full-mode low-energy lower
bound}\label{weak-coupling-part-10-compact-local-energy-radial-electric-observables-and-a-full-mode-low-energy-lower-bound}

8 September 2026.

The compact local-energy state and the global radial electric observable
admit an exact finite-regulator double-commutator relation and an exact
common lowest-mode quadratic representation. Their different raw powers
are computed below with the original coefficients and measures. For a
nonnegative nonzero compact profile, a growing family of the original
transverse modes then gives a strictly positive lower bound for the
unprojected local-energy spectral weight in every window
\((0,\epsilon)\), \(\epsilon>0\). A specified refinement of the
positive-coupling diagonal transfers that bound to the nonlinear
finite-box Hamiltonians. These statements do not identify an interacting
continuum theory or verify a Navier--Stokes singularity.

\subsection{1. Original objects and operator
domains}\label{weak-coupling-part-10-original-objects-and-operator-domains}

Vertices remain the integer points of \(\{-L,\ldots,L\}^3\), with all
positively oriented links whose endpoints belong to this cube and all
its elementary faces. Physical coordinates are \(x=an\), edge midpoints
are \(am(e)\), and \[
 N=2L+1,\qquad \ell=aN.
\] The endpoint distance in each coordinate is \(2La=\ell-a\). No
formula below replaces this distance by \(\ell\); the latter is the
length appearing in the exact open difference basis.

On product Haar probability measure, with \(T_\alpha=-i\sigma_\alpha/2\)
and the original left fields \(X_{e,\alpha}\), retain \[
 E_e=-\sum_{\alpha=1}^3X_{e,\alpha}^2,\qquad
 H=\kappa\sum_eE_e+b\sum_p(2-\operatorname{tr}U_p),\qquad
 \kappa=\frac{2g^2}{a},\quad b=\frac1{2g^2a}.
\tag{A1}
\] The complete Wilson scalar is \(2b|\mathsf P_L|\). The physical
Hilbert space is the vertex-gauge invariant subspace, with the full
residual simultaneous colour rotation constraint after the exact
rooted-tree quotient. The two source gauge conventions are identical
under \(k_v=h_v^{-1}\), with the identity map on link configurations and
wavefunctions.

The Hamiltonian has domain \(H^2\) and form domain \(H^1\) on the
compact product group, restricted to the invariant subspace. Its
positive unit vacuum is \(\psi_g\), its energy is \(\mathcal E_g\), and
its excitation operator is \(A_g=H-\mathcal E_g\). For real edge
weights, \[
 f_p=\frac14\sum_{e\in\partial p}f_e,\qquad
 D_f=\kappa\sum_ef_eE_e+b\sum_pf_p(2-\operatorname{tr}U_p),
 \quad C_{g,f}=D_f-\langle D_f\rangle_{\psi_g},\quad
 \Xi_f=C_{g,f}\psi_g.
\tag{A2}
\] The self-adjoint domain of the signed kinetic sum is its Peter--Weyl
diagonal domain: the squared coefficients times
\(\kappa^2(\sum_e f_ej_e(j_e+1))^2\) must be summable. It need not be
\(H^2\). Adding the bounded real Wilson multiplier preserves this
domain. Smooth physical functions belong to it and are preserved by
every operator in the commutator identities below. Thus those identities
have a specified common domain without an assertion that every signed
\(D_f\) is elliptic.

The comparison oscillator uses all transverse modes and all three
colours: \[
 H_{\rm osc}=\frac1a\sum_{\nu,\alpha}
 \left(-2\partial_{z_{\nu,\alpha}}^2+
       \frac{\sigma_\nu^2z_{\nu,\alpha}^2}{8}\right),\qquad
 A_0=\sum_{\nu,\alpha}\frac{\sigma_\nu}{a}
                         a_{\nu,\alpha}^\dagger a_{\nu,\alpha}.
\tag{A3}
\] The tree, Haar-density, dilation and constant-Jacobian
linear-coordinate maps are the source isometries, not replacements of
the counting or Haar measures. Colour-contracted creation pairs are
invariant under the residual gauge group.

Let \(V\) be an orthonormal transverse edge eigenbasis and
\(Y=d_1V\Sigma^{-1}\), \(\Sigma=\operatorname{diag}(\sigma_\nu)\). Both
\(V\) and \(Y\) are isometries for the original counting inner products.
Define \[
 A^f=V^{\mathsf T}\operatorname{diag}(f_e)V,\qquad
 F^f=Y^{\mathsf T}\operatorname{diag}(f_p)Y.
\] The source full matrices can equivalently be written \[
 \mathsf R_f=\Sigma^{1/2}(F^f-A^f)\Sigma^{1/2},\qquad
 \mathsf Q_f=\Sigma^{1/2}(F^f+A^f)\Sigma^{1/2}.
\tag{A4}
\] In particular both electric and magnetic contributions remain in
\(\mathsf R_f\). For constant \(f\), \(F^f=A^f=fI\), and its vacuum
creation matrix vanishes exactly, as required by \(D_f=fH\).

\subsection{2. Compact-support powers and
constants}\label{weak-coupling-part-10-compact-support-powers-and-constants}

Retain \(L_j=j^2\), \(a_j=1/(100j)\), \(N_j=2j^2+1\), and \[
 \ell_j=\frac{2j^2+1}{100j},\quad
 \sigma_j=\sqrt8\sin\frac{\pi}{2N_j},\quad
 \Delta_j=\frac{2\sigma_j}{a_j},\quad
 \frac{\ell_j\sigma_j}{a_j}\longrightarrow\sqrt2\pi,\quad
 \frac{\ell_j}{j}\longrightarrow\frac1{50}.
\tag{A5}
\] For \(h\in C_c(\mathbb R^3;\mathbb R)\), let
\(f_{h,j}(e)=h(a_jm(e))\), with the full quarter-face average, and put
\[
 H_h=\int h(x)\,dx,\qquad H_{h,ik}=\int x_ix_kh(x)\,dx.
\] The three source lowest modes \(V_i\) have the original cyclic signs.
On the face plane normal to \(i\), their unit curls obey exactly \[
 Y_i(n)^2=\frac4{N_j^3}
 \cos^2\frac{\pi a_j(n_p+1/2)}{\ell_j}
 \cos^2\frac{\pi a_j(n_q+1/2)}{\ell_j}.
\tag{A6}
\] Different \(Y_i\) have disjoint face-plane supports. A face-center
Riemann sum therefore gives \[
 \ell_j^3F^h_{ii}\longrightarrow4H_h,\qquad F^h_{ik}=0\ (i\ne k).
\tag{A7}
\] The face weight differs from the value at its center by at most the
modulus of continuity of \(h\) at \(a_j/2\). Its support is enlarged by
at most that distance. Multiplication by the number of nonzero samples
and by \(a_j^3\) bounds this error by a constant times that modulus. The
limiting cosine factors are uniformly one on that enlarged compact
support. These statements justify (A7) with the open endpoint grids and
the quarter-face average intact.

The two electric directional squares for \(V_i\) each have coefficient
\(2/N_j^3\), one cosine square and one sine square. The same Riemann
argument, using \(\ell_j\sin(\pi x/\ell_j)\to\pi x\) uniformly on a
compact set, proves the additional exact leading electric moments \[
 \ell_j^5 A^h_{ii}\longrightarrow
 2\pi^2\int(x_p^2+x_q^2)h(x)\,dx.
\tag{A8}
\] For \(i\ne k\), only the remaining edge direction is shared. The
original product is negative: \[
 V_i(n,d)V_k(n,d)=-\frac2{N_j^3}
 \sin\frac{\pi a_jn_i}{\ell_j}
 \sin\frac{\pi a_jn_k}{\ell_j}
 \cos^2\frac{\pi a_j(n_d+1/2)}{\ell_j}.
\] Consequently \[
 \ell_j^5A^h_{ik}\longrightarrow-2\pi^2H_{h,ik}.
\tag{A9}
\] Since the lowest frequencies coincide, (A4) gives the exact
identities \(\mathsf R_h^{\min}=\sigma_j(F^h-A^h)\) and
\(\mathsf Q_h^{\min}=\sigma_j(F^h+A^h)\). Combining (A5)--(A9) proves \[
 \frac{\ell_j^4}{a_j}\mathsf R_h^{\min},\quad
 \frac{\ell_j^4}{a_j}\mathsf Q_h^{\min}
 \longrightarrow4\sqrt2\pi H_h I_3,
\qquad
 \frac{\ell_j^6}{a_j}(\mathsf R_h^{\min})_{ik}
 \longrightarrow2\sqrt2\pi^3H_{h,ik}\ (i\ne k).
\tag{A10}
\] The leading compact-support contribution in (A10) is magnetic. The
electric term is present and has its explicit two-power smaller
contribution (A8)--(A9). Its smallness for these global lowest modes
does not say that local electric observables vanish on all modes.

For a symmetric lowest-mode matrix \(B\), retain the raw vector \[
 \Phi(B)=\sum_{i,k,\alpha}B_{ik}
     a_{i,\alpha}^\dagger a_{k,\alpha}^\dagger\Phi_0,
 \qquad \langle\Phi(B),\Phi(E)\rangle=6\operatorname{tr}(B^*E).
\tag{A11}
\] The source lowest-band vector is \(\Phi(\mathsf R_h^{\min}/(4a_j))\).
Its raw mass is exactly \[
 d_{L_j,a_j}(f_{h,j})=\frac3{8a_j^2}
                      \operatorname{tr}((\mathsf R_h^{\min})^2).
\] Thus \[
 \ell_j^8d_{L_j,a_j}(f_{h,j})\longrightarrow36\pi^2H_h^2.
\tag{A12}
\] The source actual Gram frame has \(G_{g_j}\to I\) and the chosen
error is smaller than every amplification used here. Hence for its
actual projected state \(v_{h,j}=P_{g_j}\Xi_{f_{h,j}}\), \[
 \ell_j^4M_{g_j}^{-1}v_{h,j}\longrightarrow
 (0,\sqrt2\pi H_h I_3),\quad
 \ell_j^8\|v_{h,j}\|^2\longrightarrow36\pi^2H_h^2.
\tag{A13}
\] Its raw mass tends to zero. Its raw first moment, keeping the
original physical Hamiltonian, satisfies \[
 \ell_j^9\langle v_{h,j},A_{g_j}v_{h,j}\rangle
 \longrightarrow72\sqrt2\pi^3H_h^2.
\tag{A14}
\] The finite-band energy is \(\Delta_j\), up to the stipulated error;
\(\ell_j\Delta_j\to2\sqrt2\pi\). Therefore fixed physical-time
correlations have the rescaled limit \(36\pi^2H_hH_k\), and at
\(t=\ell_j\theta\) they acquire the factor \(e^{-2\sqrt2\pi\theta}\).
These are precisely the source powers. No division of a state by its
norm is involved.

When \(H_h=0\), (A12) has zero right side; it supplies no positive
leading mass. The mixed moments (A9)--(A10), the other moments (A8), and
the full matrices are still retained. In particular vanishing integral
does not imply that \(\Xi_f\) is zero.

The coefficients in source (144)--(145) also agree with its frame and
metric. In particular \(\operatorname{tr}S_*=337/64\) and
\(\operatorname{tr}B_*=-8425\sqrt2/(32\pi)\) give
\(6\sqrt2\pi H_h\operatorname{tr}B_*=-25275H_h/8\). These equations
concern \(w_j=P_{g_j}\Xi_{f_{S_j}}\), defined in source (127). They do
not compute the expectation in the unprojected vector \(\Xi_{f_{S_j}}\),
although the weight retains the entire original material tensor.

\subsection{3. Exact finite-regulator relation to the radial electric
observable}\label{weak-coupling-part-10-exact-finite-regulator-relation-to-the-radial-electric-observable}

For a selected original lowest mode \(i\), let \(Q_{g,i}=g^{-2}F_i\) be
the Section 22 cutoff quadratic. It is a globally smooth bounded
gauge-invariant multiplier for every fixed \(g>0,L,a\). Define \[
 R_{g,i}=\frac12[Q_{g,i},[A_g,Q_{g,i}]]
        =\kappa\sum_{e,\alpha}(X_{e,\alpha}Q_{g,i})^2.
\tag{A15}
\] Every derivative of the cutoff remains in the right side. The full
Wilson multiplication commutes with \(Q_{g,i}\); this explains its zero
commutator without removing it from \(A_g\).

There is an exact weighted version involving the unchanged full local
energy: \[
 \boxed{\frac12[Q_{g,i},[D_f,Q_{g,i}]]
       =\kappa\sum_{e,\alpha} f_e(X_{e,\alpha}Q_{g,i})^2
       =:R_{g,i;f}.}
\tag{A16}
\] For its proof, on a smooth physical function \(u\), \[
 [E_e,Q]u=-\sum_\alpha\{(X_{e,\alpha}^2Q)u
                    +2(X_{e,\alpha}Q)X_{e,\alpha}u\}.
\] Commuting \(Q\) with this expression leaves precisely
\(2\sum_\alpha(X_{e,\alpha}Q)^2u\). Multiplication by \(\kappa f_e/2\)
and summation proves (A16). The magnetic term in (A2) commutes with
\(Q\) exactly. The right side extends to a bounded multiplier at each
fixed regulator. For \(f_e\ge0\), it is nonnegative; for arbitrary real
weights, \[
 |R_{g,i;f}|\le\|f\|_\infty R_{g,i}
\tag{A17}
\] as a pointwise inequality of functions. This supplies an exact
relation on the original physical Hilbert space before any projection or
limit.

The cutoff is a configuration-space cutoff in all chord logarithms. It
is not a compact physical-space support condition. The selected
transverse cochain is global; taking a derivative square does not by
itself convert this phase coordinate into a local field. Conversely
(A16) gives its exact local energy-weighted derivative distribution, so
global dependence is not a claim of unrelatedness.

At fixed box the Section 22 tail proof gives, with the cutoff and its
derivative terms retained, \[
 \mathcal B_g(R_{g,i}^m\psi_g)\longrightarrow
 (\lambda q_i)^m\Phi_0,\qquad
 \lambda=\frac{2\sigma}{a},\quad q_i=\frac\sigma4|z_i|^2.
\tag{A18}
\] It uses the global bound \(S\le C_SW\), all actual vacuum moments of
\(W/g^2\), and the weighted chart-tail inequality; its constants depend
on the fixed box and cutoff. It is not a volume-uniform estimate.

The same explicit tangent metric computation and (A17) give the
fixed-box weighted limit \[
 R_{g,i;f}\quad\leadsto\quad A^f_{ii}\lambda q_i
\tag{A19}
\] on the vacuum, with strong transport of this multiplier times the
vacuum. Indeed \(\nabla q_i=(\sigma/2)z_i\) in the selected kinetic
mode, the weighted tangent Gram entry is \(A^f_{ii}\), and the unchanged
kinetic coefficient is \(2/a\). Their product is
\((2/a)(\sigma^2/4)A^f_{ii}|z_i|^2=A^f_{ii}\lambda q_i\). The tails
follow by (A17) and the already proved tails of \(R_{g,i}\psi_g\). For
compact \(h\), (A8) specifies the \(\ell^{-5}\) factor in this weighted
radial distribution. That factor differs from the leading full
local-energy factor because (A16) measures the electric phase-gradient
contribution, whereas (A2) includes the magnetic energy as well.

\subsection{4. Exact oscillator bridge, including its raw vacuum
image}\label{weak-coupling-part-10-exact-oscillator-bridge-including-its-raw-vacuum-image}

This section proves an operator relation before taking the
compact-support asymptotic. It is confined to the common lowest-mode
oscillator factor and its projection; it is not an assertion about
arbitrary full operator words.

Let \(B\) be any real symmetric \(3\)-by-\(3\) matrix and define on the
three lowest modes \[
 q_B=\frac\sigma4\sum_\alpha z_\alpha^{\mathsf T}Bz_\alpha,
 \qquad \mathcal N_B=\sum_{i,k,\alpha}B_{ik}
                         a_{i,\alpha}^\dagger a_{k,\alpha}.
\] The polynomial multiplier \(q_B\) has expectation
\(\langle q_B\rangle=(3/2)\operatorname{tr}B\). On the invariant
Schwartz core, the creation formulas give \[
 \lambda\left(q_B-\frac32\operatorname{tr}B\right)
 =\frac\sigma a\sum_{i,k,\alpha}B_{ik}
       (a_{i,\alpha}a_{k,\alpha}
          +a_{i,\alpha}^\dagger a_{k,\alpha}^\dagger)
   +\frac{2\sigma}a\mathcal N_B.
\tag{A20}
\] For signed \(B\), this is the equality of the indicated polynomial
operators on that core; no positivity of \(B\) is assumed. The
multiplier has its maximal multiplication domain. The finite-mode
creation expressions are well-defined on every polynomial times Gaussian
used below.

Let \(C^{\rm low}_{0,f}\) mean the terms in the source full centered
oscillator whose two mode indices both belong to the three lowest modes.
With the entire full vacuum scalar already subtracted, the exact
identity is \[
 \boxed{
 C^{\rm low}_{0,f}
 =\frac14\lambda\left(q_{F^f-A^f}
                 -\frac32\operatorname{tr}(F^f-A^f)\right)
       +\frac\sigma a\mathcal N_{A^f}.}
\tag{A21}
\] Here \(F^f,A^f\) denote their lowest blocks. To prove it, the pair
coefficient from (A20) is \(\sigma(F^f-A^f)/(4a)\), and the number
coefficient is \(\sigma(F^f-A^f)/(2a)+\sigma A^f/a
 =\sigma(F^f+A^f)/(2a)\). These are exactly source (134), with
\(\mathsf R=\sigma(F-A)\) and \(\mathsf Q=\sigma(F+A)\). Thus (A21) has
retained both contributions and the number-preserving term.

Terms with indices outside these three modes vanish after compression on
both sides to vacuum plus the lowest band; their uncompressed action is
the source higher-mode map (147)--(148). In particular
\(\mathcal N_{A^f}\Phi_0=0\), so the exact lowest-band vacuum state is
already the radial quadratic image \[
 P_{\min}C_{0,f}\Phi_0
 =\frac14\lambda\left(q_{F^f-A^f}
             -\frac32\operatorname{tr}(F^f-A^f)\right)\Phi_0.
\tag{A22}
\] Every such symmetric quadratic is a real linear combination of
single-mode radial quadratics: an orthogonal eigenbasis \(b_r\) of its
coefficient matrix gives
\(q_B=\sum_r\beta_r(\sigma/4)|\sum_i(b_r)_iz_i|^2\). This is a proved
map within the degenerate lowest transverse space, preserving the three
colour coordinates and the displayed raw coefficients.

For the compact weights, (A7)--(A9) imply on the finite vacuum-plus-band
space \[
 \ell_j^4C^{\rm comp}_{0,h,j}
 -H_h\ell_j\sum_{i=1}^3
     \left(R^{\rm comp}_{0,i}-\langle R_{0,i}\rangle\right)
 \longrightarrow0,
 \qquad R_{0,i}=\lambda_jq_i.
\tag{A23}
\] Indeed the leading coefficient of \(F^h-A^h\) is
\(4H_h\ell_j^{-3}I\); the factor \(1/4\) in (A21) cancels this four. The
number correction has size
\(\ell_j^4(\sigma_j/a_j)O(\ell_j^{-5})=O(\ell_j^{-2})\). The remaining
matrix error tends to zero by (A7) and \(\ell_j\lambda_j\to2\sqrt2\pi\).
The same proof applies to every component of the finite compressed
operator, not only its vacuum vector.

Write \[
 \mathcal K=\mathbb C\oplus\operatorname{Sym}_3(\mathbb C),\qquad
 \langle(c,B),(d,E)\rangle=\overline c d+6\operatorname{tr}(B^*E),
 \quad E_{ii}=\text{the matrix unit},\quad c_\Gamma=100\sqrt2\pi.
\] The limiting compressed centered radial operator associated to
\(jR_{g,i}\), whenever all relevant finite-box tests are included on the
same diagonal, is \[
 \mathcal R_i(c,B)=c_\Gamma\left(
       3B_{ii},\ \frac c2E_{ii}+E_{ii}B+BE_{ii}\right).
\tag{A24}
\] Its vacuum creation coefficient is \(c_\Gamma E_{ii}/2\). Pair
annihilation contributes \(3c_\Gamma B_{ii}\) because of the factor-six
inner product (A11); the number commutator gives the final two terms.
Consequently the source compact operator has the precise relation \[
 \boxed{\mathcal T_h=\frac{H_h}{50}\sum_{i=1}^3\mathcal R_i,
 \qquad
 \mathcal T_h(c,B)=\sqrt2\pi H_h(6\operatorname{tr}B,cI_3+4B).}
\tag{A25}
\] The coefficient \(1/50\) is the actual limit \(\ell_j/j\). Formula
(A25) would have a wrong coefficient if \(a_j\), \(N_j\), or the
amplitude \(jR_g\) were replaced without recording that change.

Section 22 explicitly writes one selected radial mode. Formula (A25)
uses the three exact lowest modes. Each of the other two cutoff
quadratics is defined by (279) with its corresponding selected row of
the same orthogonal mode matrix. The fixed-box cutoff, potential-moment
and finite spectral arguments apply separately to these two smooth
invariant functions. There are only three functions at a fixed stage.
Their proved tests, the source compact-energy tests, and mixed
finite-frame tests can therefore be included together in a refined
positive dyadic sequence. This is not a claim that either earlier
least-dyadic choice already satisfies their union.

For one selected mode, \(q_i\) has Gamma law with shape \(k=3/2\), so
its centered raw vacuum mass is \[
 \|(R_{0,i}-\langle R_{0,i}\rangle)\Phi_0\|^2
   =\frac32\lambda_j^2,
 \qquad
 \|j(R_{0,i}-\langle R_{0,i}\rangle)\Phi_0\|^2
   \longrightarrow\frac32c_\Gamma^2=30000\pi^2.
\tag{A26}
\] The three centered radial vectors are orthogonal, since the
oscillator factors are independent and each is centered. Thus their
summed limiting mass is \(90000\pi^2\). Multiplying by \((H_h/50)^2\)
gives \(36\pi^2H_h^2\), exactly (A12), and verifies the raw-mass
compatibility of (A25). The mixed limiting Gram entry for one radial
mode is \[
 6\operatorname{tr}\left((\sqrt2\pi H_hI_3)^*
                                      \frac{c_\Gamma}2E_{ii}\right)
 =600\pi^2H_h.
\tag{A27}
\] On an actual common refined diagonal this is the limit of
\(\ell_j^4\langle P_{g_j}\Xi_{f_{h,j}},
 j(R_{g_j,i}-\langle R_{g_j,i}\rangle)\psi_{g_j}\rangle\), once its
finite-frame and weighted-vector tests are included. The unprojected
radial vector may have a small remainder at each fixed coupling; the
fixed-box strong limit and included mixed tests control that remainder.
No unsupported rate on the earlier diagonal is used.

The radial centered vacuum vector tends spectrally to a positive atom at
zero after multiplication by \(j\), while its high-amplitude
phase-carrier family has the continuous Gamma excitation limit. Those
are two explicitly different state images of the same multiplier and
phase construction. Neither is identical to the compact local energy
state without the maps and factors in (A21)--(A27).

\subsection{5. A full-mode low-energy calculation from the actual
open-box
matrices}\label{weak-coupling-part-10-a-full-mode-low-energy-calculation-from-the-actual-open-box-matrices}

The full comparison raw spectral measure of the unprojected
compact-energy state is \[
 \nu_{h,0,j}
 =\frac3{8a_j^2}\sum_\mu (\mathsf R_h)_{\mu\mu}^2
                         \delta_{2\sigma_\mu/a_j}
 +\frac3{4a_j^2}\sum_{\mu<\nu}(\mathsf R_h)_{\mu\nu}^2
                         \delta_{(\sigma_\mu+\sigma_\nu)/a_j}.
\tag{A28}
\] For any Borel energy set \(I\), the equivalent ordered sum is \[
 \nu_{h,0,j}(I)=\frac38\sum_{\mu,\nu}
 \frac{\sigma_\mu\sigma_\nu}{a_j^2}
 \left[
   \sum_pf_{h,j}(p)Y_\mu(p)Y_\nu(p)
  -\sum_ef_{h,j}(e)V_\mu(e)V_\nu(e)
 \right]^2
 \mathbf1_I\left(\frac{\sigma_\mu+\sigma_\nu}{a_j}\right).
\tag{A29}
\] This finite sum includes every original transverse mode and both
terms before squaring. Restricting the two indices gives a valid lower
bound because each summand is nonnegative. An orthogonal change of basis
inside a degenerate frequency space preserves this ordered
Hilbert--Schmidt sum and the energy test. Thus one may choose the
explicit transverse eigenvectors below and extend them to the complete
eigenbasis without changing the full quantity.

Here is an evaluated lower bound for (A29), rather than only a proposed
spectral test. It applies to any fixed nonnegative nonzero
\(h\in C_c(\mathbb R^3)\). Put \(H_h=\int h>0\). Choose \(R\ge1\) so
that \(\operatorname{supp}h\subset[-R+1,R-1]^3\), increasing \(R\) if
necessary, and fix any \(\epsilon>0\). Define \[
 \delta=\min\left\{\frac\epsilon{4\sqrt3},\frac1{1000R}\right\}>0.
\tag{A30}
\] All inequalities below hold for all sufficiently large \(j\), with
this fixed \(h,R,\epsilon,\delta\).

\subsubsection{5.1. Exact mode family and its gauge
constraint}\label{weak-coupling-part-10-exact-mode-family-and-its-gauge-constraint}

For \(1\le m\le N-1\), define the full open one-dimensional basis \[
 v_m(n)=\sqrt{\frac2N}\cos\left(\frac{\pi m}2+\frac{\pi mn}N\right),
 \quad -L\le n\le L,
\] \[
 w_m(n)=-\sqrt{\frac2N}\sin\left(\frac{\pi m}2+
                                      \frac{\pi m(n+1/2)}N\right),
 \quad -L\le n<L,
 \qquad s_m=2\sin\frac{\pi m}{2N}.
\tag{A31}
\] Here are the complete difference and counting-inner-product checks.
Set \(u=n+L\) and \(\theta_m=\pi m/N\). The vertex expression becomes
\(\sqrt{2/N}\cos(\theta_m(u+1/2))\), while the edge expression is
\(-\sqrt{2/N}\sin(\theta_m(u+1))\). The cosine subtraction identity
therefore gives \(dv_m=s_mw_m\). At an interior vertex, \[
 (d^*w_m)(u)=w_m(u-1)-w_m(u)
 =2\sqrt{2/N}\sin(\theta_m/2)\cos(\theta_m(u+1/2))
 =s_mv_m(u).
\] At the first vertex
\((d^*w_m)(0)=-w_m(0)=\sqrt{2/N}\sin\theta_m=s_mv_m(0)\); at the last
vertex
\((d^*w_m)(N-1)=w_m(N-2)=(-1)^m\sqrt{2/N}\sin\theta_m=s_mv_m(N-1)\).
These retain both open endpoint terms.

The finite geometric identities \[
 \sum_{u=0}^{N-1}\cos((u+1/2)\theta)
 =\frac{\sin(N\theta)}{2\sin(\theta/2)},
\qquad
 \sum_{v=1}^{N-1}\cos(\pi r v/N)
 =-\frac{1+(-1)^r}{2}
\tag{A31a}
\] hold respectively when the displayed denominator is nonzero and when
\(r\) is a nonzero integer with \(|r|<2N\). For the first, multiply the
sum of \(e^{i(u+1/2)\theta}\) by \(1-e^{i\theta}\) and take real parts.
For the second, sum \(e^{i\pi rv/N}\) from \(v=0\) to \(N-1\), take real
parts, and subtract the \(v=0\) term; even \(r\) gives \(-1\), and odd
\(r\) gives zero. The vertex cosine product identity with \(r=m-n,m+n\)
now proves orthogonality for distinct \(m,n\) and squared norm one for
each \(v_m\). In the equal-index case the constant part sums to \(N/2\)
and the \(r=2m\) part vanishes. For the edges, use
\(2\sin A\sin B=\cos(A-B)-\cos(A+B)\): for distinct \(m,n\), the two
sums in (A31a) cancel because \(m-n\) and \(m+n\) have the same parity;
for equal indices the sum is \(((N-1)-(-1))/2=N/2\). Thus every \(w_m\)
also has squared counting norm one. The vertex function \(v_0=N^{-1/2}\)
is orthogonal to every \(v_m\) by the first identity. The \(N\) vertex
and \(N-1\) edge functions are consequently complete orthonormal bases
by their respective dimensions. At \(m=1\), (A31) is exactly the source
signed pair \(v_1=-\sqrt{2/N}\sin(\pi n/N)\),
\(w_1=-\sqrt{2/N}\cos(\pi(n+1/2)/N)\).

Select triples \(\mathbf m=(m_1,m_2,m_3)\) such that \[
 m_1\equiv0\pmod4,\quad m_2\equiv m_3\equiv1\pmod4,\qquad
 \frac\delta2\le p_r:=\frac{\pi m_r}{\ell_j}\le\delta
 \quad(r=1,2,3).
\tag{A32}
\] Write \(s_r=s_{m_r}\), \(t_{\mathbf m}=\sqrt{s_2^2+s_3^2}\), and
\(\sigma_{\mathbf m}=\sqrt{s_1^2+s_2^2+s_3^2}\). Define a cochain with
the three original edge directions by \[
 (V_{\mathbf m})_1=0,\quad
 (V_{\mathbf m})_2=\frac{s_3}{t_{\mathbf m}}v_{m_1}w_{m_2}v_{m_3},
 \quad
 (V_{\mathbf m})_3=-\frac{s_2}{t_{\mathbf m}}v_{m_1}v_{m_2}w_{m_3}.
\tag{A33}
\] The factors are evaluated on their respective original vertex or edge
grids. Its divergence is zero because the two terms are
\((s_2s_3/t_{\mathbf m})v_{m_1}v_{m_2}v_{m_3}\) and its negative. Its
squared counting norm is \((s_3^2+s_2^2)/t_{\mathbf m}^2=1\). Distinct
triples are orthogonal by the tensor basis orthogonality. To check its
eigenvalue, for one fixed positive index triple write \(e_r\) for the
tensor edge cochain with the factor \(w_{m_r}\) in direction \(r\) and
the other two vertex factors. The curl of \(\sum_rb_re_r\) has
plane-\(rs\) coefficient \(s_rb_s-s_sb_r\). Expanding its squared norm
gives \[
 \sum_{r<s}|s_rb_s-s_sb_r|^2
 =\left(\sum_rs_r^2\right)\left(\sum_r|b_r|^2\right)
                       -\left|\sum_rs_rb_r\right|^2.
\tag{A33a}
\] The relations \(dv_m=s_mw_m\) and \(d^*w_m=s_mv_m\) show that \(d_1\)
and \(d_1^*\) preserve this same tensor-index block between edge and
face components. Polarization therefore identifies the complete block of
\(d_1^*d_1\) as \(|s|^2I-ss^{\mathsf T}\). The coefficient vector in
(A33) is \((0,s_3/t_{\mathbf m},-s_2/t_{\mathbf m})\), whose dot product
with \(s\) is zero. Hence
\(d_1^*d_1V_{\mathbf m}=\sigma_{\mathbf m}^2V_{\mathbf m}\), with the
asserted actual frequency. Explicitly, the three positively ordered face
components are \[
 (d_1V_{\mathbf m})_{23}=-t_{\mathbf m}v_{m_1}w_{m_2}w_{m_3},
\] \[
 (d_1V_{\mathbf m})_{12}=\frac{s_1s_3}{t_{\mathbf m}}w_{m_1}w_{m_2}v_{m_3},
 \qquad
 (d_1V_{\mathbf m})_{13}=-\frac{s_1s_2}{t_{\mathbf m}}w_{m_1}v_{m_2}w_{m_3}.
\tag{A34}
\] These also give its curl squared norm \(\sigma_{\mathbf m}^2\). The
cochain is therefore an actual transverse mode. Its two-quantum
colour-contracted creation states in (A28) satisfy the residual gauge
constraint: simultaneous orthogonal rotation of the three
creation-operator colour indices leaves their contracted sum unchanged.

\subsubsection{5.2. Signs and lower bounds before the electric--magnetic
subtraction}\label{weak-coupling-part-10-signs-and-lower-bounds-before-the-electricmagnetic-subtraction}

Put \(c_N=\sqrt{2/N}\). The congruences in (A32) preserve definite
original trigonometric signs: \[
 v_{m_1}=c_N\cos(p_1x_1),\quad w_{m_1}=-c_N\sin(p_1x_1),
\] \[
 v_{m_r}=-c_N\sin(p_rx_r),\quad w_{m_r}=-c_N\cos(p_rx_r)
 \quad(r=2,3).
\tag{A35}
\] Here \(x_r\) is the appropriate physical vertex or midpoint
coordinate in each factor. For all edges or faces with nonzero weight,
\(|x_r|\le R\) once \(a_j<1\); the enlarged face support is included by
the margin in the definition of \(R\). Hence
\(|p_rx_r|\le\delta R\le1/1000\). In particular every cosine is at least
\(1/2\), and the sign of every sine is the sign of \(x_r\).

It follows directly from (A34)--(A35) that for any two selected triples
\(\mu,\nu\), \[
 Y_\mu(p)Y_\nu(p)\ge0
\tag{A36}
\] on every weighted face, in all three face planes. On plane 23 the two
curls are negative multiples of products of three positive cosines. On
plane 12 their signs are both minus the sign of \(x_1x_3\); on plane 13
their signs are both the sign of \(x_1x_2\). This proves (A36) even when
the support crosses a coordinate plane.

For large \(j\), \(a_j\delta\le\pi\), and \[
 \frac{2}{\pi}p_r\le\frac{s_r}{a_j}\le p_r.
\] These inequalities follow from \((2/\pi)x\le\sin x\le x\) on
\([0,\pi/2]\). Consequently \[
 \frac{\sigma_{\mathbf m}}{a_j}\ge\frac\delta4,
 \qquad \frac{\sigma_{\mathbf m}}{a_j}\le\sqrt3\delta,
 \qquad \frac{t_{\mathbf m}}{\sigma_{\mathbf m}}\ge\frac18.
\tag{A37}
\] For the first and third bounds it suffices to use \(\pi<4\), one of
the lower bounds \(s_r/a_j\ge\delta/\pi\), and \(\sqrt3<2\). They are
deliberately weak inequalities with the exact frequency still specified
by (A31)--(A34).

On plane 23, (A34)--(A37) give \[
 |Y_{\mathbf m}(p)|\ge\frac{c_N^3}{64},\qquad
 Y_\mu(p)Y_\nu(p)\ge\frac1{512N_j^3}.
\tag{A38}
\] For each edge direction \(r\), its shifted Riemann sum satisfies
\(a_j^3\sum_{e\parallel r}f_{h,j}(e)\to H_h\). Thus it lies between
\(H_h/2\) and \(2H_h\) for large \(j\). The support is eventually away
from the open spatial boundary. Each weighted direction-2 or direction-3
edge then belongs to two plane-23 faces, so the exact quarter-average
gives \[
 \sum_{p\parallel23}f_{h,j}(p)
 =\frac12\left(\sum_{e\parallel2}f_{h,j}(e)
                   +\sum_{e\parallel3}f_{h,j}(e)\right)
 \ge\frac{H_h}{2a_j^3}.
\tag{A39}
\] Combining (A36), (A38), and (A39), the full magnetic matrix entry
obeys \[
 F^h_{\mu\nu}\ge\frac{H_h}{1024\ell_j^3}.
\tag{A40}
\]

For each of the two nonzero edge directions in (A33), the magnitude of a
cochain factor is at most \(c_N^3\delta R\): one sine occurs and the
polarization coefficient is at most one. Therefore \[
 |A^h_{\mu\nu}|
 \le\frac{8\delta^2R^2}{N_j^3}
          \left(\sum_{e\parallel2}f_{h,j}(e)
                     +\sum_{e\parallel3}f_{h,j}(e)\right)
 \le\frac{32H_h\delta^2R^2}{\ell_j^3}.
\tag{A41}
\] This explicitly bounds the full electric contribution before
subtracting it. Since \(\delta R\le1/1000\) and \(32/10^6<1/2048\),
(A40)--(A41) imply \[
 F^h_{\mu\nu}-A^h_{\mu\nu}
 \ge\frac{H_h}{2048\ell_j^3},\qquad
 \frac{(\mathsf R_h)_{\mu\nu}}{a_j}
 \ge\frac{\delta H_h}{8192\ell_j^3}.
\tag{A42}
\] Every retained selected pair has physical energy at most
\(2\sqrt3\delta\le\epsilon/2\), by (A30), (A37). Its energy is strictly
positive.

\subsubsection{5.3. Counting the growing family and its raw
mass}\label{weak-coupling-part-10-counting-the-growing-family-and-its-raw-mass}

The interval of allowed integers for each component in (A32) has length
\(\delta\ell_j/(2\pi)\). The count in any one residue class modulo four
is at least one fourth of this length minus one. Once its length is at
least eight, this is at least \(\delta\ell_j/(16\pi)\). All such
integers belong to \(1,\ldots,N_j-1\) for large \(j\), since
\(a_j\delta\to0\). There are therefore at least \[
 M_j\ge\left(\frac{\delta\ell_j}{16\pi}\right)^3
\tag{A43}
\] selected orthonormal modes. Applying (A42) to all \(M_j^2\) ordered
pairs in the exact full mass formula (A29) proves \[
 \boxed{
 \nu_{h,0,j}((0,\epsilon))
 \ge\nu_{h,0,j}((0,\epsilon/2])
 \ge \frac{3\delta^8H_h^2}{2^{53}\pi^6}>0
 \quad\text{for all sufficiently large }j.}
\tag{A44}
\] The coefficient follows from
\((3/8)(16\pi)^{-6}8192^{-2}=3/(2^{53}\pi^6)\). This is a lower bound on
the full raw measure, not on a state divided by its norm. It comes from
a number of mode pairs of order \(\ell_j^6\), each of squared amplitude
of order \(\ell_j^{-6}\). The lowest three modes alone do not have that
growing count and instead have the exact \(\ell_j^{-8}\) mass in (A12).
These are compatible statements about two explicitly included parts of
(A29).

Nonnegativity of \(h\) is used in (A36)--(A41). This proof does not
assert (A44) for arbitrary signed \(h\). It covers the nonzero
nonnegative compact curvature profiles described in the source without
using a singular time or asserting anything about their fluid evolution.

\subsection{6. What is proved at positive coupling, and what the current
diagonal does not
supply}\label{weak-coupling-part-10-what-is-proved-at-positive-coupling-and-what-the-current-diagonal-does-not-supply}

At each fixed \(j\), the source full-vector and finite-projection
theorems prove convergence of the entire actual raw measure of
\(\Xi_{f_{h,j}}\) to (A28), including its total mass. A continuous
compactly supported spectral test therefore converges, regardless of
degeneracies in the comparison spectrum.

For a concrete transfer of (A44), choose a continuous function
\(\varphi_\epsilon\) on \([0,\infty)\) with values in \([0,1]\), equal
to one on \([0,\epsilon/2]\), and zero on \([3\epsilon/4,\infty)\). Both
actual and comparison centered measures have no vacuum atom at zero at
each fixed regulator. Thus \[
 \nu_{h,g,j}((0,\epsilon))
 \ge\int\varphi_\epsilon(E)\,d\nu_{h,g,j}(E),\qquad
 \int\varphi_\epsilon\,d\nu_{h,0,j}
 \ge\frac{3\delta^8H_h^2}{2^{53}\pi^6}
\tag{A45}
\] for large \(j\). The exact finite calculation to preserve on the
nonlinear side is the Gram entry \[
 \left\langle\Xi_{f_{h,j}},
       \varphi_\epsilon(H_{g,L_j,a_j}-\mathcal E_{g,L_j,a_j})
                         \Xi_{f_{h,j}}\right\rangle.
\tag{A46}
\] Its fixed-box comparison is the full finite sum (A29) with
\(\mathbf1_I\) replaced by \(\varphi_\epsilon\). This test does not
insert the first-band projection.

The Section 15 dyadic conditions control its seven-dimensional
compressed operators, frame energies, and projected Fabel coordinates.
They do not contain (A46). Section 22 contains its own weighted radial
tests. Neither displayed finite list alone proves the joint full-mode
compact-energy conclusion on the other section's already chosen least
dyadic.

There is nevertheless a proved explicit continuation, with no missing
fixed-box theorem: for this fixed \(h\), at stage \(j\) include in the
union of the existing tests \[
 \left|\int\varphi_{1/n}\,d\nu_{h,g,j}
             -\int\varphi_{1/n}\,d\nu_{h,0,j}\right|<\frac1j,
 \qquad 1\le n\le j.
\tag{A47}
\] There are finitely many added scalar tests, and each converges at
fixed box by the proved raw-measure theorem. Hence they hold for every
sufficiently small positive coupling. Taking the least positive dyadic
exponent satisfying their union defines a refined sequence
\(g_j^{\sharp}\). This establishes existence rather than postulating a
rate. It retains \(\kappa_j=200j(g_j^{\sharp})^2\),
\(b_j=50j/(g_j^{\sharp})^2\), \(\xi_j=1/(4(g_j^{\sharp})^4)\), and the
complete Wilson scalar.

For each fixed positive integer \(n\), (A44)--(A47) then prove the
actual nonlinear raw bound \[
 \liminf_{j\to\infty}\nu_{h,g_j^{\sharp},j}((0,1/n))
 \ge\frac{3\delta_{1/n}^8H_h^2}{2^{53}\pi^6}>0,
 \qquad
 \delta_{1/n}=\min\left\{\frac1{4\sqrt3n},\frac1{1000R}\right\}.
\tag{A48}
\] Every window \((0,\epsilon)\), \(\epsilon>0\), contains one of these
windows \((0,1/n)\). Thus the same refined sequence has strictly
positive raw low-energy liminf in each such window. The refinement in
(A47) is a construction proved here; it is not asserted to coincide with
either section's earlier least-dyadic definition. Its existence uses
proved fixed-box convergence and does not provide a numerical
convergence rate or a numerical value of the exponent.

The construction can preserve any fixed finite collection of profiles by
including their finite scalar tests. A simultaneous assertion for an
unspecified uncountable profile class requires its own uniform or
dense-family argument; it is not inferred merely from the separately
available fixed-profile selections.

\subsection{7. Consequences and remaining continuum
limits}\label{weak-coupling-part-10-consequences-and-remaining-continuum-limits}

The compact lowest-band raw mass is
\(36\pi^2H_h^2\ell_j^{-8}+o(\ell_j^{-8})\); the radial centered raw mass
after multiplication by \(j\) tends to \((3/2)c_\Gamma^2\). Their exact
common quadratic map (A21)--(A27) accounts for both powers. The domain
of signed weighted energy remains the Peter--Weyl domain, with the
smooth common core used for the commutators; it is not asserted to be
\(H^2\).

The precise relations established here are the original nonlinear
weighted double commutator (A16), its weighted oscillator image (A19),
the exact three-mode decomposition (A21), the compact/radial compression
map (A25), its raw mass and mixed Gram checks (A26)--(A27), and the
full-mode low-energy lower bound (A44) with the explicit actual-coupling
refinement (A47)--(A48). The electric, magnetic, number-preserving,
gauge, orientation, endpoint, and raw-amplitude factors have all been
retained in those calculations.

The remaining continuum issue is not removed by either a positive radial
Gamma measure or the stronger local raw bound (A48). No common limiting
representation of the full local observables, vacuum, time evolution,
and all required domains has been constructed here. No uniform
full-state norm or ultraviolet limit is supplied by (A44). The
unprojected products and higher-mode insertions still need their own
continuum analysis; the compression formulas do not determine them. The
spatially global phase coordinate and its configuration-space cutoff
remain as specified. Neither a prescribed renormalized running-coupling
law nor an interacting four-dimensional continuum theory is identified
by choosing sufficiently small positive finite-box couplings. No claim
about a Navier--Stokes singular endpoint follows from the regular-time
profile map.

\clearpage
\part{Local-energy spacetime kernel and nonlinear coupling response}
\section{Local-energy spacetime kernel and the exact nonlinear coupling
response}\label{weak-coupling-part-11-local-energy-spacetime-kernel-and-the-exact-nonlinear-coupling-response}

8 September 2026.

The complete open-mode calculation in the companion \emph{Full
local-energy spectral measure in the expanding-box limit} produces a
spatially local energy two-point function. This note evaluates its
position-space kernel, its positive-energy continuation, and its precise
correspondence with the energy-density covariance of three free vector
fields. It then derives an exact response formula in the original
positive-coupling compact-link Hamiltonian. That formula retains the
electric term, the Wilson term, the motion of the true vacuum, the
vacuum subtraction, and the change of time evolution. It supplies a
quantitative map from the computed comparison covariance to the
nonlinear regulator covariance.

The correspondence calculated here does not identify the selected
continuum covariance with an interacting Yang--Mills theory. In
particular a free stress-tensor covariance is not advertised as a new
interacting spectrum or as a mass-gap counterexample.

\subsection{1. Original local states and the full-mode
input}\label{weak-coupling-part-11-original-local-states-and-the-full-mode-input}

On the original open cubic lattice keep all positively oriented
contained edges and elementary faces, product Haar probability, the
generators \(T_\alpha=-i\sigma_\alpha/2\), and \[
 H_g=\frac{2g^2}{a}\sum_eE_e+
                 \frac1{2g^2a}\sum_p(2-\operatorname{tr}U_p),
 \qquad E_e=-\sum_{\alpha=1}^3X_{e,\alpha}^2.
 \tag{1}
\] The Hilbert space is the invariant subspace for the full vertex gauge
group. Its unique positive unit vacuum, vacuum energy, and excitation
operator are \(\psi_g,\mathcal E_g,A_g=H_g-\mathcal E_g\). For real
\(h\in C_c^\infty(\mathbb R^3)\) take the exact weights \[
 h_e=h(am(e)),\qquad h_p=\tfrac14\sum_{e\in\partial p}h_e,
 \quad D_{h,g}=\frac{2g^2}{a}\sum_eh_eE_e+
          \frac1{2g^2a}\sum_ph_p(2-\operatorname{tr}U_p).
 \tag{2}
\] Put \(m_{h,g}=\langle\psi_g,D_{h,g}\psi_g\rangle\) and
\(\Xi_{h,g}=(D_{h,g}-m_{h,g})\psi_g\). The scalar Wilson terms in
(1)--(2) remain included. The centered states are exactly physical and
vacuum orthogonal.

The complete full-mode companion proves, with
\(L_j=j^2,\ a_j=1/(100j),\ \ell_j=a_j(2L_j+1)\), and its explicitly
selected positive dyadic couplings, that for \(t>0\) \[
 \langle\Xi_{h,g_j},e^{-tA_{g_j}}\Xi_{k,g_j}\rangle
 \longrightarrow
 C_{h,k}(t)=\frac1{320\pi^5t}\int_{\mathbb R^3}
 |p|^4e^{-t|p|}\overline{\widehat h(p)}\,\widehat k(p)\,dp,
 \quad \widehat h(p)=\int e^{-ip\cdot x}h(x)\,dx .
 \tag{3}
\] Its real-profile formula extends sesquilinearly to complex profiles.
The same companion gives the ordered two-creation map and the raw
measure \[
 d\nu_h(\omega)=\frac1{320\pi^5}
       \int_{|p|\le\omega}|p|^4|\widehat h(p)|^2\,dp\,d\omega .
 \tag{4}
\] These are inputs with full proofs in the companion, not an assumed
nonlinear continuum equation. The positive-time states
\(e^{-tA_g/2}\Xi_{h,g}\) have finite norms, with limit (3). The measure
(4) has infinite total mass for nonzero compact smooth \(h\); no
unsmeared continuum vector \(\Xi_h\) is asserted.

\subsection{2. Exact position-space
kernel}\label{weak-coupling-part-11-exact-position-space-kernel}

Define \(r=|x|\). With the Fourier convention in (3), the kernel in \[
 C_{h,k}(t)=\int_{\mathbb R^3}\int_{\mathbb R^3}
                  \overline{h(x)}\,k(y)\,K_t(x-y)\,dx\,dy
 \tag{5}
\] is exactly \[
 \boxed{K_t(x)=
 \frac{9t^4-30t^2r^2+9r^4}{\pi^4(t^2+r^2)^6}
 =\frac1{40\pi^4}\Delta_x^2(t^2+r^2)^{-2}.}
 \tag{6}
\] Here \(t\) remains the physical excitation time of (1).

To prove this, spherical integration, first for \(r>0\), gives \[
 \int_{\mathbb R^3}e^{-t|p|}e^{ip\cdot x}\,dp
 =\frac{4\pi}{r}\int_0^\infty se^{-ts}\sin(rs)\,ds
 =\frac{8\pi t}{(t^2+r^2)^2}.
 \tag{7}
\] The last integral follows by differentiating in \(t\) the imaginary
part of \(\int_0^\infty e^{-(t-ir)s}ds=(t-ir)^{-1}\). Absolute
exponential domination permits these operations and gives the value at
\(r=0\) by continuity. Multiplication by \(|p|^4\) is \(\Delta_x^2\)
under this transform. Thus the factors \((320\pi^5t)^{-1}\) and
\(8\pi t\) in (3),(7) leave \((40\pi^4)^{-1}\).

For completeness set \(u=t^2+r^2\). On radial functions in three space
dimensions, \(\Delta f(r)=f''(r)+2f'(r)/r\). It yields \[
 \Delta u^{-2}=12(r^2-t^2)u^{-4},\qquad
 \Delta^2u^{-2}
       =120(3t^4-10t^2r^2+3r^4)u^{-6}.
 \tag{8}
\] The expressions are smooth at \(r=0\); hence the radial computation
proves (6) everywhere when \(t>0\). Fubini in (3), using compact
integrable profiles and the exponentially integrable momentum weight,
proves (5).

Three exact consequences are \[
 K_t(0)=\frac9{\pi^4t^8},\qquad
 K_{bt}(bx)=b^{-8}K_t(x)\ (b>0),\qquad
 \int_{\mathbb R^3}K_t(x)\,dx=0 .
 \tag{9}
\] For the last identity \(K_t\) is integrable, its Fourier value at
zero is zero by the \(|p|^4\) factor, and (7) fixes the Fourier factors.
Alternatively integrate the two Laplacians by parts on balls and use the
explicit derivative decay. This is the continuum counterpart of
\(D_{1,g}=H_g\), whose centered vacuum vector is exactly zero. The
kernel in (6) need not be pointwise nonnegative; covariance positivity
follows from the nonnegative Fourier density, not from the sign of this
function.

At fixed \(t>0\) and for fixed compact profiles, a displacement \(Rn\),
\(|n|=1,\ R\to\infty\), has \[
 \int\!\!\int \overline{h(x)}\,k(y)
                   K_t(x-y-Rn)\,dx\,dy
 =\frac{9\overline{\int h}\int k}{\pi^4R^8}+O_{h,k,t}(R^{-9}).
 \tag{10}
\] The rational formula (6) expands uniformly for \(x,y\) in their fixed
compact supports: \(R^8K_t(x-y-Rn)=9/\pi^4+O(R^{-1})\). Integration
against \(|h(x)k(y)|\) proves (10), with no assumption of positivity of
either profile.

\subsection{3. The precise free stress-tensor
correspondence}\label{weak-coupling-part-11-the-precise-free-stress-tensor-correspondence}

For \(X=(t,x)\in\mathbb R^4\setminus\{0\}\), set \[
 I_{\mu\nu}(X)=\delta_{\mu\nu}
                     -\frac{2X_\mu X_\nu}{|X|^2},\qquad
 \mathcal I_{\mu\nu,\rho\sigma}
 =\tfrac12(I_{\mu\rho}I_{\nu\sigma}+I_{\mu\sigma}I_{\nu\rho})
                         -\tfrac14\delta_{\mu\nu}\delta_{\rho\sigma}.
 \tag{11}
\] Direct substitution, not a fit of exponents, gives \[
 \boxed{K_t(x)=
       \frac{12}{\pi^4|X|^8}\mathcal I_{00,00}(X).}
 \tag{12}
\] Indeed \(\mathcal I_{00,00}=(1-2t^2/|X|^2)^2-1/4\); putting it over
the denominator \(|X|^{12}\) produces exactly the numerator in (6).

Osborn and Petkou give (11) as their Eq. (2.23) and the free vector
coefficient \(C_T=16/S_4^2\) in Eq. (5.16), with \(S_4=2\pi^2\). Thus
each free vector contributes \(4/\pi^4\); the three retained adjoint
colour coordinates contribute \(3(4/\pi^4)=12/\pi^4\), exactly (12). The
matched coefficient is established literature, not a new value of
\(C_T\). Our new correspondence here evaluates the complete
retained-regulator covariance in that convention.

Equation (12) identifies this one energy-density covariance. It does not
prove convergence of the other stress components, their Ward identities
as operator distributions, or higher products of the interacting
regulator fields. Equality of this two-point function also does not
assert that every correlation in any theory containing it is free.

\subsection{4. Positive energy, reflected covariance, and real
time}\label{weak-coupling-part-11-positive-energy-reflected-covariance-and-real-time}

The retained joint momentum measure is \[
 d\varrho(\omega,p)=
       \frac{|p|^4}{320\pi^5}\mathbf1_{\{\omega\ge|p|\}}
                             \,d\omega\,d^3p .
 \tag{13}
\] It is a positive tempered measure: on a Euclidean ball of radius
\(R\) its mass is at most a constant times \(R^8\), by direct
integration; the same estimate controls every Schwartz test. Its support
lies in the closed forward light cone. Its Laplace transform in
\(\omega\) is precisely the Fourier kernel in (3), since
\(\int_{|p|}^\infty e^{-t\omega}d\omega=e^{-t|p|}/t\).

For compact smooth test functions \(f(t,x)\) supported at strictly
positive times, the reflected two-point quadratic form is \[
 \int_{t,s>0}\!\!\int
       \overline{f(t,x)}K_{t+s}(x-y)f(s,y)\,dx\,dy\,dt\,ds
 =
 \int\left|\int_0^\infty
              e^{-t\omega}\widehat f(t,p)\,dt\right|^2
                                      d\varrho(\omega,p)\ge0 .
 \tag{14}
\] Compact positive-time support supplies an exponential factor at
infinity, so Fubini is justified. Spatial Fourier signs agree by the
evenness in \(p\). Formula (14) is reflection positivity of this
two-point covariance; it is not the full hierarchy of
reflection-positivity inequalities for unspecified higher fields.

The positive-energy real-time distribution is \[
 W(s,x)=\lim_{\epsilon\downarrow0}K_{\epsilon+is}(x)
       =\int e^{-is\omega+ip\cdot x}\,d\varrho(\omega,p),
 \tag{15}
\] where the second integral is a distribution and fixes the sign of the
boundary value. To justify the limit, pair its Fourier form with a
Schwartz spacetime test; the factor \(e^{-\epsilon\omega}\le1\) and the
polynomial growth of (13) give dominated convergence. For
\(\operatorname{Re}z>0\) the integral is analytic and equals the
rational formula (6) with \(t=z\), first for real \(z>0\) and then by
analytic continuation.

Away from \(r^2=s^2\), the boundary values at \(s\) and \(-s\) are the
same real rational function. Therefore \[
 \operatorname{supp}\{W(s,x)-W(-s,-x)\}
                         \subset\{(s,x):r^2=s^2\}.
 \tag{16}
\] This support assertion is a direct statement about the computed
two-point commutator distribution. It does not infer an operator
commutator identity from its vacuum expectation. The behaviour at the
cone tip is determined by (13),(15); no unverified product of singular
delta distributions is used.

Thus this local covariance has an exact four-dimensional spacetime map,
positive energy and a light-cone-supported commutator expectation. Its
continuous spectrum near zero is exhibited by (4). These statements
close the two-point spacetime calculation while retaining its precisely
stated scope.

\subsection{5. Exact differentiation of the nonlinear
regulator}\label{weak-coupling-part-11-exact-differentiation-of-the-nonlinear-regulator}

Fix the original finite box and \(a>0\). No coupling limit is taken in
the identities of this section. Abbreviate \[
 K=\sum_eE_e,\quad W=\sum_p(2-\operatorname{tr}U_p),\quad
 K_h=\sum_eh_eE_e,\quad
 W_h=\sum_ph_p(2-\operatorname{tr}U_p).
 \tag{17}
\] Here \(K\) denotes an operator, not the position-space kernel
\(K_t(x)\). For \(g>0\) the exact derivatives are \[
 V_g:=\partial_gH_g=\frac{4g}{a}K-\frac1{g^3a}W,\qquad
 Z_{h,g}:=\partial_gD_{h,g}
                  =\frac{4g}{a}K_h-\frac1{g^3a}W_h.
 \tag{18}
\] Both signs and all scalar Wilson pieces are present.

The original \(H_g\) has common domain \(H^2\) on every compact positive
interval of \(g\); \(W\) is bounded and the electric coefficient is
positive. The graph norms are equivalent there. Thus
\(H_{g+\epsilon}-H_g\), composed with a fixed nonreal resolvent of
\(H_g\), is a bounded operator tending to zero in norm. More precisely,
with \(\mathscr X=H^2\) carrying a fixed equivalent graph norm and
\(R_0(z)=(H_{g_0}-z)^{-1}\), \[
 (H_g-z)^{-1}
 =R_0(z)[I+(H_g-H_{g_0})R_0(z)]^{-1}.
\] The first factor maps \(L^2\) boundedly into \(\mathscr X\), and the
bracket is differentiable in bounded operators on \(L^2\). The Neumann
expansion therefore proves differentiability in
\(\mathcal B(L^2,\mathscr X)\), not merely in \(\mathcal B(L^2,L^2)\).
Integrating the resolvent on a small circle containing only the simple
vacuum eigenvalue proves differentiability of its rank-one projection.
Applying that projection to \(\psi_g\) and taking its positive real unit
representative gives a differentiable vacuum. This use of the
unit-vacuum convention does not divide any trial state \(\Xi_{h,g}\) by
its norm. Elliptic regularity gives smooth vacuum functions, and the
differentiated equations hold in the original operator domain.

Define the reduced inverse on the actual physical space, \[
 P_g=I-|\psi_g\rangle\langle\psi_g|,\qquad
 R_g=(A_g|_{\psi_g^\perp})^{-1}P_g,\qquad
 v_g=\langle\psi_g,V_g\psi_g\rangle,\qquad
 \eta_g=R_g(V_g-v_g)\psi_g .
 \tag{19}
\] Compact resolvent and vacuum simplicity imply a strictly positive
first excitation at this fixed regulator, so \(R_g\) is bounded. No
uniform bound in \(L,a,g\) is asserted. Gauge invariance of every
operator in (19) preserves Gauss' law exactly. Differentiating
\(H_g\psi_g=\mathcal E_g\psi_g\), and using the real positive
unit-vacuum convention, proves \[
 \mathcal E_g'=v_g,\qquad
 \psi_g'=-\eta_g,\qquad
 m_{h,g}'=\langle\psi_g,Z_{h,g}\psi_g\rangle
                    -2\operatorname{Re}
                           \langle\eta_g,D_{h,g}\psi_g\rangle .
 \tag{20}
\] The identity \(\langle\psi_g,\psi_g'\rangle=0\) follows because both
vectors are real and differentiation of the unit norm gives zero real
part. Applying \(P_g\) to the differentiated eigenvalue equation then
proves (19)--(20).

Consequently the full derivative of the centered local state is \[
 \boxed{\Xi_{h,g}'=
        (Z_{h,g}-m_{h,g}')\psi_g
                       -(D_{h,g}-m_{h,g})\eta_g.}
 \tag{21}
\] The signed weighted energy has its Peter--Weyl diagonal self-adjoint
domain plus bounded Wilson multiplication, not necessarily all of
\(H^2\). Its application in (21) is legitimate: \(\eta_g\) solves an
elliptic equation with smooth right side and is smooth at this finite
regulator. All displayed vectors belong to the required common smooth
domain.

\subsection{6. A response formula retaining all intermediate
states}\label{weak-coupling-part-11-a-response-formula-retaining-all-intermediate-states}

For real \(h\), put
\(C_{g;h}(t)=\langle\Xi_{h,g},e^{-tA_g}\Xi_{h,g}\rangle\). At any
\(g,t>0\), the exact derivative is \[
 \boxed{\begin{split}
 \partial_gC_{g;h}(t)
 ={}&2\operatorname{Re}
       \langle\Xi_{h,g}',e^{-tA_g}\Xi_{h,g}\rangle\\
 &-\int_0^t
 \langle\Xi_{h,g},e^{-(t-s)A_g}
                    (V_g-v_g)e^{-sA_g}\Xi_{h,g}\rangle\,ds .
 \end{split}}
 \tag{22}
\] To verify Duhamel's term without assuming a bounded \(V_g\), work on
the smooth vectors above. Graph norm equivalence makes \(V_g\) bounded
from \(\mathscr X=H^2\) to \(L^2\). Write
\(S_\epsilon(u)=e^{-uA_{g+\epsilon}}\), \(S(u)=e^{-uA_g}\). The product
rule on a fixed smooth vector \(\xi\) gives \[
 \frac d{ds}\{S_\epsilon(t-s)S(s)\xi\}
  =S_\epsilon(t-s)(A_{g+\epsilon}-A_g)S(s)\xi,
\] and integration gives the exact sign \[
 S_\epsilon(t)\xi-S(t)\xi
  =-\int_0^tS_\epsilon(t-s)
                (A_{g+\epsilon}-A_g)S(s)\xi\,ds .
 \tag{22a}
\] Spectral calculus bounds \(\|S(s)\xi\|\) and \(\|A_gS(s)\xi\|\)
uniformly on \(0\le s\le t\). Thus this trajectory is bounded in
\(\mathscr X\). Equation (22a) first proves strong semigroup
convergence, uniform on bounded time intervals for \(\xi\in\mathscr X\);
density and the contraction bounds extend it to each fixed \(L^2\)
vector. Also \((A_{g+\epsilon}-A_g)/\epsilon\to V_g-v_g\) in
\(\mathcal B(\mathscr X,L^2)\). The set
\(\{(V_g-v_g)S(s)\xi:0\le s\le t\}\) is compact in \(L^2\), by
graph-norm continuity and compactness of the parameter interval. A
finite net and the contraction bounds make strong semigroup convergence
uniform on this set. We may therefore divide (22a) by \(\epsilon\) and
pass to the limit in its Bochner integral, including both endpoints.
This proves the unbounded-insertion derivative used in (22).
Differentiating the two external vectors supplies the first line of
(22). No projection onto a low spectral band has been inserted in its
integral or in \(R_g\).

The constant-profile check is exact even though the profile is then not
compact in the continuum: on a finite box take \(h_e=1\) for every edge.
Then \(D_h=H_g,\ m_h=\mathcal E_g,\ Z_h=V_g\). The expectation
\(\langle\eta_g,H_g\psi_g\rangle\) vanishes, so \(m_h'=v_g\). Equations
(19),(21) give \[
 \Xi_{1,g}=0,\qquad
 \Xi_{1,g}'=(V_g-v_g)\psi_g-A_g\eta_g=0.
 \tag{23}
\] Both sides of (22) vanish. Omitting the vacuum response
\(-(D_h-m_h)\eta_g\) would fail this elementary consistency test.

For each fixed box the full weak-coupling graph theorem proves
\(C_{g;h}(t)\to C_{0;h}^{L,a}(t)\) as \(g\downarrow0\), where the right
side is the complete oscillator sum, not its lowest band. The
fundamental theorem of calculus therefore gives \[
 \boxed{C_{g;h}(t)-C_{0;h}^{L,a}(t)
       =\lim_{\epsilon\downarrow0}
                    \int_\epsilon^g \partial_vC_{v;h}(t)\,dv,}
 \tag{24}
\] with the integrand explicitly (18)--(22). The integral is the
displayed improper limit. No absolute integrability or volume-uniform
derivative bound at zero coupling is inferred.

Formula (24) is the requested exact relation between the comparison
covariance and the original nonlinear one. On the specific dyadic
trajectory used in (3), the defining tests force this difference to tend
to zero at every fixed positive time. That is a calculated property of
this trajectory and these observables, not a general theorem that
non-Abelian interactions vanish on all continuum trajectories.

The next interaction calculation has a definite integrand: retain
\(V_g-v_g\) in the full semigroup, \(R_g\) in the vacuum response, and
the two derivative terms of \(D_h\), then control (24) jointly with
\(a\downarrow0\) and \(\ell\uparrow\infty\) on a physical-scale
trajectory. The present work does not supply that bound or identify an
interacting reconstruction. It does supply the exact response whose
survival must be computed, rather than replaced by an assumed remainder.

\subsection{Literature and proof
dependencies}\label{weak-coupling-part-11-literature-and-proof-dependencies}

The full-mode lattice-to-spectral calculation, original graph
convergence, state domains and Fourier factors are proved in Sections
14--16 of the included \emph{Extensive quantum blocking and local
energy} manuscript. The exact compact/radial comparison is proved
separately in \emph{Compact local energy, radial electric observables,
and a full-mode low-energy lower bound}.

H. Osborn and A. C. Petkou, \emph{Implications of conformal invariance
in field theories for general dimensions}, Annals of Physics 231 (1994),
311--362,
\href{https://doi.org/10.1006/aphy.1994.1045}{doi:10.1006/aphy.1994.1045},
\href{https://arxiv.org/abs/hep-th/9307010v2}{arXiv:hep-th/9307010v2},
Eqs. (2.23), (5.3), (5.14)--(5.16). The actual source TeX was read for
these equations. It supplies the known free-vector stress-tensor
convention, not a new proof of our regulator limit.

A. Jaffe and E. Witten, \emph{Quantum Yang--Mills Theory}, in \emph{The
Millennium Prize Problems} (2006), 129--152,
\href{https://www.claymath.org/wp-content/uploads/2022/06/yangmills.pdf}{official
problem text}, Sections 3--4. Its local-field, spacetime and
short-distance requirements explain why the two-point correspondence and
the nonlinear response must both be retained.

S. Chatterjee, \emph{Yang--Mills for probabilists},
\href{https://arxiv.org/abs/1803.01950}{arXiv:1803.01950}, the section
entitled \emph{The problem of defining the continuum limit}. Its source
TeX was consulted for the scale/correlation and Wilson-loop discussion.
No result here assigns a continuum interpretation to a coupling sequence
merely because each member has positive coupling.

\clearpage
\part{The spatial energy measure of the radial carrier}
\section{Where the radial carrier's finite energy is
located}\label{weak-coupling-part-12-where-the-radial-carriers-finite-energy-is-located}

8 September 2026.

The radial carrier has a nonzero limiting energy and a Gamma spectral
law. The full compact local-energy observable has an exact relation to
that carrier, not merely a comparison of their names or dimensions. This
note calculates its spatial distribution. The energy is spread over a
box whose physical side length tends to infinity. Every fixed compact
region receives vanishing carrier energy, while the original box-scale
coordinate \(x/\ell_j\) carries an explicit nonzero limiting density.
Both the uncentered unit phase state and the original vacuum-orthogonal
centered state are retained.

\subsection{1. Original model and
observable}\label{weak-coupling-part-12-original-model-and-observable}

Use the original open box, Haar measure, full gauge-invariant Hilbert
space, and \[
 H_g=\kappa K+bW,\quad \kappa=\frac{2g^2}{a},\
 b=\frac1{2g^2a},\quad
 K=\sum_eE_e,\quad W=\sum_p(2-\operatorname{tr}U_p).
 \tag{1}
\] The letter \(b\) here is the original Wilson coefficient, not a
radial carrier amplitude. Let \(\psi_g>0\) be the actual unit vacuum and
\(A_g=H_g-\mathcal E_g\). For real edge weights \(f\), put \[
 f_p=\tfrac14\sum_{e\in\partial p}f_e,\quad
 D_f=\kappa\sum_ef_eE_e+b\sum_pf_p(2-\operatorname{tr}U_p),
 \quad C_f=D_f-\langle\psi_g,D_f\psi_g\rangle .
 \tag{2}
\] All faces adjacent to a weighted edge occur, with their original
quarter coefficients. For constant \(f=1\), \(C_1=A_g\).

Choose one of the three specified lowest transverse modes, denoted
\(i\in\{1,2,3\}\) by its normal plane. Its actual smooth gauge-invariant
cutoff phase is \(Q_{g,i}\), exactly the quadratic chart multiplier of
the radial companion. Define \[
 B_g(\beta)=e^{i\beta Q_{g,i}},\quad
 \phi_g(\beta)=B_g(\beta)\psi_g,\quad
 m_g(\beta)=\langle\psi_g,B_g(\beta)\psi_g\rangle,\quad
 v_g(\beta)=(B_g(\beta)-m_g(\beta))\psi_g .
 \tag{3}
\] There is no rescaling by \(\|v_g\|\). The unit phase state has norm
one; \(v_g\perp\psi_g\).

\subsection{2. Exact positive local energy measure at finite
coupling}\label{weak-coupling-part-12-exact-positive-local-energy-measure-at-finite-coupling}

On the common smooth physical core, \[
 R_{g,i;f}:=\tfrac12[Q_{g,i},[D_f,Q_{g,i}]]
           =\kappa\sum_{e,\alpha}f_e(X_{e,\alpha}Q_{g,i})^2 .
 \tag{4}
\] The Wilson multiplication commutes with the phase exactly. Applying
the product rule to each electric square, including the first
commutator, gives \[
 B_g(\beta)^*D_fB_g(\beta)
 =D_f+i\beta[D_f,Q_{g,i}]+\beta^2R_{g,i;f}.
 \tag{5}
\] Its vacuum expectation of the linear term is zero. In fact
\(\psi_g,D_f\psi_g,Q_{g,i}\psi_g\) are real, and symmetry gives
\(\langle\psi_g,[D_f,Q_{g,i}]\psi_g\rangle=0\). Consequently \[
 \boxed{\langle\phi_g(\beta),C_f\phi_g(\beta)\rangle
           =\beta^2\langle\psi_g,R_{g,i;f}\psi_g\rangle .}
 \tag{6}
\] This is an exact identity in the full nonlinear model, not a removal
of the magnetic energy from its Hamiltonian.

With original edge midpoints \(am(e)\), define the positive finite
measure \[
 \mu_{g,\beta}=\beta^2\kappa\sum_e
       \left(\sum_\alpha\int\psi_g^2
                     (X_{e,\alpha}Q_{g,i})^2\,dU\right)
                       \delta_{am(e)} .
 \tag{7}
\] For \(f_e=h(am(e))\), equation (6) is exactly
\(\langle\phi_g,C_f\phi_g\rangle=\int h\,d\mu_{g,\beta}\). Every
coefficient in (7) is nonnegative. Its total mass is
\(\langle\phi_g,A_g\phi_g\rangle
 =\langle v_g,A_gv_g\rangle\), since the vacuum subtraction in (3) has
zero excitation form.

For the centered state define instead the signed finite measure \[
 \widetilde\mu_{g,\beta}
     =\sum_e\langle v_g(\beta),C_{\mathbf1_e}v_g(\beta)\rangle
                                                  \delta_{am(e)} .
 \tag{8}
\] The local operator \(C_{\mathbf1_e}\) includes its adjacent
quarter-weight faces. Linearity proves that (8) integrates any
edge-sampled \(h\) to \(\langle v_g,C_fv_g\rangle\), and its total mass
agrees with that of (7). No local positivity of (8) is assumed.

\subsection{3. Full comparison coefficients and the centered
correction}\label{weak-coupling-part-12-full-comparison-coefficients-and-the-centered-correction}

At a fixed box write \(V_i\) for the original counting-unit transverse
edge cochain and \(Y_i=d_1V_i/\sigma\) for its counting-unit curl. For
every edge weight set \[
 A_i^f=\sum_ef_eV_i(e)^2,\qquad
 F_i^f=\sum_pf_pY_i(p)^2,\qquad
 \lambda=\frac{2\sigma}{a},\quad k=\frac32 .
 \tag{9}
\] The fixed-box weighted phase-gradient proof in the compact/radial
companion gives \[
 \langle\psi_g,R_{g,i;f}\psi_g\rangle
                       \longrightarrow k\lambda A_i^f .
 \tag{10}
\] It retains the cutoff derivatives and controls them by the complete
vacuum moment and chart-tail estimates, rather than discarding the
cutoff outside a formal oscillator calculation.

The centered local correction can also be evaluated exactly in that
comparison. Put \(q=\sigma|z_i|^2/4\). Its vacuum density is
\(q^{k-1}e^{-q}/\Gamma(k)\), so \[
 m_0(\beta)=(1-i\beta)^{-k}.
 \tag{11}
\] Only the pair with both modes equal to \(i\) contributes to
\(\langle C_f\Phi_0,e^{i\beta q}\Phi_0\rangle\). The exact pair
coefficient is \(\sigma(F_i^f-A_i^f)/(4a)\), and its creation vector is
\((2q-2k)\Phi_0\). Direct Gamma integration therefore gives \[
 \langle\Phi_0,C_f e^{i\beta q}\Phi_0\rangle
 =\frac{ik\beta\sigma}{2a}(F_i^f-A_i^f)
                                      (1-i\beta)^{-k-1}.
 \tag{12}
\] For example the integral used is
\(2k[(1-i\beta)^{-k-1}-(1-i\beta)^{-k}]\). This fixes the sign and the
factor two. Expanding both vacuum subtractions in (3), using (6),(10),
and taking the real part of (12), proves the full formula \[
 \boxed{\langle v_0(\beta),C_fv_0(\beta)\rangle
   =k\lambda\beta^2 A_i^f+
       \frac{k\sigma\beta^2}{a}
          (1+\beta^2)^{-k-1}(F_i^f-A_i^f).}
 \tag{13}
\] In taking the real part, \(\operatorname{Re}\{i\beta/(1-i\beta)\}
 =-\beta^2/(1+\beta^2)\). The sign of the extra term in (13) is thus
positive times \(F_i^f-A_i^f\). This term would be missed by identifying
the centered and uncentered configuration measures.

For \(f=1\), \(F_i^1=A_i^1=1\), so the correction vanishes. For the
single-edge weights, both \(A_i^{\mathbf1_e}\) and \(F_i^{\mathbf1_e}\)
are nonnegative and their sums over all edges are one. Hence the total
variation of the comparison centered correction is at most \[
 k\lambda\beta^2(1+\beta^2)^{-k-1}.
 \tag{14}
\] This inequality retains every edge and the quarter-face averages.

\subsection{4. A common actual positive-coupling
refinement}\label{weak-coupling-part-12-a-common-actual-positive-coupling-refinement}

Take the original sequence \[
 L_j=j^2,\quad a_j=\frac1{100j},\quad N_j=2j^2+1,\quad
 \ell_j=a_jN_j,\quad
 \sigma_j=\sqrt8\sin\frac{\pi}{2N_j},\quad
 \beta_j=\sqrt j,\quad c_\Gamma=100\sqrt2\pi .
 \tag{15}
\] Then \(j\lambda_j\to c_\Gamma\) and \(\ell_j/j\to1/50\). The original
carrier amplitude, coordinates and physical time have not changed.

At each stage \(j\), add to the finite union of the previous kernel,
weighted-moment and local-energy tests the following tests for every
original edge: the expectations
\(\langle\phi_g(\beta_j),C_{\mathbf1_e}\phi_g(\beta_j)\rangle\) and
\(\langle v_g(\beta_j),C_{\mathbf1_e}v_g(\beta_j)\rangle\) differ from
their complete fixed-box values (6),(10),(13) by less than
\(j^{-10}/|\mathsf E_{L_j}|\). All these are finitely many scalar tests.

For the uncentered term, (6) and the proved weighted gradient
convergence (10) establish convergence at this fixed \(j,\beta_j\). For
the centered correction, the graph convergence
\(C_{\mathbf1_e}\psi_g\to C_{\mathbf1_e,0}\Phi_0\) and the bounded
phase-vacuum convergence prove convergence of the two mixed matrix
elements in its exact expansion. The vacuum multiplier expectation
converges as well. Thus the additional finite list holds for every
sufficiently small positive \(g\).

Choose the least positive dyadic exponent satisfying this union and the
earlier positive-coupling bound. This defines an actual refinement
\(g_j^\sharp\), not a claim that a previously selected least exponent
already satisfies new tests. The total variation distances between
(7),(8) and their respective comparison edge measures are at most
\(j^{-10}\). This gives one sequence for every bounded edge weight at
stage \(j\), with error at most \(j^{-10}\|f\|_\infty\). In particular
it covers both fixed physical profiles and the expanding profiles below
without a separate coupling choice.

\subsection{5. Vanishing energy in fixed compact physical
regions}\label{weak-coupling-part-12-vanishing-energy-in-fixed-compact-physical-regions}

For \(h\in C_c(\mathbb R^3)\), write \(p,q\) for the two spatial
directions other than \(i\), and define \[
 M_{h,i}=\int_{\mathbb R^3}(x_p^2+x_q^2)h(x)\,dx .
 \tag{16}
\] The original lowest cochain components have squares \[
 \frac2{N_j^3}
 \cos^2\frac{\pi a_j(n_p+1/2)}{\ell_j}
 \sin^2\frac{\pi a_j n_q}{\ell_j},\qquad
 \frac2{N_j^3}
 \sin^2\frac{\pi a_j n_p}{\ell_j}
 \cos^2\frac{\pi a_j(n_q+1/2)}{\ell_j}.
 \tag{17}
\] The remaining coordinate is the unit constant vertex factor. The
original midpoint/vertex positions in (17) are retained. On a fixed
compact set, \(\ell_j\sin(\pi x/\ell_j)\to\pi x\) and
\(\cos(\pi x/\ell_j)\to1\), uniformly. Multiplication by the original
spacing volume \(a_j^3\) in the shifted Riemann sums therefore proves \[
 \ell_j^5 A_i^{h,j}\longrightarrow2\pi^2M_{h,i}.
 \tag{18}
\] This argument uses the modulus of continuity of \(h\), an enlarged
compact support for face weights, and a uniform bound on the number of
sampled cells times \(a_j^3\). The companion's exact curl squares also
give \(F_i^{h,j}=O_h(\ell_j^{-3})\).

Combining (6),(10),(15),(18) and the total-variation error gives \[
 \boxed{\ell_j^5\int h\,d\mu_{g_j^\sharp,\sqrt j}
             \longrightarrow300\sqrt2\pi^3M_{h,i}.}
 \tag{19}
\] For the centered correction in (13), \(\sigma_j/a_j=O(\ell_j^{-1})\)
and \(j(1+j)^{-k-1}=O(j^{-3/2})\). Its compact-profile size is therefore
\(O_h(\ell_j^{-4}j^{-3/2})\). Multiplication by \(\ell_j^5\) makes it
tend to zero because \(\ell_j j^{-3/2}\to0\). The error
\(j^{-10}\ell_j^5\) also tends to zero. Thus (19) holds for
\(\widetilde\mu_{g_j^\sharp,\sqrt j}\) as well.

In particular both measures converge vaguely to zero on fixed compact
physical sets. Their total masses instead converge to \[
 kc_\Gamma=150\sqrt2\pi>0 .
 \tag{20}
\] For the positive measures (7), the mass outside every fixed compact
set tends to this total mass. This follows by choosing a compactly
supported continuous majorant of the set and applying (19). It is not a
conclusion drawn merely from a zero low-band projection.

\subsection{6. Exact box-scale coordinate map and retained limiting
density}\label{weak-coupling-part-12-exact-box-scale-coordinate-map-and-retained-limiting-density}

Define the coordinate pushforward \(s_j(x)=x/\ell_j\). It is invertible
on \(\mathbb R^3\), with inverse \(s_j^{-1}(\xi)=\ell_j\xi\). Apply it
to the actual measures; no state or measure is divided by its mass. The
finite support is the original midpoint box divided by \(\ell_j\); its
endpoints tend to \([-1/2,1/2]^3\).

For continuous \(h\) on a neighbourhood of this cube, the same cochain
formula (17) is now an ordinary counting Riemann sum with mesh
\(1/N_j\). It proves \[
 \sum_e h(am(e)/\ell_j)V_i(e)^2
 \longrightarrow\int_{[-1/2,1/2]^3}h(\xi)\rho_i(\xi)\,d\xi,
 \tag{21}
\] where \[
 \boxed{\rho_i(\xi)=
 2\{\cos^2(\pi\xi_p)\sin^2(\pi\xi_q)
       +\sin^2(\pi\xi_p)\cos^2(\pi\xi_q)\}
 =1-\cos(2\pi\xi_p)\cos(2\pi\xi_q).}
 \tag{22}
\] The constant coordinate \(i\) has length one in this cube. Each
sine-square and cosine-square integral on \([-1/2,1/2]\) is \(1/2\), so
\(\int\rho_i=1\). This value follows from the original counting-unit
cochain, not from a new normalization of the energy measure.

The total-variation bound (14) tends to zero because \(j\lambda_j\)
stays bounded and \((1+j)^{-5/2}\to0\). The actual edge-test errors
vanish as well. Equations (10),(15),(21) therefore prove the two weak
finite-measure limits \[
 \boxed{(s_j)_*\mu_{g_j^\sharp,\sqrt j},\
        (s_j)_*\widetilde\mu_{g_j^\sharp,\sqrt j}
 \ \Longrightarrow\
 150\sqrt2\pi\,\rho_i(\xi)
                    \mathbf1_{[-1/2,1/2]^3}\,d^3\xi .}
 \tag{23}
\] For the signed centered measures weak convergence follows from their
vanishing total-variation distance to the positive comparison measure.
Equations (19),(20),(23) specify both spatial scales exactly. The
limiting Gamma energy is not concentrated at a point of the original
physical space; its nonzero density is visible on the growing box scale.

\subsection{7. Consequence for the local spectral
programme}\label{weak-coupling-part-12-consequence-for-the-local-spectral-programme}

The spatial measures (7),(8) encode local-energy expectations, that is,
spatially resolved first moments. They are not a joint positive measure
of position and spectral energy, and (23) does not identify such a
measure with the Gamma spectral law. Local energies and the full
Hamiltonian need not commute. Both distributions have been calculated at
their stated levels.

The full compact local-energy spectrum proved in the companion uses a
growing family of transverse modes and retains nonzero positive-time
norms. Equations (19)--(23) instead evaluate the energy of one specified
global radial carrier with its increasing occupation. They are two
computations within the same original local energy and phase-gradient
identities, not claims that the two sectors are unrelated.

This spatial measure calculation rules out interpreting the carrier's
finite total energy as a nonzero limiting energy localized in a fixed
compact region, on this specified sequence. It does not rule out other
states, mixed local excitations, or another retained-coupling continuum
trajectory. The full local products and the nonlinear coupling-response
integral are the remaining concrete routes to test that larger
construction.

The proof inputs are the included \emph{Compact local energy, radial
electric observables, and a full-mode low-energy lower bound} (weighted
phase-gradient convergence and original cochain squares), \emph{Spatial
continuum}, Sections 22--23 (actual radial moment tails and common-mode
maps), and the complete full-mode local-spectrum companion. All formulas
above concern the original finite Hamiltonian and its proved fixed-box
maps; no outside fluid or Yang--Mills existence theorem is assumed here.

\clearpage
\part{Local spectrum and interaction continuation checkpoint}
\section{Yang--Mills feedback: full local spectrum completed; retain the
interaction
next}\label{weak-coupling-part-13-yangmills-feedback-full-local-spectrum-completed-retain-the-interaction-next}

9 September continuation: the full first coupling coefficients, cubic
singlet inverse, vacuum correction and six-dimensional first
physical-gap matrix are now explicitly calculated. The new bounded task
is WEB\_PRO\_FIRST\_PHYSICAL\_GAP\_MATRIX\_FEEDBACK\_20260909.md,
together with cubic\_singlet\_vacuum\_energy.md. The earlier request
below is retained as the preceding checkpoint, not an instruction to
repeat completed work.

8 September 2026. This updates the earlier request to calculate the full
local-energy two-mode spectrum: that calculation is now complete. Do not
repeat only the oscillator spectrum or relabel it as an interacting
mass-gap counterexample.

\subsection{Exact retained objects}\label{weak-coupling-part-13-exact-retained-objects}

\[
 H_g=\frac{2g^2}{a}\sum_eE_e+\frac1{2g^2a}
                  \sum_p(2-\operatorname{tr}U_p),\qquad
 L_j=j^2,\quad a_j=\frac1{100j},\quad\ell_j=a_j(2j^2+1).
\] All original edges, faces, SU(2) colour coordinates, gauge
projection, Haar measure, actual positive vacuum \(\psi_g\) and vacuum
energy \(\mathcal E_g\) are retained. Set \(A_g=H_g-\mathcal E_g\). For
\(h\in C_c^\infty(\mathbb R^3)\) put \[
 h_e=h(am(e)),\quad h_p=\tfrac14\sum_{e\in\partial p}h_e,\quad
 D_{h,g}=\frac{2g^2}{a}\sum_eh_eE_e+
     \frac1{2g^2a}\sum_ph_p(2-\operatorname{tr}U_p),\quad
 \Xi_{h,g}=(D_{h,g}-\langle D_{h,g}\rangle_{\psi_g})\psi_g .
\]

\subsection{Completed all-mode
calculation}\label{weak-coupling-part-13-completed-all-mode-calculation}

With \(\widehat h(k)=\int e^{-ik\cdot x}h(x)\,dx\), the complete open
standing-wave, parity/sign and continuum calculation gives \[
 \rho_h(\omega)=\frac1{320\pi^5}
          \int_{|k|\le\omega}|k|^4|\widehat h(k)|^2\,dk,\qquad
 C_h(t)=\frac1{320\pi^5t}
          \int |k|^4e^{-t|k|}|\widehat h(k)|^2\,dk.
\] This is the covariance of the unprojected local-energy state at
positive Euclidean time. Both electric and magnetic pair terms are
included; opposite momenta cancel, exactly as \(D_1=H\) requires. For
\(H_h=\int h\ne0\), \[
 \rho_h(\omega)\sim\frac{H_h^2}{560\pi^4}\omega^7,\quad
 C_h(t)\sim\frac{9H_h^2}{\pi^4t^8},\quad
 -C_h'(t)\sim\frac{72H_h^2}{\pi^4t^9}.
\] The raw spectral measure has infinite total mass, but
\(e^{-t\omega}\rho_h(\omega)d\omega\) is finite with every moment for
each \(t>0\). No vector is divided by its norm.

All finite-edge covariance tests can be imposed on one positive dyadic
\(g_j\), using the proved fixed-box nonlinear graph convergence. They
give the same actual sequence for every compact smooth \(h\) and every
positive time. Its limiting state map is explicitly the equal-colour
two-creation map in transverse Fock space with Hamiltonian
\(d\Gamma(|p|)\). Thus this is a proved free spectral correspondence
from the specified nonlinear regulators, not an identified interacting
continuum theory.

The position kernel is now evaluated: \[
 K_t(x)=\frac{9t^4-30t^2|x|^2+9|x|^4}
                    {\pi^4(t^2+|x|^2)^6}.
\] It matches the 00,00 free-vector stress covariance with
\(C_T=12/\pi^4\), in Osborn--Petkou (1994), Eqs. (2.23),(5.16). The
joint measure is
\(|k|^4\mathbf1_{\{\omega\ge|k|\}}d\omega\,d^3k/(320\pi^5)\). Its
reflected two-point form is positive and its real-time commutator
expectation is supported on the light cone. These assertions do not
replace missing higher local products.

\subsection{What the global Gamma carrier does
spatially}\label{weak-coupling-part-13-what-the-global-gamma-carrier-does-spatially}

Keep its original amplitude \(\beta_j=\sqrt j\), mode phase \(Q_{g,i}\),
and centered state
\(v_j=(e^{i\sqrt jQ_{g,i}}-\langle e^{i\sqrt jQ_{g,i}}\rangle)\psi_g\).
Its total excitation mean tends to \(150\sqrt2\pi\). The spatial measure
of its local-energy expectations tends vaguely to zero on fixed compact
physical sets. Under the exact coordinate pushforward \(\xi=x/\ell_j\),
its measure instead tends to \[
 150\sqrt2\pi
 [1-\cos(2\pi\xi_p)\cos(2\pi\xi_q)]
                 \mathbf1_{[-1/2,1/2]^3}\,d^3\xi,\quad p,q\ne i.
\] This holds for the actual centered state on a further specified
finite-edge-test refinement. It is a spatial first-moment statement, not
a joint position/spectral measure. The full local spectrum above uses
the growing family of modes, not just this delocalized carrier.

\subsection{Exact finite-coupling response now
available}\label{weak-coupling-part-13-exact-finite-coupling-response-now-available}

At fixed \(L,a,g>0\), abbreviate
\(K=\sum E_e,\ W=\sum(2-\operatorname{tr}U_p)\). Keep \[
 V_g=H_g'=\frac{4g}{a}K-\frac1{g^3a}W,\quad
 Z_{h,g}=D_{h,g}'=\frac{4g}{a}K_h-\frac1{g^3a}W_h,\quad
 v_g=\langle V_g\rangle_{\psi_g}.
\] With \(R_g=A_g^{-1}\) on the vacuum complement and zero on the
vacuum, put \(\eta_g=R_g(V_g-v_g)\psi_g\). Then \[
 \psi_g'=-\eta_g,\quad
 m_{h,g}'=\langle Z_{h,g}\rangle_{\psi_g}
                   -2\operatorname{Re}\langle\eta_g,D_{h,g}\psi_g\rangle,
\] \[
 \Xi_{h,g}'=(Z_{h,g}-m_{h,g}')\psi_g
                          -(D_{h,g}-m_{h,g})\eta_g .
\] The full response is \[
 \partial_g C_{g;h}(t)=
 2\operatorname{Re}\langle\Xi_{h,g}',e^{-tA_g}\Xi_{h,g}\rangle
 -\int_0^t\langle\Xi_{h,g},e^{-(t-s)A_g}
                    (V_g-v_g)e^{-sA_g}\Xi_{h,g}\rangle\,ds .
\] All domains and the unbounded Duhamel insertion are proved at fixed
regulator. The constant weight \(h_e=1\) makes both the state and its
derivative exactly zero, checking the vacuum-motion terms. The actual
difference from the full fixed-box comparison is the improper integral
of this derivative from zero to \(g\).

\subsection{One bounded demanding
task}\label{weak-coupling-part-13-one-bounded-demanding-task}

Calculate the first interaction-sensitive contribution to this complete
local covariance response, including the reduced-resolvent vacuum terms
and both Wilson/electric derivatives. Work from the full retained graph
operator and its native curl/mode tensors; do not replace it by a
homogeneous truncation. Give an explicit finite-mode sum or kernel and
its dependence on \(a,L,g,t\) for a fixed nonzero compact smooth
profile, then determine a regime where the remainder is actually
bounded.

The useful output is an evaluated correction with its scale dependence,
not the statement that an unknown correction should be small. If a
nonuniform term appears, retain it and determine its spatial/momentum
map rather than selecting \(g\) so small that the term disappears. Do
not infer a full interacting construction or its gap from one
perturbative coefficient. The currently proved identities supply the
exact relation to be continued.

Complete proof sources in the canonical AI bundle:
FULL\_LOCAL\_SPECTRUM.md;
local\_energy\_spacetime\_and\_coupling\_response.md;
radial\_carrier\_spatial\_energy\_measure.md; the original finite-box,
graph, state, and complete quantum-blocking companions.

\clearpage
\part{Full original operator and weighted coupling coefficients}
\section{The first two coupling coefficients of the full tree-coordinate
Hamiltonian}\label{weak-coupling-part-14-the-first-two-coupling-coefficients-of-the-full-tree-coordinate-hamiltonian}

8 September 2026.

This note computes the terms of orders g and g squared after the exact
maximal-tree, Haar-density, and dilation map of the original finite-box
SU(2) Hamiltonian. Every original edge and face remains in the formulas.
The coefficients are finite differential expressions in all chord
coordinates. No spatial mode, Wilson word, density term, or scalar
energy has been removed. A quantitative remainder estimate is proved on
compact chart cores. This is a fixed-box operator calculation; it does
not assert a uniform-volume perturbation expansion or an interacting
continuum construction.

\subsection{1. Original operator, weights, and graph
data}\label{weak-coupling-part-14-original-operator-weights-and-graph-data}

Fix an integer L\textgreater=2 and physical spacing a\textgreater0. The
vertices are \(\{-L,\ldots,L\}^3\). Let E be every positively oriented
contained nearest-neighbor edge and P every contained positively
oriented elementary face. Keep

\[
 T_\alpha=-\frac i2\sigma_\alpha,\qquad
 [T_\alpha,T_\beta]=\epsilon_{\alpha\beta\gamma}T_\gamma,
 \qquad E_e=-\sum_{\alpha=1}^3X_{e,\alpha}^2,
\]

and the full operator

\[
 H_g=\frac{2g^2}{a}\sum_{e\in E}E_e
       +\frac1{2g^2a}\sum_{p\in P}(2-\operatorname{tr}U_p),
 \qquad g>0.                                                   \tag{1}
\]

The measure is product Haar probability, and the physical subspace is
the invariant subspace for all vertex gauge transformations. In
particular the coefficients remain \(\kappa=2g^2/a\), \(b=1/(2g^2a)\),
and \(\xi=1/(4g^4)\), with the full scalar Wilson contribution \(2b|P|\)
included in (1).

For fixed real edge weights h\_e retain the original face average and
weighted operator

\[
 h_p=\frac14\sum_{e\in\partial p}h_e,\qquad
 D_{h,g}=\frac{2g^2}{a}\sum_eh_eE_e
       +\frac1{2g^2a}\sum_ph_p(2-\operatorname{tr}U_p).          \tag{2}
\]

For a physical profile h one may take h\_e=h(am(e)), with the original
edge midpoint m(e). The weights and a are fixed throughout this coupling
expansion; their derivatives with respect to g are zero. The
calculations allow signed weights. All operator identities below are
first stated on smooth compact chart functions, which lie in the
original weighted operators' domains.

Use the fixed maximal tree rooted at o=(-L,-L,-L), with parent obtained
by decreasing the first coordinate in the order 1,2,3 that exceeds -L.
Let \(\mathcal T\) be its tree edges and \(\mathcal C\) its chords, with
\(r=|\mathcal C|=2(2L)^3+3(2L)^2\). If t\_v(U) is the ordered root-to-v
tree holonomy, the exact chord is

\[
 Z_c=t_{s(c)}U_ct_{t(c)}^{-1}.                              \tag{3}
\]

The inverse retains the tree links and sets U\_c=t\_s\^{}\{-1\}Z\_ct\_t.
Independent left and right Haar translations in each chord preserve the
full product measure. Based gauge averaging removes the tree variables;
residual physical invariance is simultaneous conjugation of every chord,
not separate conjugation in each chord.

For every original edge e and chord c define two signed coefficients
\(\ell_{ec},r_{ec}\) by

\[
 (\ell_{ec},r_{ec})=
 \begin{cases}
 (\mathbf1_{\{c=e\}},0),&e\in\mathcal C,\\
 (\epsilon(s(c),e),-\epsilon(t(c),e)),&e\in\mathcal T,
 \end{cases}                                                \tag{4}
\]

where epsilon(v,e) is one when the original tree edge e lies on the
specified root-to-v path and zero otherwise. Put

\[
 t_{ec}=\ell_{ec}+r_{ec},\qquad
 s_{ec}=\ell_{ec}-r_{ec}.                                  \tag{5}
\]

Thus a tree edge has t\_ec=epsilon(s(c),e)-epsilon(t(c),e) and
s\_ec=epsilon(s(c),e)+epsilon(t(c),e). A chord edge has t\_ec=s\_ec
=delta\_ec. These signs are fixed by the inverse terminal holonomy in
(3). The additive tree map has matrix T\_\{ce\}=t\_ec; consequently

\[
 G_{cd}=\sum_e t_{ec}t_{ed},\qquad
 (G_h)_{cd}=\sum_eh_e t_{ec}t_{ed}.                         \tag{6}
\]

The original oriented face curl composed with zero-tree chord insertion
is denoted C, as in the fixed-box companion. No metric on the chords has
been substituted for G.

\subsection{2. Exact left and right fields, with their logarithmic
signs}\label{weak-coupling-part-14-exact-left-and-right-fields-with-their-logarithmic-signs}

Write \(Z_c=\exp(y_c^\alpha T_\alpha)\), \(|y_c|<2\pi\). For one chord
use \(A_yv=y\times v\), which is exactly the coefficient vector of
\([y^\alpha T_\alpha,v^\beta T_\beta]\). Left and right multiplication
give coordinate vector fields

\[
 L_\alpha=\left(I-\frac12A_y+b(|y|)A_y^2\right)e_\alpha
                      \mathbin{\cdot}\nabla_y,
 \qquad
 R_\alpha=\left(I+\frac12A_y+b(|y|)A_y^2\right)e_\alpha
                      \mathbin{\cdot}\nabla_y,               \tag{7}
\]

where e\_alpha is the standard color basis vector and

\[
 b(u)=\frac1{u^2}-\frac{\cot(u/2)}{2u}\quad(0<u<2\pi),
 \qquad b(0)=\frac1{12}.                                   \tag{8}
\]

Here b(u) is a scalar analytic function, distinct from the original
Hamiltonian coefficient b=1/(2g\^{}2a).

To prove (7) and its signs, differentiate the exponential:

\[
 (d\exp)_y(v)\exp(-y)=\int_0^1e^{sA_y}v\,ds.
\]

For exp(t e\_alpha) exp(y), its left-translated derivative is e\_alpha;
the logarithmic derivative is therefore \(A_y(e^{A_y}-I)^{-1}e_\alpha\),
with the analytic value on the kernel of A\_y. On the line spanned by y
this is the identity. On the perpendicular plane
A\_y\textsuperscript{2=-\textbar y\textbar{}}2 I, and direct inversion
of

\[
 \frac{\sin |y|}{|y|}I+
           \frac{1-\cos |y|}{|y|^2}A_y
\]

gives \((|y|/2)\cot(|y|/2)I-A_y/2\). This is exactly the first
expression in (7). For exp(y) exp(t e\_alpha), the same left-translated
derivative is exp(A\_y)e\_alpha; multiplying by that factor changes the
sign of A\_y/2 and gives the right field. Thus the left sign is negative
and the right sign positive in the original T\_alpha convention.

Taylor division of cosine by sine gives

\[
 b(u)=\frac1{12}+\frac{u^2}{720}+O(u^4).                    \tag{9}
\]

In particular the first logarithmic coefficients are I minus or plus
A\_y/2 plus A\_y\^{}2/12, and there is no cubic coefficient in an
individual left or right vector field.

The exact quotient fields of the original kinetic operator are

\[
 \mathcal X_{e,\alpha}
 =\sum_c\{\ell_{ec}L_{c,\alpha}+r_{ec}R_{c,\alpha}\}.
                                                                  \tag{10}
\]

Each displayed field is divergence-free for product Haar measure. The
exact quotient kinetic operator is
\(-(2g^2/a)\sum_{e,\alpha}\mathcal X_{e,\alpha}^2\); the weighted
version inserts h\_e in the same original-edge sum. This follows edge by
edge from the invariance of the sum over the three generators under the
adjoint rotations induced by changing the tree representative. It does
not require the weights to be constant across edges.

\subsection{3. The exact Haar/dilation map and its density
contribution}\label{weak-coupling-part-14-the-exact-haardilation-map-and-its-density-contribution}

Retain the complete density and isometry

\[
 \mathcal J(y)=(16\pi^2)^{-r}
   \prod_c\left[\frac{\sin(|y_c|/2)}{|y_c|/2}\right]^2,
 \qquad
 (\mathcal B_gF)(x)=g^{3r/2}\mathcal J(gx)^{1/2}F(\exp(gx))
                                                               \tag{11}
\]

on \(\Omega_g=\{x:|x_c|<2\pi/g\}\), with zero extension outside. Its
adjoint includes the reciprocal density. Compactly supported x-functions
whose support lies inside Omega\_g pull back to smooth functions in the
original operator domain. The factor (16pi\textsuperscript{2)}-r and the
dilation power in (11) remain in this exact map.

Put \(l_g(x)=\log\mathcal J(gx)\). Its first terms are

\[
 l_g(x)=-r\log(16\pi^2)
   -\frac{g^2}{12}\sum_c|x_c|^2
   -\frac{g^4}{1440}\sum_c|x_c|^4+O(g^6)
                                                               \tag{12}
\]

on fixed compact x-sets. The constants follow by expanding 2
log(sin(u/2)/(u/2)); the remainder and its derivatives are bounded on
each smaller chart.

For each e,alpha define the exact first-order transformed field
\(\mathcal Q_{e,\alpha}(g)=\mathcal B_g(g\mathcal X_{e,\alpha})
\mathcal B_g^*\) on these compact cores. The chain and product rules
give the vector part from (7), followed by subtraction of one half of
its derivative of l\_g. Because \(x\cdot(x\times e_\alpha)=0\) and
\(x\cdot[x\times(x\times e_\alpha)]=0\), the exact result is

\[
 \boxed{\begin{split}
 \mathcal Q_{e,\alpha}(g)
 ={}&\sum_c t_{ec}\partial_{c,\alpha}
   -\frac g2\sum_c s_{ec}(x_c\times e_\alpha)\cdot\nabla_c\\
 &+g^2\sum_c t_{ec}b(g|x_c|)
       \left([x_c\times(x_c\times e_\alpha)]\cdot\nabla_c
                           +x_c^\alpha\right).
 \end{split}}                                                \tag{13}
\]

The last x\_c\^{}alpha is multiplication, with a positive sign. In fact
the radial derivative of l\_g is
\([g\cot(g|x_c|/2)/|x_c|-2/|x_c|^2]x_c\); multiplication by -1/2 and the
radial part t\_ec e\_alpha gives exactly g\^{}2 t\_ec
b(g\textbar x\_c\textbar)x\_c\^{}alpha. Thus (13) includes the full
density term, not only its leading approximation.

Write its first three evaluated differential coefficients as

\[
 \begin{split}
 Q^{(0)}_{e\alpha}&=\sum_ct_{ec}\partial_{c,\alpha},\\
 Q^{(1)}_{e\alpha}&=-\frac12\sum_cs_{ec}
                         (x_c\times e_\alpha)\cdot\nabla_c,\\
 Q^{(2)}_{e\alpha}&=\frac1{12}\sum_ct_{ec}
       \left((x_c x_c^\alpha-|x_c|^2e_\alpha)\cdot\nabla_c
                           +x_c^\alpha\right).
 \end{split}                                                \tag{14}
\]

Then
Q(g)=Q\textsuperscript{(0)+gQ}(1)+g\textsuperscript{2Q}(2)+O(g\^{}4) in
the coefficient C1 norm on compact sets. Each coefficient is formally
skew-adjoint in dx: Q\^{}(0) and Q\^{}(1) have divergence zero; the
vector part of Q\^{}(2) has divergence \(\sum_ct_{ec}x_c^\alpha/6\), and
its multiplication term is one half of that divergence. This
independently checks the sign of the half-density term.

\subsection{4. Evaluated kinetic tensors and the retained
scalar}\label{weak-coupling-part-14-evaluated-kinetic-tensors-and-the-retained-scalar}

Let composite indices be (c,beta), with beta=1,2,3, and sum repeated
color components explicitly or by their stated finite ranges. The exact
principal tensor is the sum of the outer products of the vector
coefficients of (13), excluding its multiplication term. Its first
coefficients are

\[
 A_0^{c\beta,d\gamma}=G_{cd}\delta_{\beta\gamma},           \tag{15}
\]

\[
 \boxed{A_1^{c\beta,d\gamma}
 =\frac12\sum_{e,\eta}\epsilon_{\eta\beta\gamma}
       (s_{ec}t_{ed}x_c^\eta-t_{ec}s_{ed}x_d^\eta),}          \tag{16}
\]

and

\[
 \boxed{\begin{split}
 A_2^{c\beta,d\gamma}=\sum_e\Big\{
 &\frac{s_{ec}s_{ed}}4
       [\delta_{\beta\gamma}(x_c\cdot x_d)
                                      -x_d^\beta x_c^\gamma]\\
 &+\frac{t_{ec}t_{ed}}{12}
       [x_c^\beta x_c^\gamma+x_d^\beta x_d^\gamma
          -( |x_c|^2+|x_d|^2)\delta_{\beta\gamma}]
                       \Big\} .
 \end{split}}                                                \tag{17}
\]

Every original edge appears, with its exact t and s coefficients. For
weighted operators define A\_\{n,h\} by inserting h\_e into each edge
sum in (15)--(17); A\_\{0,h\}=G\_h tensor I\_3. The tensors are real and
symmetric in the two composite indices.

For a direct derivation, write the vector coefficients of (14) as
a\_0,a\_1,a\_2. The order-one tensor is sum(a\_0 a\_1\^{}T+a\_1
a\_0\^{}T), giving (16). The order-two tensor is sum(a\_1 a\_1\^{}T+a\_0
a\_2\^{}T+a\_2 a\_0\^{}T). The color contraction

\[
 \sum_\alpha (x_c\times e_\alpha)^\beta
                 (x_d\times e_\alpha)^\gamma
 =\delta_{\beta\gamma}(x_c\cdot x_d)-x_d^\beta x_c^\gamma
\]

gives the first line of (17); inserting x\_c times x\_c\^{}alpha minus
\textbar x\_c\textbar\^{}2 e\_alpha gives its second line.

The exact divergence-form kinetic operator after (11) is

\[
 \frac2a\left[-\partial_i A^{ij}(gx)\partial_j
   +\frac12\partial_i(A^{ij}(gx)\partial_j l_g)
   +\frac14(\partial_i l_g)A^{ij}(gx)(\partial_j l_g)\right].
                                                               \tag{18}
\]

In the first term the derivatives act on the function to their right as
well as on A. The next two terms are multiplication operators. Formula
(18) follows by conjugating -J\^{}-1 partial\_i(J A\^{}\{ij\}
partial\_j) with J\^{}1/2 and expanding both derivatives. It is
equivalently -(2/a) sum Q\_ealpha(g)\^{}2.

Since partial\_j l\_g=-g\^{}2 x\_j/6+O(g\^{}4), the order-two density
scalar in the bracket is -g\^{}2 trace(A\_0)/12=-g\^{}2 trace(G)/4.
Therefore the kinetic coefficients are exactly

\[
 \begin{split}
 K_0&=-\frac2a\partial_i A_0^{ij}\partial_j,\\
 K_1&=-\frac2a\partial_i A_1^{ij}\partial_j,\\
 K_2&=-\frac2a\partial_i A_2^{ij}\partial_j
                       -\frac{\operatorname{tr}G}{2a}I.
 \end{split}                                                \tag{19}
\]

The corresponding weighted scalar is -trace(G\_h)/(2a). It has not been
absorbed into a changed vacuum energy. Notice also that partial\_i
A\_1\^{}\{ij\}=0: for each e,alpha, its constant vector field
differentiates the linear field only in the same alpha direction, and
partial\_alpha(x times e\_alpha)=0, while the linear field has
divergence zero. Thus K\_1 has no hidden first-order drift. The
divergence notation in K\_2 retains all the quadratic metric's
first-order derivatives.

As a check for one chord and one left Casimir, t=s=1 gives A\_1=0 and
A\_2=(\textbar x\textbar\^{}2 I-xx\^{}T)/12. Its order-two kinetic
coefficient is \(-\Delta_{S^2}/(6a)-1/(2a)\), using
\(\Delta_{S^2}=|x|^2\Delta-(x\cdot\nabla)^2-x\cdot\nabla\). This follows
from (17)--(19) without changing the group radius or the original factor
2/a.

\subsection{5. All ordered Wilson coefficients through degree
four}\label{weak-coupling-part-14-all-ordered-wilson-coefficients-through-degree-four}

For a face p=(n;i,j), i\textless j, retain its original four ordered
slots

\[
 ((n,i),+),\quad ((n+e_i,j),+),\quad
 ((n+e_j,i),-),\quad ((n,j),-).                            \tag{20}
\]

In slot m put v\_\{p,m\}=epsilon\_m x\_c if its edge is chord c, and
v\_\{p,m\}=0 if it is a tree edge. Thus tree slots remain in the ordered
list with their exact identity factor. The face word in the tree-I
representative is exactly

\[
 U_p(gx)=\prod_{m=1}^4\exp(g v_{p,m}^\alpha T_\alpha),
                                                               \tag{21}
\]

in the order (20). In particular inverse traversal is represented by the
negative vector in that same slot; no factors are commuted or reordered.
For this section fix p and write v\_m for its four vectors,
r\_m\^{}2=v\_m dot v\_m, and d\_mn=v\_m dot v\_n.

The Pauli multiplication rule in the retained convention gives

\[
 \begin{split}
 \operatorname{tr}(vT\,wT)&=-\tfrac12 v\cdot w,\\
 \operatorname{tr}(vT\,wT\,zT)&=-\tfrac14 v\cdot(w\times z),\\
 \operatorname{tr}(vT\,wT\,zT\,uT)
 &=\tfrac18[(v\cdot w)(z\cdot u)-(v\cdot z)(w\cdot u)
                                      +(v\cdot u)(w\cdot z)],
 \end{split}                                                \tag{22}
\]

where vT=v\^{}alpha T\_alpha. These identities follow by applying
T\_alpha
T\_beta=-(delta\_alpha\_beta/4)I+(epsilon\_alpha\_beta\_gamma/2)
T\_gamma and taking the trace. In particular
(vT)\textsuperscript{2=-\textbar v\textbar{}}2I/4.

Define w\_\{n,p\} by expanding the exact expression 2-tr U\_p(gx). Its
constant and linear terms vanish, and the next coefficients are

\[
 \boxed{w_{2,p}=\frac14\left|\sum_{m=1}^4v_m\right|^2,
 \qquad
 w_{3,p}=\frac14\sum_{m<n<l}v_m\cdot(v_n\times v_l).}       \tag{23}
\]

The complete quartic coefficient is

\[
 \boxed{\begin{split}
 w_{4,p}=-\Bigg[\,
 &\sum_m\frac{r_m^4}{192}
 +\sum_{m<n}\left\{\frac{(r_m^2+r_n^2)d_{mn}}{48}
                                  +\frac{r_m^2r_n^2}{32}\right\}\\
 &+\sum_{m<n<l}
       \frac{r_m^2d_{nl}+r_n^2d_{ml}+r_l^2d_{mn}}{16}\\
 &+\sum_{m<n<l<q}
       \frac{d_{mn}d_{lq}-d_{ml}d_{nq}+d_{mq}d_{nl}}8
                                                        \Bigg].
 \end{split}}                                                \tag{24}
\]

All sums use the retained slot order. There is at most one
four-distinct-slot term, but it is left in the displayed ordered form so
every position and inverse sign is explicit.

To verify completeness, the degree-n coefficient of the product in (21)
is the finite sum

\[
 \sum_{a_1+\cdots+a_4=n}
 \frac{(v_1T)^{a_1}(v_2T)^{a_2}
       (v_3T)^{a_3}(v_4T)^{a_4}}{a_1!a_2!a_3!a_4!}.
\]

At degree three, the patterns (3) and (2,1) have zero trace; only the
three-distinct-slot trace in (22) survives, with the positive sign in
w\_3 because W=2-tr U. At degree four the patterns (4), (3,1), (2,2),
(2,1,1), and (1,1,1,1) give respectively the five displayed types in
(24), with the negative sign from W. This proves that none of the
quartic word contributions is missing.

For a single nonzero slot, (23)--(24) give r\^{}2/4 and -r\^{}4/192, the
direct expansion of 2-2 cos(g r/2). For the four-slot test (v,w,-v,-w),
(23) has zero quadratic and cubic terms and (24) gives
\(|v\times w|^2/4\). This agrees with the commutator word's leading
logarithm g\^{}2{[}vT,wT{]}. This check does not replace the independent
original face letters in (20). In particular w\_4 need not have one sign
for arbitrary noncommuting letters.

Set

\[
 W_n(x)=\sum_{p\in P}w_{n,p}(x),\qquad
 W_{n,h}(x)=\sum_{p\in P}h_p w_{n,p}(x),\qquad n=2,3,4.
                                                               \tag{25}
\]

The signed sum in w\_2 is the original curl (Cx)\_p, so W\_2=one quarter
sum\_alpha
\textbar\textbar Cx\textsuperscript{alpha\textbar\textbar{}}2. The
expansion of the full Wilson multiplication term is therefore

\[
 \frac{W(gx)}{2g^2a}
 =\frac1{2a}\{W_2+gW_3+g^2W_4\}+O(g^3),                 \tag{26}
\]

with the weighted version using W\_\{n,h\}. This is an expansion of the
full face terms 2-tr U\_p: their constant at the identity cancels inside
each original expression. It is not deletion of the original scalar
2b\textbar P\textbar{} from the operator.

\subsection{6. The full coefficients and the weighted
observable}\label{weak-coupling-part-14-the-full-coefficients-and-the-weighted-observable}

On compact chart cores define \(\widetilde H_g=\mathcal B_gH_g
\mathcal B_g^*\), including the exact tree quotient already specified.
The evaluated coefficients are

\[
 \boxed{\begin{split}
 H_0&=-\frac2a\partial_i A_0^{ij}\partial_j+\frac{W_2}{2a},\\
 H_1&=-\frac2a\partial_i A_1^{ij}\partial_j+\frac{W_3}{2a},\\
 H_2&=-\frac2a\partial_i A_2^{ij}\partial_j
                     -\frac{\operatorname{tr}G}{2a}I
                     +\frac{W_4}{2a}.
 \end{split}}                                                \tag{27}
\]

Here A\_0,A\_1,A\_2 are exactly (15)--(17), and W\_2,W\_3,W\_4 are the
explicit four-letter polynomials (23)--(25). Thus

\[
 \widetilde H_g f=H_0f+gH_1f+g^2H_2f+\mathcal R_g f.
                                                               \tag{28}
\]

The leading coefficient is precisely the complete companion oscillator
in the original chord coordinates:

\[
 H_0=\frac1a\left[-2\sum_{\alpha,c,d}G_{cd}
          \partial_{c,\alpha}\partial_{d,\alpha}
                 +\frac18\sum_\alpha\|Cx^\alpha\|^2\right].
\]

No diagonalization or selection of oscillator modes was used to obtain
either higher coefficient.

Likewise \(\widetilde D_{h,g}=\mathcal B_gD_{h,g}\mathcal B_g^*\) has
coefficients

\[
 \boxed{\begin{split}
 D_{h,0}&=-\frac2a\partial_i A_{0,h}^{ij}\partial_j
                                    +\frac{W_{2,h}}{2a},\\
 D_{h,1}&=-\frac2a\partial_i A_{1,h}^{ij}\partial_j
                                    +\frac{W_{3,h}}{2a},\\
 D_{h,2}&=-\frac2a\partial_i A_{2,h}^{ij}\partial_j
                     -\frac{\operatorname{tr}G_h}{2a}I
                                    +\frac{W_{4,h}}{2a}.
 \end{split}}                                                \tag{29}
\]

The constant profile gives D\_\{1,n\}=H\_n for each n=0,1,2, as required
by the original identity D\_\{1,g\}=H\_g. These are operator
coefficients, before subtraction of the moving actual vacuum energy or
weighted vacuum expectation. Centering and vacuum-response derivatives
remain the exact operations in the coupling-response companion; they are
not silently incorporated into H\_2 or D\_\{h,2\}.

\subsection{7. A quantitative compact-core remainder
estimate}\label{weak-coupling-part-14-a-quantitative-compact-core-remainder-estimate}

This section makes the remainder in (28) an actual norm bound. Put N=3r
and take R\textgreater=1. Assume f is smooth and compactly supported in
the Euclidean ball \textbar x\textbar\textless R, and
\(0<g\le\min(1,R^{-1})\). Then every
g\textbar x\_c\textbar\textless=1\textless2pi, so the support is
strictly inside the exact logarithm chart. Define

\[
 \|f\|_{2,\Sigma}=\sum_{|\alpha|\le2}\|\partial^\alpha f\|_2,
 \quad C_N=(N+1)(2N+1),
\]

where the multi-index alpha here ranges over all N real coordinates. For
the explicit scalar in (8) put

\[
 d(u)=\frac{b(u)-1/12}{u^2},\quad d(0)=\frac1{720},\qquad
 B_* =1+\max_{0\le u\le1}(|d(u)|+|d'(u)|).                  \tag{30}
\]

The even Taylor expansion (9) makes d smooth at zero, with d'(0)=0; its
elementary trigonometric expression is smooth on the rest of this
interval. Thus B\_* is a specified finite positive real constant,
independent of the graph or regulator. It can be evaluated from (8); it
is not an unknown operator remainder or a mathematical hypothesis.

For each edge put

\[
 T_e=\sum_c|t_{ec}|,\quad S_e=\sum_c|s_{ec}|,
 \quad M_{0,e}=T_e,\quad
 M_{1,e}=\frac{S_e(1+R)}2,\quad
 M_{2,e}=\frac{T_e(1+R)^2}{12},\quad
 M_{4,e}=10T_eB_*(1+R)^4.
                                                               \tag{31}
\]

Let q\_p be the number of chord slots among the four original slots of
face p.~For the weighted expansion in (29) a valid explicit remainder
constant is

\[
 \begin{split}
 \mathcal C_{h,L,a,R}={}&\frac{6C_N}{a}\sum_e|h_e|
 \left[2M_{1,e}M_{2,e}+M_{2,e}^2
  +2(M_{0,e}+M_{1,e}+M_{2,e})M_{4,e}+M_{4,e}^2\right]\\
 &+\frac1{a\,5!}\sum_p|h_p|\left(\frac{q_pR}{2}\right)^5.
                                                               \tag{32}
\end{split}
\]

Every quantity in this finite sum is now defined by the original graph,
weights, a, R, and the explicit scalar function (8). The estimate is

\[
 \boxed{
 \|[\widetilde D_{h,g}-D_{h,0}-gD_{h,1}-g^2D_{h,2}]f\|_2
       \le g^3\mathcal C_{h,L,a,R}\|f\|_{2,\Sigma}.}
                                                               \tag{33}
\]

Taking every h\_e=h\_p=1 proves the same estimate for (28).

Here is the proof. For a first-order operator P, let M(P) be the largest
absolute value of a coefficient or one of its first coordinate
derivatives on \textbar x\textbar\textless=R, including its
multiplication coefficient. Expanding the composition of two such
operators produces N\^{}2 second-order terms, N\^{}2+2N first-order
terms, and N+1 multiplication terms. Thus

\[
 \|PQf\|_2\le C_N M(P)M(Q)\|f\|_{2,\Sigma}.              \tag{34}
\]

The explicit polynomials in (14) obey
M(Q\^{}(i)\_ealpha)\textless=M\_i,e for i=0,1,2. For example
differentiating x\^{}beta x\textsuperscript{alpha-\textbar x\textbar{}}2
delta\_alpha\_beta costs at most 4R, and 4R\textless=(1+R)\^{}2. The
multiplication coefficient in Q\^{}(2) is bounded by T\_e R/12 with the
same bound on its derivatives.

Formula (13) gives the exact field remainder

\[
 Q_{e\alpha}(g)=Q^{(0)}_{e\alpha}+gQ^{(1)}_{e\alpha}
                 +g^2Q^{(2)}_{e\alpha}+g^4F_{e\alpha,g},
\]

where

\[
 F_{e\alpha,g}=\sum_ct_{ec}|x_c|^2d(g|x_c|)
       \left([x_c\times(x_c\times e_\alpha)]\cdot\nabla_c
                                +x_c^\alpha\right).
\]

The chain rule, g\textless=1, and (30) bound every coefficient and first
derivative by M\_4,e. The factor 10 in (31) bounds, in particular, the
derivative terms
6R\textsuperscript{3\textbar d\textbar+gR}4\textbar d'\textbar{} from
the vector coefficient and
3R\textsuperscript{2\textbar d\textbar+gR}3\textbar d'\textbar{} from
its multiplication coefficient. The even smooth extension handles x\_c=0
as well.

Let Q\_app=Q\textsuperscript{(0)+gQ}(1)+g\textsuperscript{2Q}(2).
Subtracting the square through order two from the exact square leaves

\[
 g^3\{Q^{(1)},Q^{(2)}\}+g^4(Q^{(2)})^2
       +g^4\{Q_{\rm app},F_g\}+g^8F_g^2,
\]

where \{P,Q\}=PQ+QP. Apply (34), use g\textless=1, multiply by
2\textbar h\_e\textbar/a, and sum over the three alpha values and all
original edges. This gives exactly the first line of (32). All
coefficient derivatives in the composition are retained.

For the potential let A\_m=v\_m\^{}alpha T\_alpha in (21). Its operator
norm is \textbar v\_m\textbar/2. At real g each exponential is unitary.
The product rule for the fifth derivative and the multinomial theorem
therefore give

\[
 \left|\frac{d^5}{dg^5}\operatorname{tr}
                   \prod_{m=1}^4e^{gA_m}\right|
 \le2\left(\sum_m\|A_m\|\right)^5
 \le2\left(\frac{q_pR}{2}\right)^5.
\]

The integral Taylor remainder after degree four is bounded by this
number times g\^{}5/5!. Multiplication by the original 1/(2g\^{}2a),
summation with \textbar h\_p\textbar, and
\textbar\textbar f\textbar\textbar{}\emph{2\textless=\textbar\textbar f\textbar\textbar{}}(2,Sigma)
give the second line of (32). This proves (33) without a uniform-volume
assumption or an estimate on unspecified unbounded states outside the
compact chart core.

\subsection{8. Residual gauge invariance, parity, and exact
scope}\label{weak-coupling-part-14-residual-gauge-invariance-parity-and-exact-scope}

The remaining SU(2) gauge action rotates every x\_c by the same adjoint
matrix in SO(3). Dot products and scalar triple products in (23)--(24)
are invariant under that action; the epsilon and delta tensors in
(16)--(17) transform covariantly in their derivative indices. The
original graph numbers t,s and the physical weights carry no color
indices. Consequently H\_n and D\_\{h,n\} preserve the exact residual
physical subspace. The density and the compact radial cores are
invariant under the same rotations.

Let parity act only on the Lie-algebra coordinates by
\((\mathcal Pf)(x)=f(-x)\). It does not change the physical edge
positions, spacings, graph orientations, or weights. The coefficients
A\_0,A\_1,A\_2 have degrees zero, one, and two; W\_2,W\_3,W\_4 have
degrees two, three, and four. Since a derivative changes sign under this
parity,

\[
 \mathcal PH_0\mathcal P=H_0,\qquad
 \mathcal PH_1\mathcal P=-H_1,\qquad
 \mathcal PH_2\mathcal P=H_2,
                                                               \tag{35}
\]

and the same identities hold for D\_\{h,0\},D\_\{h,1\},D\_\{h,2\}. All
these expressions are formally symmetric on the smooth compact core;
this follows also by expanding the exact symmetric weighted operator
with the norm remainder (33).

Simultaneous x-\textgreater-x is not an adjoint SO(3) rotation. It
therefore does not remove the odd cubic gauge-invariant terms. For the
even full Gaussian Phi\_0 of H\_0, \textless Phi\_0,H\_1
Phi\_0\textgreater=0 follows exactly from (35): H\_1 has polynomial
coefficients, so this Gaussian matrix element is integrable, and
changing x to -x reverses its sign. The same is true of
\textless Phi\_0,D\_\{h,1\}Phi\_0\textgreater. These calculated matrix
elements are not an assumption of differentiable eigenvalues or of a
globally convergent perturbation series at g=0.

Equations (27)--(33) give the full requested operator and
weighted-observable coefficients with a controlled local remainder.
Using them in vacuum, spectral, or simultaneous continuum perturbation
theory requires actual estimates for the relevant states and tails; the
compact-core statement does not replace those estimates. The original
nonlinear operator and every finite graph interaction remain the objects
from which these coefficients were computed.

\subsection{Proof dependencies}\label{weak-coupling-part-14-proof-dependencies}

The exact maximal-tree quotient, Haar measure, and full finite-box
operator are proved in \emph{The original finite-box SU(2) Hamiltonian
at small positive coupling}. The full-measure dilation isometry and
physical Gaussian are proved in \emph{Actual weak-coupling vacuum,
observable and spectral maps}. The original weighted local-energy
operator and its exact vacuum/covariance response at positive coupling
are retained in \emph{Local-energy spacetime kernel and the exact
nonlinear coupling response}. The present note evaluates the higher
graph coefficients directly from those exact maps.

\clearpage
\part{Cubic singlets, vacuum energy and the first physical gap matrix}
\section{The cubic color singlet and the first vacuum-energy
correction}\label{weak-coupling-part-15-the-cubic-color-singlet-and-the-first-vacuum-energy-correction}

9 September 2026. All formulas refer to the original open-box SU(2)
Hamiltonian, including every edge and plaquette. This calculation
evaluates its first interaction coefficients; it does not replace that
Hamiltonian by a three-mode model. A sum over three distinct modes below
ranges over all such triples in the complete finite box.

\subsection{1. Original operator and the coordinate
dictionary}\label{weak-coupling-part-15-original-operator-and-the-coordinate-dictionary}

Vertices are \(n\in\{-L,\ldots,L\}^3\), \(L\ge2\), and edges are all
positive contained nearest-neighbor edges. Faces carry the oriented word
\(U_i(n)U_j(n+e_i)U_i(n+e_j)^{-1}U_j(n)^{-1}\), \(i<j\). Use
\(T_\alpha=-i\sigma_\alpha/2\), Haar probability measure, and \[
 H_g={2g^2\over a}\sum_e E_e+
       {1\over2g^2a}\sum_p(2-\operatorname{tr}U_p),\qquad
 E_e=-\sum_\alpha X_{e,\alpha}^2,\quad g,a>0.
 \tag{1}
\] The physical Hilbert space is the invariant subspace under all vertex
gauge transformations. Its vacuum is the positive unit eigenfunction. No
trial state is divided by its norm in this calculation.

Use the rooted tree whose parent decreases the first available
coordinate in the order \(1,2,3\). Write \(\mathcal C\) for its chords
and \(r=|\mathcal C|=2(2L)^3+3(2L)^2\). For a vertex \(v\), \(t_v\) is
its ordered root-to-vertex holonomy. The exact chord variable is
\(Z_c=t_{s(c)}U_ct_{t(c)}^{-1}\). Tree links and these chord variables
give a Haar-measure-preserving bijection, by independent left/right
translations of the chord factors for fixed tree links. Vertex gauge
reduction leaves simultaneous conjugation of all \(Z_c\).

For an edge cochain \(u\), let \(p_v(u)\) be its additive integral on
the same tree path and define the full-row-rank matrix \[
 (\mathsf T u)_c=p_{s(c)}(u)+u_c-p_{t(c)}(u),\qquad
 G=\mathsf T\mathsf T^*,\qquad C=d_1j.
 \tag{2}
\] Here \(j\) inserts chord values and zero tree values; \(d_1\) is the
oriented face sum in (1). The kernel of \(\mathsf T\) consists exactly
of gradients: path telescoping proves one inclusion and the vertex
potentials \(p_v(u)\) prove the converse. Likewise
\(\ker d_1=\operatorname{im}d_0\), by interchanging adjacent coordinate
steps in a monotone path. Consequently \(G>0\), \(C\) is injective, and
the eigenvalues of \(G^{1/2}C^*CG^{1/2}\) are strictly positive.

Choose an orthogonal diagonalization \[
 O^*G^{1/2}C^*CG^{1/2}O=\Sigma^2,\qquad
 \Sigma=\operatorname{diag}(\sigma_1,\ldots,\sigma_r),\quad \sigma_\mu>0.
 \tag{3}
\] This is the complete spectrum, not a selection of low modes.
Explicitly \(\sigma^2=4\sum_{i=1}^3\sin^2(\pi j_i/[2(2L+1)])\), with
\(j_i\in\{0,\ldots,2L\}\), at least two positive indices, and
multiplicity one less than their number of positive indices. Tensoring
the interval cosine vertex modes and sine edge modes proves this
formula: the gradient at each triple is the vector of its three positive
square roots, and the curl squared is their squared norm times the
projection perpendicular to that gradient.

Set \[
 \mathsf A=G^{1/2}O,\quad
 B={1\over4}G^{-1/2}O\Sigma O^*G^{-1/2},\quad
 \mathsf N=B\mathsf A={1\over4}G^{-1/2}O\Sigma,\quad M={1\over2}B^{-1}.
 \tag{4}
\] These matrices act on chord indices, identically in each of the three
color coordinates. The exact local maps are \(Z_c=\exp(gx_c^\alpha
T_\alpha)\) and \(x=\mathsf A z\). The Haar density is \[
 \mathcal J(y)=(16\pi^2)^{-r}
  \prod_c\left({\sin(|y_c|/2)\over |y_c|/2}\right)^2.
 \tag{5}
\] The local norm-preserving map is
\(F(\exp(gx))=g^{-3r/2}\mathcal J(gx)^{-1/2}f(x)\); the linear
\(x=\mathsf A z\) map retains the absolute Jacobian \(\det(G)^{3/2}\).
The comparison vacuum in \(x\) is \[
 \Phi(x)=\det(B)^{3/4}\pi^{-3r/4}
           \exp\!\left(-{1\over2}\sum_\alpha(x^\alpha)^*Bx^\alpha\right).
 \tag{6}
\] Thus \(\partial_{x_c^\alpha}\Phi=-(Bx)_c^\alpha\Phi\). Its exact
covariance is
\(\langle x_c^\alpha x_d^\beta\rangle=M_{cd}\delta_{\alpha\beta}\). In
\(z\), the covariance of each color component of mode \(\mu\) is
\(2/\sigma_\mu\). The comparison operator and its vacuum energy are \[
 H_0=\sum_{\mu=1}^r
       \left(-{2\over a}\Delta_{z_\mu}
                    +{\sigma_\mu^2\over8a}|z_\mu|^2\right),\qquad
 e_0={3\over2a}\sum_\mu\sigma_\mu.
 \tag{7}
\]

\subsection{2. All tree incidence and first kinetic
coefficients}\label{weak-coupling-part-15-all-tree-incidence-and-first-kinetic-coefficients}

Let \(\ell_{ec},r_{ec}\) be the coefficients of the left and right chord
derivatives in the exact tree-reduced edge field. For a chord edge
\(e\), \(\ell_{ec}=\delta_{ec},r_{ec}=0\). For a tree edge, they are
respectively \(\epsilon(s(c),e)\) and \(-\epsilon(t(c),e)\), where
\(\epsilon(v,e)\) records membership in the root-to-\(v\) path. Put \[
 t_{ec}=\ell_{ec}+r_{ec},\qquad s_{ec}=\ell_{ec}-r_{ec}.
 \tag{8}
\] In particular \(\sum_e t_{ec}t_{ed}=G_{cd}\). Differentiating
\(\log(e^{uT_\alpha}e^y)\) and \(\log(e^ye^{uT_\alpha})\) at \(u=0\)
gives \(I-\operatorname{ad}_y/2+\operatorname{ad}_y^2/12+O(|y|^3)\) and
\(I+\operatorname{ad}_y/2+\operatorname{ad}_y^2/12+O(|y|^3)\),
respectively. The signs follow by differentiating the exponential series
and solving its tangent equation order by order, using
\([v^\alpha T_\alpha,w^\beta T_\beta]=(v\times w)^\gamma T_\gamma\).

For the principal matrix after \(y=gx\), write
\(A(gx)=G\otimes I_3+gA_1(x)+g^2A_2(x)+O(g^3)\). Squaring all the
preceding edge fields gives \[
 (A_1)^{c\beta,d\gamma}
 ={1\over2}\sum_e\epsilon_{\eta\beta\gamma}
   \{s_{ec}t_{ed}x_c^\eta-t_{ec}s_{ed}x_d^\eta\},
 \tag{9}
\] \[
 \begin{split}
 (A_2)^{c\beta,d\gamma}=\sum_e\{&
 {s_{ec}s_{ed}\over4}
 [\delta_{\beta\gamma}x_c\cdot x_d-x_d^\beta x_c^\gamma]\\
 &+{t_{ec}t_{ed}\over12}
 [x_c^\beta x_c^\gamma+x_d^\beta x_d^\gamma
                 -( |x_c|^2+|x_d|^2)\delta_{\beta\gamma}]\}.
 \end{split}
 \tag{10}
\] The cross-product identity \(\sum_\alpha(x_c\times e_\alpha)_\beta
(x_d\times e_\alpha)_\gamma
=\delta_{\beta\gamma}x_c\cdot x_d-x_d^\beta x_c^\gamma\) proves the
first term of (10). The identity
\(x\times(x\times e_\alpha)=x(x\cdot e_\alpha)-|x|^2e_\alpha\) proves
its second term.

The derivative contraction \(\partial_i(A_1)^{ij}\) is zero: each
derivative identifies two indices in the alternating tensor in (9). Also
\(\operatorname{tr}(A_1(B\otimes I_3))=0\), because that trace
identifies the two color indices. Put \(p_c=(Bx)_c\). Substitution in
(9), followed by exchanging \(c,d\) in its second summand, yields the
exact polynomial identity \[
 p^*A_1p=\sum_{e,c,d}s_{ec}t_{ed}
                  x_c\cdot(p_c\times p_d).
 \tag{11}
\] No integration or parity argument is needed for this identity.

\subsection{3. Retained ordered Wilson
terms}\label{weak-coupling-part-15-retained-ordered-wilson-terms}

For each original face retain its four signed chord letters
\(v_{p,i}=\varepsilon_{p,i}x_{c(p,i)}\); a tree letter is the zero
vector, and an inverse chord has sign \(-1\). Write
\(W(gx)=g^2W_2+g^3W_3+g^4W_4+O(g^5)\). Direct multiplication gives \[
 W_2={1\over4}\sum_p\left|\sum_{i=1}^4v_{p,i}\right|^2,\qquad
 W_3={1\over4}\sum_p\sum_{i<j<k}
                   \det(v_{p,i},v_{p,j},v_{p,k}).
 \tag{12}
\] For a single face, its quartic contribution is the negative of \[
 \begin{split}
 &\sum_i {|v_i|^4\over192}
 +\sum_{i<j}{(|v_i|^2+|v_j|^2)(v_i\cdot v_j)\over48}
 +\sum_{i<j}{|v_i|^2|v_j|^2\over32}\\
 &+\sum_{i<j<k}
 { |v_i|^2(v_j\cdot v_k)+|v_j|^2(v_i\cdot v_k)
                         +|v_k|^2(v_i\cdot v_j)\over16}\\
 &+{(v_1\cdot v_2)(v_3\cdot v_4)
       -(v_1\cdot v_3)(v_2\cdot v_4)
       +(v_1\cdot v_4)(v_2\cdot v_3)\over8}.
 \end{split}
 \tag{13}
\] Indeed \(V_i=v_i^\alpha T_\alpha\) satisfies \(V_i^2=-|v_i|^2I/4\).
Moreover \(\operatorname{tr}V_iV_j=-v_i\cdot v_j/2\),
\(\operatorname{tr}V_iV_jV_k=-\det(v_i,v_j,v_k)/4\), and the four-factor
trace is the last bracket in (13) divided by \(8\). The degree-four
distributions among ordered letters are \(4\), \(3+1\), \(2+2\),
\(2+1+1\), and \(1+1+1+1\). Their exponential factorials give precisely
(13).

The half-density map (5) contributes the scalar
\(-\operatorname{tr}G/(2a)\) at order \(g^2\). For clarity,
\(\log\mathcal J(gx)=\mathrm{constant}
-g^2\sum_c|x_c|^2/12+O(g^4|x|^4)\). In the exact kinetic quadratic form
its shifted derivative is \(\partial+g^2x/12+O(g^4)\). The two cross
terms with \(G\otimes I_3\), integrated by parts, give
\(-(2/a)g^2\operatorname{tr}(G\otimes I_3)/12
=-g^2\operatorname{tr}G/(2a)\). There is no order-\(g\) density term.

Consequently the two interaction coefficients are \[
 H_1=-{2\over a}\partial_i(A_1)^{ij}\partial_j+{W_3\over2a},
 \qquad
 H_2=-{2\over a}\partial_i(A_2)^{ij}\partial_j
                  -{\operatorname{tr}G\over2a}+{W_4\over2a}.
 \tag{14}
\] They preserve simultaneous SO(3) color rotations. Simultaneous
inversion \(x\mapsto-x\) changes the sign of \(H_1\), not of \(H_2\).
Inversion is not an adjoint SU(2) rotation; its parity does not remove
the odd rotation singlets.

\subsection{4. Evaluated three-mode singlets, not an unspecified
inverse}\label{weak-coupling-part-15-evaluated-three-mode-singlets-not-an-unspecified-inverse}

Equations (6),(9),(11),(12),(14) give \[
 {H_1\Phi\over\Phi}
 =-{2\over a}\sum_{e,c,d}s_{ec}t_{ed}\det(x_c,p_c,p_d)
     +{1\over8a}\sum_{p,i<j<k}\det(v_{p,i},v_{p,j},v_{p,k}).
 \tag{15}
\] Define \(a_{p,i,\mu}=\varepsilon_{p,i}\mathsf A_{c(p,i),\mu}\) for a
chord letter and zero for a tree letter. The fully specified tensor \[
 Q_{\mu\nu\rho}=
 -{2\over a}\sum_{e,c,d}s_{ec}t_{ed}
                      \mathsf A_{c\mu}\mathsf N_{c\nu}\mathsf N_{d\rho}
 +{1\over8a}\sum_{p,i<j<k}
                         a_{p,i,\mu}a_{p,j,\nu}a_{p,k,\rho}
 \tag{16}
\] uses only the original incidences and the complete eigenbasis (3).
For \(\mu<\nu<\rho\) set \[
 \Theta_{\mu\nu\rho}
       =\sum_{\pi\in S_3}\operatorname{sgn}(\pi)
                   Q_{\pi(\mu,\nu,\rho)}.
 \tag{17}
\] This is the six-term sum, with no factor \(1/6\). Multilinearity of
determinant proves \[
 H_1\Phi=\sum_{\mu<\nu<\rho}
              \Theta_{\mu\nu\rho}\det(z_\mu,z_\nu,z_\rho)\Phi.
 \tag{18}
\] Repeated spatial indices vanish because the corresponding determinant
has two equal vector arguments. Equation (18) is not a claim that every
coefficient is nonzero; coefficients are exactly (16)--(17).

Put \(D_{\mu\nu\rho}=\det(z_\mu,z_\nu,z_\rho)\Phi\). It is an SO(3)
singlet and is orthogonal to \(\Phi\). Applying (7) to a polynomial
times \(\Phi\) gives the exact conjugated operator \[
 (H_0-e_0)(P\Phi)
 =\left[-{2\over a}\sum_\eta\Delta_{z_\eta}P
       +{1\over a}\sum_\eta\sigma_\eta z_\eta\cdot\nabla_{z_\eta}P\right]\Phi.
 \tag{19}
\] Each determinant is linear in each of its three distinct vector
arguments, hence its Laplacian vanishes in every mode. Thus \[
 (H_0-e_0)D_{\mu\nu\rho}
      ={\sigma_\mu+\sigma_\nu+\sigma_\rho\over a}D_{\mu\nu\rho},
 \qquad
 \|D_{\mu\nu\rho}\|^2={48\over\sigma_\mu\sigma_\nu\sigma_\rho}.
 \tag{20}
\] For the norm, expand both determinants as alternating sums.
Independence of the three modes forces matching color indices; six
permutations remain, each contributing \(8/(\sigma_\mu\sigma_\nu
\sigma_\rho)\). Different unordered spatial triples are orthogonal: a
mode occurring in only one triple has one centered coordinate factor.
This also proves \(\langle\Phi,H_1\Phi\rangle=0\).

Let \(R_0\) be the reduced inverse of \(H_0-e_0\). Its value on (18) is
therefore explicitly \[
 R_0H_1\Phi=\sum_{\mu<\nu<\rho}
 {a\Theta_{\mu\nu\rho}\over\sigma_\mu+\sigma_\nu+\sigma_\rho}
                         D_{\mu\nu\rho},
 \tag{21}
\] and \[
 \langle H_1\Phi,R_0H_1\Phi\rangle
 =\sum_{\mu<\nu<\rho}
 {48a\,\Theta_{\mu\nu\rho}^2\over
      \sigma_\mu\sigma_\nu\sigma_\rho(\sigma_\mu+\sigma_\nu+\sigma_\rho)}.
 \tag{22}
\] The first vacuum-vector coefficient is minus (21). Its squared norm
has the same summands with \(a\) replaced by \(a^2\) and the last
frequency sum squared. These are sums over the entire graph, not three
chosen oscillators.

\subsection{5. The evaluated quartic and measure
contributions}\label{weak-coupling-part-15-the-evaluated-quartic-and-measure-contributions}

For each face define the real symmetric four-letter covariance \[
 m_{ij}=\varepsilon_{p,i}\varepsilon_{p,j}
                             M_{c(p,i),c(p,j)};
 \tag{23}
\] a tree letter makes that entry zero. Define the following polynomial:
\[
 \begin{split}
 \mathcal W_4(m)=-\bigg\{&
 {15\over192}\sum_i m_{ii}^2
 +{15\over48}\sum_{i<j}(m_{ii}+m_{jj})m_{ij}\\
 &+{1\over32}\sum_{i<j}(9m_{ii}m_{jj}+6m_{ij}^2)\\
 &+{1\over16}\sum_{i<j<k}
 [9(m_{ii}m_{jk}+m_{jj}m_{ik}+m_{kk}m_{ij})
       +6(m_{ij}m_{ik}+m_{ij}m_{jk}+m_{ik}m_{jk})]\\
 &+{9m_{12}m_{34}-3m_{13}m_{24}+9m_{14}m_{23}\over8}
 \bigg\}.
 \end{split}
 \tag{24}
\] Then \(\langle W_4\rangle=\sum_p\mathcal W_4(m^{(p)})\). To verify it
directly, Gaussian integration gives \[
 \begin{split}
 \langle |v_i|^4\rangle&=15m_{ii}^2,\\
 \langle |v_i|^2(v_i\cdot v_j)\rangle&=15m_{ii}m_{ij},\\
 \langle |v_i|^2|v_j|^2\rangle&=9m_{ii}m_{jj}+6m_{ij}^2,\\
 \langle |v_i|^2(v_j\cdot v_k)\rangle
                         &=9m_{ii}m_{jk}+6m_{ij}m_{ik},\\
 \langle(v_i\cdot v_j)(v_k\cdot v_l)\rangle
                         &=9m_{ij}m_{kl}+3m_{ik}m_{jl}+3m_{il}m_{jk}.
 \end{split}
 \tag{25}
\] These identities follow by differentiating the exact Gaussian
generating function
\(\langle e^{u\cdot v}\rangle=e^{u^*\operatorname{Cov}(v)u/2}\) four
times; the three pairings are all retained. Applying (25) to each term
of (13) proves (24), including its middle negative coefficient.

The kinetic expectation can also be evaluated, without differentiating
an unspecified Gaussian integral. Put \(p=Bx\); then
\(\langle p_c^\alpha p_d^\beta\rangle=B_{cd}\delta_{\alpha\beta}/2\) and
\(\langle x_c^\alpha p_d^\beta\rangle=
\delta_{cd}\delta_{\alpha\beta}/2\). Substitution in (10) gives \[
 \begin{split}
 \langle p^*A_2p\rangle=\sum_{e,c,d}\bigg[
 &s_{ec}s_{ed}\left({3\over4}B_{cd}M_{cd}-{3\over8}\delta_{cd}\right)\\
 +&t_{ec}t_{ed}\left({1\over4}\delta_{cd}
                       -{1\over4}B_{cd}(M_{cc}+M_{dd})\right)\bigg].
 \end{split}
 \tag{26}
\] For example the first bracket is one quarter of the difference
between \(\langle(p_c\cdot p_d)(x_c\cdot x_d)\rangle\) and
\(\langle(p_c\cdot x_d)(p_d\cdot x_c)\rangle\). Its two Wick expansions
give respectively \(9B_{cd}M_{cd}/2+3/4+3\delta_{cd}/4\) and
\(3B_{cd}M_{cd}/2+3/4+9\delta_{cd}/4\). For either \(v=x_c\) or
\(v=x_d\), the difference
\(\langle(p_c\cdot v)(p_d\cdot v)-|v|^2(p_c\cdot p_d)\rangle\) is
\(3\delta_{cd}/2-3B_{cd}M_{vv}\). Adding these and dividing by twelve
proves the second bracket.

Integration by parts in (14) now yields the fully evaluated coefficient
\[
 \boxed{\begin{split}
 e_2={}&{2\over a}\langle p^*A_2p\rangle
             -{\operatorname{tr}G\over2a}
             +{1\over2a}\sum_p\mathcal W_4(m^{(p)})\\
 &-\sum_{\mu<\nu<\rho}
 {48a\,\Theta_{\mu\nu\rho}^2\over
 \sigma_\mu\sigma_\nu\sigma_\rho(\sigma_\mu+\sigma_\nu+\sigma_\rho)} .
 \end{split}}
 \tag{27}
\] In (27), (26),(24),(16),(17) are finite explicitly specified sums;
there is no uncomputed reduced resolvent or intermediate-state symbol.
The separate Haar scalar is displayed, not absorbed into the energy
origin. It contributes to the absolute vacuum energy although scalar
shifts cancel in vacuum-subtracted spectral differences.

\subsection{6. What coefficient has been
calculated}\label{weak-coupling-part-15-what-coefficient-has-been-calculated}

At a fixed finite box, the exact chart expansion with Gaussian cutoffs
gives the actual vacuum expansion
\(\mathcal E_g=e_0+g^2e_2+O_{L,a}(g^3)\). Here is a direct proof with
the cutoff and actual spectral projection. The coefficient equations
follow by inserting \(\Phi+g\phi_1+g^2\phi_2\) and \(e_0+g^2e_2\): \[
 \phi_1=-R_0H_1\Phi,\qquad
 (H_0-e_0)\phi_2=-(H_1\phi_1+(H_2-e_2)\Phi).
 \tag{28}
\] Its right side is perpendicular to \(\Phi\) precisely for (27). The
unit-vacuum convention fixes
\(\langle\Phi,\phi_2\rangle=-\|\phi_1\|^2/2\). All vectors in (28) are
finite Gaussian polynomials; (19) determines the inverse coefficient by
coefficient in the complete Hermite basis. In particular the odd cubic
singlet changes the vacuum at order \(g\), even though its vacuum-energy
expectation vanishes at that order.

Choose a smooth invariant cutoff \(\chi(y)\), equal to one for
\(|y|\le\rho/2\) and zero for \(|y|\ge\rho\), with \(0<\rho<1\). On the
support of \(\chi_g(x)=\chi(gx)\), the exact field formulas give the
pointwise remainder bound \[
 |(\widetilde H_g-H_0-gH_1-g^2H_2)f(x)|
 \le C_{L,a}g^3(1+|x|)^8
                    \sum_{|\alpha|\le2}|\partial^\alpha f(x)|.
 \tag{28a}
\] To verify its polynomial degree, the exact logarithmic-field
remainder is \(g^4F_g\); its coefficients and first derivatives are
bounded by \(C(1+|x|)^4\) when \(g|x|<1\). Squaring the fields leaves
products of degrees at most eight, with powers of \(g\) at least three
after subtracting the first three coefficients. The Wilson fifth
derivative is bounded by \(2(\sum_i|v_i|/2)^5\), since each exponential
is unitary on the real integration interval. Its Taylor remainder,
multiplied by \(1/(2g^2a)\), obeys the same bound. These are finitely
many original graph terms, so their constants sum to \(C_{L,a}\).

Every derivative through order two of \(\chi_g(\Phi+g\phi_1+g^2\phi_2)\)
is bounded by a fixed polynomial times
\(e^{-\lambda_{\min}(B)|x|^2/2}\), uniformly for \(g\le1\). This follows
from the product rule, the finite Gaussian-polynomial form of
\(\phi_1,\phi_2\), and the factor \(g^{|\alpha|}\) in each cutoff
derivative. Integrating (28a) therefore gives an \(O(g^3)\) norm
remainder with no growing logarithmic cutoff factor. Commutators of
\(H_0,H_1,H_2\) with \(\chi_g\) are supported on \(|x|\ge\rho/(2g)\);
their coefficients and derivatives grow only polynomially. In polar
coordinates each such Gaussian tail is bounded by \(C_Ng^N\) for every
fixed \(N\), since a polynomial times \(e^{-c/g^2}\) is smaller than
every power of \(g\). This proves the cutoff error as well as the
coefficient error.

Let \(U_g\) be the inverse of the exact tree/Haar/dilation map on this
chart, and set \[
 q_g=U_g\chi_g(\Phi+g\phi_1+g^2\phi_2),\qquad
 \lambda_g=e_0+g^2e_2.
 \tag{28b}
\] The original vectors \(q_g\) are smooth and physical. The coefficient
equations, the preceding bounds, and odd/even parity give \[
 \|(H_g-\lambda_g)q_g\|=O_{L,a}(g^3),\qquad
                     \|q_g\|^2=1+O_{L,a}(g^4).
 \tag{28c}
\] The full fixed-box physical eigenvalue convergence isolates its
actual ground state \(\psi_g\) from the rest of the physical spectrum by
a positive fixed-box constant for small \(g\). The spectral theorem
applied to the residual in (28c) thus gives
\(\|(I-|\psi_g\rangle\langle\psi_g|)q_g\|=O(g^3)\). The real coefficient
\(\beta_g=\langle\psi_g,q_g\rangle\) has absolute value \(1+O(g^4)\). It
is positive: the leading chart vector \(U_g\chi_g\Phi\) is nonnegative
and its difference from \(q_g\) has norm \(O(g)\), excluding a
coefficient near \(-1\). Taking the vacuum component of the residual
proves \(|\mathcal E_g-\lambda_g|=O(g^3)\), and its orthogonal component
proves \[
 \|\psi_g-q_g\|+\|H_g(\psi_g-q_g)\|=O_{L,a}(g^3).
 \tag{28d}
\] For the graph term explicitly use
\((H_g-\mathcal E_g)q_g=(H_g-\lambda_g)q_g+
(\lambda_g-\mathcal E_g)q_g\) and boundedness of \(\mathcal E_g\) at
this fixed box. No differentiability assumption at \(g=0\) or
self-adjointness of a global cubic-polynomial truncation was used.

The graph estimate controls the original weighted energies too. Write
\(K=\sum E_e\) and \(W=\sum_p(2-\operatorname{tr}U_p)\). Each face trace
is a fundamental matrix coefficient in each of its four independent edge
variables. Its Casimir sum is three times that trace, so \(KW=3W-6|P|\).
Integration by parts gives, for every smooth original \(u\), \[
 \begin{split}
 \|H_gu\|^2={}&\|(2g^2/a)Ku\|^2+\|Wu/(2g^2a)\|^2\\
 &+{2\over a^2}\sum_{e,\alpha}\int W|X_{e,\alpha}u|^2
                 +{1\over a^2}\int(3W-6|P|)|u|^2.
 \end{split}
 \tag{28e}
\] Indeed \(\operatorname{Re}\langle Ku,Wu\rangle
=\sum\int W|Xu|^2+\tfrac12\int(KW)|u|^2\); the product of the two
original Hamiltonian coefficients is \(a^{-2}\). The sum of the first
two squared norms in (28e) is at most \(\|H_gu\|^2+6|P|\|u\|^2/a^2\).
Simultaneous Peter--Weyl diagonalization of the commuting nonnegative
\(E_e\) proves \(\|\sum h_eE_eu\|\le\|h\|_\infty\|Ku\|\), and pointwise
\(|W_h|\le\|h\|_\infty W\), with the original face average. Thus \[
 \|D_{h,g}u\|\le\sqrt2\,\|h\|_\infty
          [\|H_gu\|^2+6|P|\|u\|^2/a^2]^{1/2}.
 \tag{28f}
\] Smooth approximation in the \(H_g\) graph extends this inequality to
that domain. In particular (28d) supplies an \(O(g^3)\) error for each
fixed weighted energy vector, without an unproved tail interpolation or
a loss of \(g^{-2}\).

Every \(\sigma_\mu,\mathsf A,B,M,s,t\) is independent of \(a\). The
tensor \(\Theta\) carries exactly one factor \(1/a\), so
\(e_2=a^{-1}\mathfrak e_2(L)\) with \(\mathfrak e_2(L)\) given by (27)
after retaining that factor. No boundedness of \(\mathfrak e_2(L)\) as
\(L\to\infty\), or of the fixed-box remainder along a continuum
trajectory, is asserted. Their dependence is the next calculation, not a
permission to discard this coefficient. A vacuum-energy coefficient
alone is not an excitation-gap coefficient: local centered states and
the vacuum-subtracted semigroup must be expanded too.

\subsection{7. The first physical gap: a specified six-by-six
interaction
matrix}\label{weak-coupling-part-15-the-first-physical-gap-a-specified-six-by-six-interaction-matrix}

The vacuum coefficient is not enough for a gap calculation. The same
full graph gives a finite explicit matrix for the first physical
cluster. Let \(\sigma_*=\min_\mu\sigma_\mu\) and \(I_*\) be its three
spatial modes. Write \(\delta_*=2\sigma_*/a\). For \(\mu,\nu\in I_*\),
define the following six real polynomials: \[
 P_{\mu\mu}={\sigma_*\over2\sqrt6}
                 (|z_\mu|^2-6/\sigma_*),\qquad
 P_{\mu\nu}={\sigma_*\over2\sqrt3}z_\mu\cdot z_\nu\quad(\mu<\nu).
 \tag{29}
\] The vectors \(S_I=P_I\Phi\) are an orthonormal basis for the full
first physical cluster. Indeed their norms and mutual inner products
follow from the independent covariance \(2I_3/\sigma_*\); (19) gives
excitation \(\delta_*\). A one-quantum vector has no SO(3) invariant
component. Every state with at least two quanta has excitation at least
\(\delta_*\), with equality only for two lowest-mode quanta. The
invariant bilinear color tensor is a multiple of
\(\delta_{\alpha\beta}\): half-turns remove off-diagonal entries and
quarter-turns equate the diagonal entries. Symmetry in the mode pair
then gives precisely (29).

The formulas below compute every matrix element using the original
coefficients, with no numerical cutoff on intermediate oscillator modes.
In \(x\) coordinates, write \(\widehat P_I(x)=P_I(\mathsf A^{-1}x)\).
Let \(P_3(x)\) denote the cubic polynomial on the right of (15). Define
\[
 F_I(x)={4\over a}(A_1(x)p)\cdot\nabla\widehat P_I(x)
                                      +\widehat P_I(x)P_3(x).
 \tag{30}
\] Then \(H_1S_I=F_I\Phi\). Expanding \(H_1(\widehat P_I\Phi)\) with
(14) proves this: the omitted-looking term
\(-2A_1:\nabla^2\widehat P_I/a\) is identically zero, not discarded. The
Hessian of each polynomial (29), after the chord-index map, is a spatial
matrix times the color identity; (9) contracts it to zero. The
divergence and trace contractions already proved after (10) remove the
other two terms.

For a multi-index \(n=(n_{\mu\alpha})\), let \[
 h_n(z)=\prod_{\mu,\alpha}
 {\operatorname{He}_{n_{\mu\alpha}}(\sqrt{\sigma_\mu/2}\,
                                          z_\mu^\alpha)
          \over\sqrt{n_{\mu\alpha}!}},\qquad
 \omega_n={1\over a}\sum_{\mu,\alpha}\sigma_\mu n_{\mu\alpha}.
 \tag{31}
\] Here the probabilists' Hermite polynomial has the explicit expansion
\[
 \operatorname{He}_k(u)
 =k!\sum_{j=0}^{\lfloor k/2\rfloor}
                  {(-1)^j u^{k-2j}\over2^j j!(k-2j)!}.
 \tag{32}
\] The complete Hermite eigenbasis is \(h_n\Phi\), with excitation
\(\omega_n\). Put \[
 c_{I,n}=\langle h_n,F_I(\mathsf A z)\rangle_{\Phi^2}.
 \tag{33}
\] This is an evaluated finite prescription: substitute (30),(32) and
take moments by pairing coordinates with
\(\langle z_\mu^\alpha z_\nu^\beta\rangle
=2\delta_{\mu\nu}\delta_{\alpha\beta}/\sigma_\mu\). For degree \(2q\) a
monomial moment is the sum over all \((2q-1)!!\) pairings of its
coordinate slots, of the products of these displayed covariances. Odd
moments are zero. Thus (33) hides neither an integral requiring a new
measure nor an unspecified state-space truncation.

Only \(|n|=3\) and \(|n|=5\) occur. Polynomial (30) has degree at most
five and odd parity, so the only other possible degree is one. Its
projection to one-quantum states is zero: (30) is invariant under
simultaneous color rotation, that projection commutes with the
rotations, and the one-quantum subspace contains no invariant vector.
This proves the exclusion of degree one, independently of any accidental
frequency equality elsewhere in the unrestricted oscillator spectrum. In
particular every remaining denominator below satisfies
\(\omega_n-\delta_*\ge\sigma_*/a>0\).

All quartic cluster matrix elements are equally explicit. For real
polynomials \(P_I,P_J\), put \[
 u_I(x)=\nabla\widehat P_I(x)-p\,\widehat P_I(x).
 \tag{34}
\] Integration by parts in the full \(H_2\) gives \[
 V_{IJ}={2\over a}\langle u_I^*A_2u_J\rangle_{\Phi^2}
       -{\operatorname{tr}G\over2a}\delta_{IJ}
       +{1\over2a}\langle\widehat P_I W_4\widehat P_J\rangle_{\Phi^2}.
 \tag{35}
\] Every expectation in (35) is a polynomial moment of degree at most
eight, computed by the same complete pairing prescription, now using
\(M_{cd}\delta_{\alpha\beta}\) in \(x\). The full order-two physical gap
matrix is \[
 \boxed{\mathsf K_{IJ}=
 V_{IJ}-\sum_{\substack{n:\ |n|=3\ {\rm or}\ 5}}
              {c_{I,n}c_{J,n}\over\omega_n-\delta_*}
                                      -e_2\delta_{IJ}.}
 \tag{36}
\] The sum in (36) ranges over all \(3r\) real oscillator coordinates.
It is finite because its total degree is fixed, not because high spatial
modes have been removed. Formula (36), together with (8)--(17),(24),
(26),(29)--(35), is a complete finite algorithm in the actual graph
matrices, signs and frequencies. It has not been replaced by an
undetermined perturbation matrix.

Here is the spectral meaning and its fixed-box error justification.
First omit the last \(-e_2I\) in (36), calling the resulting matrix
\(\mathsf K^{\rm abs}\). With \(R_*\) the reduced inverse of
\(H_0-e_0-\delta_*\) off the six-dimensional physical cluster, set \[
 \eta_I=-R_*H_1S_I,\qquad
 (H_0-e_0-\delta_*)\zeta_I
 =-H_1\eta_I-H_2S_I+\sum_JS_J\mathsf K^{\rm abs}_{JI}.
 \tag{37}
\] Equation (33) explicitly evaluates \(\eta_I\) through its positive
denominators; orthogonality of the right side in the second equation
follows from (35)--(36). The right side is a finite Gaussian polynomial.
The inverse on its orthogonal complement is obtained by dividing its
Hermite coefficients by their nonzero excitation differences. Physical
invariance removes any nonphysical resonant components; in the physical
space the eigenspace at this energy is precisely the six vectors (29).
Add the cluster component
\(-\tfrac12\sum_JS_J\langle\eta_J,\eta_I\rangle\) to \(\zeta_I\). The
Gram matrix of \(S_I+g\eta_I+g^2\zeta_I\) is then \(I+O(g^3)\).

Apply a fixed original-chart radial cutoff and the exact map (5) to
these six polynomial Gaussians. The compact coefficient remainder and
Gaussian-tail argument (28a)--(28c) apply to this finite list exactly as
to the vacuum: its residual as a six-column map is \(O_{L,a}(g^3)\) for
the matrix \((e_0+\delta_*)I+g^2\mathsf K^{\rm abs}\). The previously
proved physical eigenvalue convergence isolates a six-dimensional actual
cluster with positive distance from the rest of the fixed-box physical
spectrum. Projecting the six columns to that cluster changes them by
\(O(g^3)\): first diagonalize the real symmetric \(\mathsf K^{\rm abs}\)
by an orthogonal six-by-six matrix, retain this change in the six
columns, and apply the off-cluster resolvent at each corresponding
\((e_0+\delta_*)+g^2\lambda_I(\mathsf K^{\rm abs})\) to its residual.
All six spectral separations remain positive at this fixed box for
sufficiently small \(g\). The projected Gram matrix is still
\(I+O(g^3)\). Multiplication by its explicitly retained positive inverse
square root makes a unitary identification with \(\mathbb R^6\) and
changes the represented energy matrix by \(O(g^3)\). The
finite-dimensional min--max principle bounds each eigenvalue error by
that matrix norm. Subtracting the actual vacuum expansion in (28)
therefore proves \[
 \boxed{\Delta_L(g,a)=
       {2\sigma_*\over a}+g^2\lambda_{\min}(\mathsf K)
                                     +O_{L,a}(g^3).}
 \tag{38}
\] This is the actual first physical gap at a fixed box. It is not a
uniform claim as the lattice and volume change. In particular (36) must
still be analyzed, rather than assigned a sign or neglected, to decide
what its correction does along a specified physical-scale trajectory.
The Haar scalar cancels between (35) and the same scalar in \(e_2\), as
it must for a gap; it remains present in the absolute vacuum and
absolute cluster energies.

\subsection{8. Exact cancellation of spectator vacuum
contributions}\label{weak-coupling-part-15-exact-cancellation-of-spectator-vacuum-contributions}

There is a further cancellation that is relevant before estimating any
large-box limit. It can be proved with the full mode set retained.
Factor the oscillator Hilbert space into the three lowest spatial modes
and all remaining modes, with vacua \(\Phi_*\) and \(\Phi_{\rm out}\).
This is only a tensor-factor identification of the complete comparison
space. The original Hamiltonian is not restricted to either factor.

To specify its operator coefficients in this identification, introduce
\[
 z_\mu^\alpha=\sqrt{2/\sigma_\mu}\,
                  (a_{\mu\alpha}+a_{\mu\alpha}^\dagger),\qquad
 \partial_{z_\mu^\alpha}=\sqrt{\sigma_\mu/8}\,
                  (a_{\mu\alpha}-a_{\mu\alpha}^\dagger).
 \tag{39}
\] Substitute \(x=\mathsf A z\) and
\(\partial_x=\mathsf A^{-*}\partial_z\) in (14), then use (39). Order
every creation operator to the left with the exact identity \[
 a_i a_j^\dagger=a_j^\dagger a_i+\delta_{ij}.
 \tag{40}
\] This is a terminating finite algorithm: each interchange reduces the
number of out-of-order pairs, and its additional contraction has two
fewer operator factors. No contraction is omitted. Each scalar left by
this procedure is retained. The resulting \(H_1\) has degrees at most
three and odd parity. A degree-one term would be a color-invariant
vector operator. Its coefficient must vanish under all color rotations,
just as a one-quantum invariant vector must vanish. Thus the nonzero
normal-ordered terms of \(H_1\) have degree three. The normal-ordered
\(H_2\) has degrees zero, two and four.

Define \(H_{n,\rm out}\) by collecting all these computed monomials
whose spatial mode indices lie outside \(I_*\); include the scalar term
in \(H_{2,\rm out}\). Define \[
 H_{n,\rm touch}=H_n-H_{n,\rm out},\qquad n=1,2.
 \tag{41}
\] Every monomial of \(H_{n,\rm touch}\) contains at least one operator
from a lowest mode. This decomposition specifies its coefficients by
(8)--(14),(39)--(40); it is not a discarded part of the calculation.
Because color rotations preserve the mode labels, each summand remains
an invariant operator.

For any first-cluster vector \(S_I=P_I\Phi_*\otimes\Phi_{\rm out}\),
\(H_{1,\rm out}S_I\) has exactly two lowest-mode quanta and three
outside quanta. Indeed only three-creation monomials of
\(H_{1,\rm out}\) survive on \(\Phi_{\rm out}\). Conversely
\(H_{1,\rm touch}S_J\) has at most two outside quanta, since every
degree-three monomial has at least one lowest-mode factor. These two
vectors are orthogonal, even after inserting the oscillator reduced
inverse \(R_*\), because that inverse is diagonal in the
occupation-number basis. No energy estimate is being substituted for
this orthogonality.

On the first vector, \(R_*\) acts on its outside factor by the inverse
outside excitation energy: the two lowest-mode quanta have energy
\(\delta_*\), exactly canceled by the reference energy in \(R_*\).
Therefore \[
 \langle H_{1,\rm out}S_I,R_*H_{1,\rm out}S_J\rangle
 =\delta_{IJ}\,
       \langle H_{1,\rm out}\Phi,R_0H_{1,\rm out}\Phi\rangle.
 \tag{42}
\] In the vacuum, \(H_{1,\rm out}\Phi\) has no lowest-mode quanta,
whereas each surviving three-creation term of \(H_{1,\rm touch}\Phi\)
has at least one. Their \(R_0\)-inserted cross terms vanish for the same
occupation-number reason.

The quartic and density contributions cancel in the corresponding way,
without an inverse: \[
 \langle S_I,H_{2,\rm out}S_J\rangle
       =\delta_{IJ}\langle\Phi,H_{2,\rm out}\Phi\rangle .
 \tag{43}
\] This identity includes every scalar arising from (40) as well as the
original Haar scalar. Normal ordering also gives
\(\langle\Phi,H_{2,\rm touch}\Phi\rangle=0\): each of its monomials
contains a lowest-mode creation or annihilation operator, and the
normal-ordered vacuum expectation of such a monomial is zero.

Combining (42)--(43) with the full, unaltered gap matrix (36) proves the
alternative exact formula \[
 \boxed{\begin{split}
 \mathsf K_{IJ}={}&
  \langle S_I,H_{2,\rm touch}S_J\rangle
       -\delta_{IJ}\langle\Phi,H_{2,\rm touch}\Phi\rangle\\
 &-\langle H_{1,\rm touch}S_I,R_*H_{1,\rm touch}S_J\rangle
       +\delta_{IJ}\langle H_{1,\rm touch}\Phi,
                                      R_0H_{1,\rm touch}\Phi\rangle .
 \end{split}}
 \tag{44}
\] The first vacuum expectation displayed here is zero as just proved.
The last expectation is (22) with precisely those triples intersecting
\(I_*\), not the sum of every outside vacuum triple. The middle inverse
is still evaluated by (31)--(33), and now its vectors have at most two
outside quanta. This proves that contributions acting exclusively on
spectator modes cancel from the physical gap, including their full
denominators. It does not assert convergence of the remaining outside
mode sums or identify a renormalized coupling.

Equations (36) and (44) retain both sides of the cancellation
explicitly. They explain why a divergent absolute vacuum-energy
coefficient cannot by itself determine the gap correction, and give a
more targeted exact sum for the next volume and physical-scale
estimates.

\subsection{Sources and proof
dependencies}\label{weak-coupling-part-15-sources-and-proof-dependencies}

The included \emph{The original finite-box SU(2) Hamiltonian at small
positive coupling} proves the global tree map, Haar identity, physical
operator, unique well and actual finite-box spectral convergence.
\emph{Full graph operator coupling coefficients} supplies the exact
chart remainder and weighted versions of (14). Sections 6--7 above
promote the coefficients to the actual finite-box vacuum, weighted graph
vectors and first physical spectral cluster.

J. Kogut and L. Susskind, \emph{Hamiltonian formulation of Wilson's
lattice gauge theories}, Physical Review D 11 (1975), 395--408,
\href{https://doi.org/10.1103/PhysRevD.11.395}{doi:10.1103/PhysRevD.11.395},
is the foundational Hamiltonian reference. Its coupling conventions are
not substituted for the explicitly retained coefficients in (1).

S. Chatterjee, \emph{Yang--Mills for probabilists},
\href{https://arxiv.org/abs/1803.01950}{arXiv:1803.01950}, sections on
Wilson loops and defining the continuum limit. The preserved primary TeX
was read for the relation between lattice correlations, physical scale
and the continuum limit; it does not prove a volume-uniform bound on the
new finite-box sums (27).

\clearpage
\part{Exact cubic-box symmetry and the three gap corrections}
\section{Exact cubic-box symmetry of the first physical gap
correction}\label{weak-coupling-part-16-exact-cubic-box-symmetry-of-the-first-physical-gap-correction}

9 September 2026. The original open box, physical spacings, all SU(2)
edge links and the full Wilson Hamiltonian are retained. This note
proves that its six-by-six order-\(g^2\) first-gap matrix has three
independent entries. It is a symmetry of the original operator, not an
approximation to that matrix or a removal of spatial modes.

\subsection{1. The original graph
automorphisms}\label{weak-coupling-part-16-the-original-graph-automorphisms}

Let \(V=\{-L,\ldots,L\}^3\), with all contained positive unit edges and
elementary square faces. For a signed permutation matrix \(R\),
\(n\mapsto Rn\) is a bijection of \(V\). If \(e=(n,n+e_i)\) is a
positive edge, its image is the directed segment from \(Rn\) to
\(Rn+Re_i\). Store its original link if this direction is positive, and
the inverse original link if it is negative. Equivalently extend link
configurations to both orientations by \(U_{-e}=U_e^{-1}\), and define
\[
 (\alpha_R U)_{Re}=U_e
 \quad\hbox{for every directed edge }e .
 \tag{1}
\] The convention on reversed edges makes (1) a bijection of the
original positive-link configuration space, with inverse
\(\alpha_{R^{-1}}\). Define the unitary on functions by \[
 (\mathscr S_R F)(U)=F(\alpha_{R^{-1}}U).
 \tag{2}
\] It is unitary for product Haar probability: it permutes the factors
and possibly applies inversion, each of which preserves Haar measure.
Composition is exactly \(\mathscr S_R\mathscr S_Q=\mathscr S_{RQ}\).

Link inversion interchanges left and right Casimirs, which coincide on
SU(2): the adjoint matrix relating their three derivatives is
orthogonal, or equivalently both are the bi-invariant group Laplacian.
Thus (2) preserves \(\sum_e E_e\). The image of an original face is the
boundary of another contained square, with a possible cyclic change of
base vertex and orientation reversal. Cyclic reordering preserves the
trace; reversal takes the inverse of its holonomy. For SU(2),
\(\operatorname{tr}U^{-1}=\operatorname{tr}U\), since the eigenvalues
are \(e^{i\theta},e^{-i\theta}\). Therefore (2) also preserves
\(\sum_p(2-\operatorname{tr}U_p)\).

Vertex gauge transformations are carried to vertex gauge transformations
by \(h_{Rn}=h_n\); substitution in (1) proves the compatibility for both
edge orientations. Consequently \(\mathscr S_R\) preserves the physical
Hilbert space and commutes with the exact original operator \[
 H_g={2g^2\over a}\sum_e E_e+
                 {1\over2g^2a}\sum_p(2-\operatorname{tr}U_p).
 \tag{3}
\] The positive unique unit vacuum is fixed by every \(\mathscr S_R\).
Neither \(g\) nor \(a\) is changed.

\subsection{2. The exact linear action on the three lowest
modes}\label{weak-coupling-part-16-the-exact-linear-action-on-the-three-lowest-modes}

Write \(m=2L\), \(N=m+1\). For the one-dimensional interval with
vertices \(l=0,\ldots,m\) and edges \(l=0,\ldots,m-1\), retain \[
 v_0(l)=N^{-1/2},\quad
 v_1(l)=\sqrt{2/N}\cos{\pi(l+1/2)\over N},\quad
 w_1(l)=-\sqrt{2/N}\sin{\pi(l+1)\over N}.
 \tag{4}
\] The difference operator obeys \(Dv_1=s w_1\), where
\(s=2\sin(\pi/(2N))\). For \(i<j\), take the two orthonormal edge
tensors at the frequency triple with ones in positions \(i,j\) and zero
elsewhere: \(b_i^{ij}\) uses \(w_1\) in direction \(i\), \(v_1\) in
direction \(j\), and \(v_0\) in the remaining direction; \(b_j^{ij}\) is
the corresponding tensor with \(i,j\) interchanged. Define the unit
transverse cochain \[
 V_{ij}={b_i^{ij}-b_j^{ij}\over\sqrt2},\qquad
                  V_{ji}=-V_{ij}.
 \tag{5}
\] The gradient at this triple has components \((s,s)\), so (5) is
perpendicular to it. The oriented curl squared on this perpendicular
line is \(2s^2\). Thus these are the three lowest modes, with
\(\sigma_*=\sqrt2\,s\). Completeness of the original interval tensor
bases proves that there are no additional modes at this value.

For a permutation \(\pi\) of coordinate axes, the induced oriented
cochain action sends \[
 V_{ij}\longmapsto V_{\pi(i)\pi(j)}.
 \tag{6}
\] This includes the sign in (5) if the new indices are reversed. For a
reflection in coordinate \(r\), a vertex index becomes \(m-l\) and an
edge index becomes \(m-1-l\), with a minus sign for an oriented edge
component in that direction. Directly from (4), \[
 v_1(m-l)=-v_1(l),\quad v_0(m-l)=v_0(l),\quad
 w_1(m-1-l)=w_1(l).
 \tag{7}
\] Thus a reflected direction in \(\{i,j\}\) changes both terms of (5)
by a minus sign: the sine edge term acquires its orientation minus, and
the other term acquires the cosine minus. A reflection in the remaining
direction does not change either. For
\(\operatorname{diag}(\eta_1,\eta_2,\eta_3)\), this is \[
 V_{ij}\longmapsto \eta_i\eta_j V_{ij}.
 \tag{8}
\] Equations (6),(8) identify the full three-dimensional lowest spatial
mode representation with the exterior square of the original signed
permutation action on the coordinate axes. This is an explicit
intertwining map \(e_i\wedge e_j\mapsto V_{ij}\), including every sign.
It does not identify spatial modes with color coordinates.

In the complete transverse mode construction, choose the first three
orthonormal mode columns to be (5), in the order \(12,13,23\), and keep
the full orthogonal complement as well. The tree-coordinate map
\(x=\mathsf A z\), \(\mathsf A=G^{1/2}O\), is the one in the cubic
singlet companion: choose \(O\) by applying the exact transverse
isometry \(G^{-1/2}\mathsf T\) to this full mode basis. Hence these
first three \(z_\mu\) are actual coefficients of the original cochains
(5), not independently rescaled chord coordinates. Every color component
follows the same spatial matrix.

\subsection{3. Transport of the symmetry through the nonlinear tree
chart}\label{weak-coupling-part-16-transport-of-the-symmetry-through-the-nonlinear-tree-chart}

The original rooted tree is not assumed to be invariant under \(R\).
There is instead an exact induced chord map. Given the tree-\(I\)
representative for chord configuration \(Z\), apply (1) and then
recompute every root-to-vertex holonomy and every rooted chord
\(t_sU_ct_t^{-1}\). Call the result \(\beta_R(Z)\). It is a finite
ordered product of the original chord variables and their inverses. On
full physical functions it implements (2). The distinction matters: a
graph symmetry may move the chosen root, so the root-fixing gauge
subgroup by itself need not be preserved. Full vertex-gauge invariance,
proved in Section 1, is what makes the re-rooted representative
implement the original unitary. Together with the retained Haar/dilation
map \(\mathcal B_g\), the exact chart operator is \[
 \mathscr S_{R,g}=\mathcal B_g\mathscr S_R\mathcal B_g^* .
 \tag{9}
\] This formula defines its density and inverse, not only its argument
map. On a smaller identity chart its argument has expansion \[
 {1\over g}\log\beta_R(\exp(gx))
                 =D\beta_R(I)x+O(g|x|^2).
 \tag{10}
\] The product rule proves the linear part directly by replacing each
link with its additive cochain. Modulo the full gradient subspace it is
the oriented cochain transformation already computed in (6),(8). The
exact Gaussian Hilbert map keeps the determinant and Haar factors. At
\(g=0\) it therefore gives the orthogonal full transverse cochain
representation, with its induced action on the oscillator functions.

For clarity, convergence on every fixed polynomial Gaussian follows from
(10), not from assuming that the nonlinear chart is an orthogonal
transformation. On a fixed compact \(x\)-set its analytic argument and
density derivatives converge to the linear ones. Inside a sufficiently
small fixed \(y=gx\) chart, the induced map and its inverse have bounded
derivatives and Jacobians, with \(|\log\beta_R(\exp y)|\) bounded above
and below by positive constants times \(|y|\); its derivative at zero is
the invertible linear tree representative of the oriented cochain map
modulo gradients. The inverse function theorem supplies these local
bounds. Polynomial Gaussian tails on both sides are then bounded by a
polynomial times \(e^{-c|x|^2}\). Dominated convergence, followed by
removing the fixed compact \(x\)-set, proves strong convergence on those
vectors. Smooth cutoffs give the same assertion for the globally defined
chart vectors used in the companion's actual spectral projection.

In particular the representation on the first physical comparison
cluster is the symmetric square of the spatial representation (6),(8).
In the orthonormal six-vector basis, write \[
 D_\mu={\sigma_*\over2\sqrt6}
           (|z_\mu|^2-6/\sigma_*)\Phi,\qquad
 O_{\mu\nu}={\sigma_*\over2\sqrt3}(z_\mu\cdot z_\nu)\Phi,\quad\mu<\nu .
 \tag{11}
\] Spatial signed permutations send \(D_\mu\) to the corresponding
diagonal basis vector and \(O_{\mu\nu}\) to the corresponding
off-diagonal basis vector with the product of its two mode signs. No
normalization of an excitation energy or coupling is involved.

\subsection{4. Why the actual order-two matrix must commute with this
action}\label{weak-coupling-part-16-why-the-actual-order-two-matrix-must-commute-with-this-action}

The companion gives an exact unitary identification
\(J_g:\mathbb R^6\to\operatorname{ran}P_g\) with the actual first
physical spectral cluster and \[
 J_g^*(H_g-\mathcal E_g)J_g
          =\delta_*I+g^2\mathsf K+O_{L,a}(g^3),\qquad
                         \delta_*={2\sigma_*\over a}.
 \tag{12}
\] It is constructed from cutoff quasimodes, the actual spectral
projector \(P_g\), and the positive square root of their Gram matrix.
Each column converges under the chart map to (11). Since
\(\mathscr S_R\) commutes with (3), it preserves
\(\operatorname{ran}P_g\). The exact finite matrix \[
 D_g(R)=J_g^*\mathscr S_RJ_g
 \tag{13}
\] is orthogonal and commutes with the matrix on the left of (12). By
Section 3 and column convergence, \(D_g(R)\to D_0(R)\), where \(D_0(R)\)
is the explicitly signed permutation on (11). Subtract \(\delta_*I\)
from (12), divide the exact commutator by \(g^2\), and pass to the
limit. The remainder tends to zero in the finite matrix norm, yielding
\[
                     [\mathsf K,D_0(R)]=0
               \quad\hbox{for every signed permutation }R.
 \tag{14}
\] This proves invariance of the correction even though the selected
tree chart itself is not invariant. No order-\(g\) chart symmetry has
been omitted: the leading energy matrix is scalar and its first-order
correction is zero, so the limiting commutator is precisely (14).

\subsection{5. Three retained numbers and all six
corrections}\label{weak-coupling-part-16-three-retained-numbers-and-all-six-corrections}

In the order \(D_1,D_2,D_3,O_{12},O_{13},O_{23}\), the reflections (8)
fix every \(D_\mu\) and give three distinct nontrivial sign characters
to the three \(O_{\mu\nu}\). Consequently (14) forces all
diagonal/off-diagonal mixing entries and all entries between different
off-diagonal basis vectors to vanish. Coordinate permutations act
transitively on the three diagonal vectors, on their distinct pairs, and
on the three off-diagonal vectors, with possible signs in the last
action. Therefore (14) and real symmetry of \(\mathsf K\) imply exactly
\[
 \boxed{\mathsf K=
 \begin{pmatrix}
 d&o&o&0&0&0\\
 o&d&o&0&0&0\\
 o&o&d&0&0&0\\
 0&0&0&c&0&0\\
 0&0&0&0&c&0\\
 0&0&0&0&0&c
 \end{pmatrix}.}
 \tag{15}
\] Here \(d=\mathsf K_{D_1D_1}\), \(o=\mathsf K_{D_1D_2}\), and
\(c=\mathsf K_{O_{12}O_{12}}\) are the full finite sums from equations
(35)--(36), or equivalently (44), of the companion. The symbol \(c\) in
this note is a matrix coefficient, not a speed-of-light constant.

The vector \((1,1,1,0,0,0)/\sqrt3\) has eigenvalue \(d+2o\). The
two-dimensional subspace \(\{(x_1,x_2,x_3,0,0,0):x_1+x_2+x_3=0\}\) has
eigenvalue \(d-o\). The three-dimensional off-diagonal subspace has
eigenvalue \(c\). These statements follow by multiplying (15), without
invoking an undeclared representation-theoretic decomposition. Thus \[
 \boxed{\Delta_L(g,a)={2\sigma_*\over a}
             +g^2\min\{d+2o,\ d-o,\ c\}+O_{L,a}(g^3).}
 \tag{16}
\] All three entries carry their original factor \(a^{-1}\). The three
numbers still involve all spatial-mode interactions. Nothing here
decides their signs or large-\(L\) asymptotics. It does reduce the
original complete \(L=2\) numerical-certificate task to three specified
entries while proving all other entries and degeneracies exactly.

The conclusion uses the full symmetric open box and the global
Hamiltonian. Arbitrary local weights \(h\) need not possess these
spatial symmetries; their state coordinates are still projected against
the complete matrix (15), not declared symmetric.

\subsection{Dependencies}\label{weak-coupling-part-16-dependencies}

The included original finite-box source proves the tree/Haar isometry
and complete mode basis. The included cubic singlet, vacuum-energy and
first-gap note proves (12) with its controlled fixed-box residual and
exact all-mode expression for \(\mathsf K\). This note constructs the
original graph symmetries and their nonlinear-chart and limiting-cluster
intertwiners explicitly. It uses neither a conjectural continuum theory
nor an assertion that different coordinate presentations are unrelated.

\clearpage
\part{Current continuation: evaluate the complete physical gap matrix}
\section{Yang--Mills continuation: evaluate the retained physical gap
matrix}\label{weak-coupling-part-17-yangmills-continuation-evaluate-the-retained-physical-gap-matrix}

9 September 2026. This supersedes the earlier request to derive only the
first coupling coefficients. Supply this note together with
cubic\_singlet\_vacuum\_energy.md and
operator\_coupling\_expansion\_full\_graph.md, or with the current
complete TeX/Bib bundle containing those proofs. Also include
cubic\_box\_gap\_symmetry.md, which proves the exact reduction of this
calculation to three entries without removing any spatial modes.

\subsection{What is now calculated}\label{weak-coupling-part-17-what-is-now-calculated}

For the original full open-box SU(2) Hamiltonian \[
 H_g={2g^2\over a}\sum_eE_e+
           {1\over2g^2a}\sum_p(2-\operatorname{tr}U_p),
\] the exact tree/Haar map gives fully evaluated \(H_1,H_2\), including
both kinetic tensors, every ordered plaquette slot, and the Haar scalar
\(-\operatorname{tr}G/(2a)\). The complete cubic vacuum vector is \[
 H_1\Phi=\sum_{\mu<\nu<\rho}\Theta_{\mu\nu\rho}
                          \det(z_\mu,z_\nu,z_\rho)\Phi ,
\] with explicitly computed graph coefficients, raw determinant norm
\(48/(\sigma_\mu\sigma_\nu\sigma_\rho)\), and exact excitation
\((\sigma_\mu+\sigma_\nu+\sigma_\rho)/a\). Every spatial mode remains.

The full first physical oscillator cluster consists of six color-singlet
quadratics in the three lowest spatial modes. Writing their orthonormal
vectors as \(S_I\), the actual fixed-box first gap satisfies \[
 \Delta_L(g,a)={2\sigma_*\over a}
                 +g^2\lambda_{\min}(\mathsf K)+O_{L,a}(g^3).
\] The source proves the cutoff residual, actual spectral projection and
weighted graph estimate; it does not assume differentiability at zero
coupling. The finite coefficient matrix is \[
 \mathsf K_{IJ}
 =\langle S_I,H_2S_J\rangle-e_2\delta_{IJ}
 -\sum_{|n|=3,5}{c_{I,n}c_{J,n}\over\omega_n-2\sigma_*/a},
 \quad c_{I,n}=\langle h_n\Phi,H_1S_I\rangle .
\] Its coefficients are specified by finite polynomials and complete
Gaussian pairings in equations (29)--(36), not left as unknown matrix
elements. Degree-one intermediate states vanish by color invariance. All
remaining denominators are strictly positive.

Equations (39)--(44) prove an exact cancellation of spectator vacuum
terms. In the finite normal-ordered polynomial, decompose \(H_n\) into
terms acting only outside the three lowest modes and terms touching at
least one lowest mode. The outside-only cubic and quartic contributions
cancel from \(\mathsf K\), with their full energy denominators and the
Haar scalar retained on both sides. The surviving cubic vectors have at
most two outside quanta. This is an exact computational reduction proved
from occupation numbers, not deletion of high spatial modes.

One further usable estimate is \[
 \|D_{h,g}u\|\le\sqrt2\|h\|_\infty
 [\,\|H_gu\|^2+6|P|\|u\|^2/a^2\,]^{1/2}.
\] It follows by integration by parts from \(KW=3W-6|P|\) and the
original \(\kappa b=a^{-2}\), so it does not lose two powers of \(g\).

\subsection{One bounded, demanding next
calculation}\label{weak-coupling-part-17-one-bounded-demanding-next-calculation}

Evaluate the complete matrix \(a\mathsf K\) on the original \(L=2\) box,
retaining all 176 spatial modes, all original edges and faces, and all
color contractions. Use the spectator cancellation to avoid enumerating
irrelevant occupation lists. The cubic-box symmetry proof now gives
exactly a \(3\times3\) diagonal-state block with entries \(d\) on its
diagonal and \(o\) off its diagonal, and an off-diagonal-state block
\(cI_3\), with no mixing. Thus only \(ad,ao,ac\) need evaluation: the
six corrections are \(a(d+2o)\) once, \(a(d-o)\) twice, and \(ac\) three
times. The proof constructs the original signed graph symmetries and
their nonlinear-chart/physical-band correspondence; do not assume the
particular rooted tree was itself invariant. Give certified rational
intervals for its six eigenvalues, or an exact algebraic representation
and a reproducible integer outward-enclosure checker. Include the actual
tree, full transverse basis, mode ordering, matrix entries, and the map
back to the physical correction \(g^2a^{-1}\lambda_{\min}(a\mathsf K)\).
Multiplying by \(a\) here records the exact scale factor; it does not
set the physical spacing or energy to one.

Do not substitute a smaller graph, one plaquette, a lowest-mode
truncation, an Abelian color sector, or floating-point eigenvalues
presented as exact signs. Do not repeat only the nonzero cubic vertex
certificate: the target is the complete physical excitation correction,
including the moving vacuum and quartic terms.

Then inspect the resulting contraction formula for the term whose
dependence on \(L\) would govern a physical-scale limit. Any asymptotic
claim must be derived from that formula with an explicit error. The sign
at \(L=2\), by itself, proves neither a mass gap nor its failure in an
interacting continuum. If the finite matrix is all that can be
completed, deliver its complete checked calculation without replacing
the remaining question by an assumed bound.

\subsection{Existing checks and scope}\label{weak-coupling-part-17-existing-checks-and-scope}

The supplied check\_cubic\_singlet\_vacuum\_energy.py currently verifies
85 exact identities: arbitrary ordered SU(2) face coefficients through
degree four; symbolic quartic covariance; cubic kinetic contraction; the
three-mode eigenvalue and raw norm; exact rational kinetic Wick tests;
and all six first-cluster norms, cross inner products and excitations.
It does not evaluate the full \(L=2\) matrix and is not a continuum
certificate.

The separate check\_cubic\_box\_gap\_symmetry.py passes 99 exact checks:
all 48 signed-coordinate actions on the lowest modes and their six
quadratic states, every commutator with the displayed matrix, the full
symmetric commutant dimension three, and the complete characteristic
polynomial. This verifies the finite matrix algebra, not the pending
three coefficient values.

\clearpage
\part{Non-Abelian fluid encoding: exact inverse, curvature and current}
\section{Non-Abelian fluid encoding: exact inverse, curvature and
current}\label{weak-coupling-part-18-non-abelian-fluid-encoding-exact-inverse-curvature-and-current}

9 September 2026. Independent calculation of the non-Abelian connection
proposed in the interacting tensor-band continuation. This note proves
identities for the original fluid field, rather than replacing it with a
different profile. The identities do not themselves construct a quantum
vacuum or a continuum spectral measure.

\subsection{1. Spaces, coefficients and
inverse}\label{weak-coupling-part-18-spaces-coefficients-and-inverse}

Let \(I\) be an open time interval and let
\(u\in C^\infty(I\times\mathbb R^3;\mathbb R^3)\) satisfy
\(\sum_i\partial_i u_i=0\). Spatial support is assumed compact whenever
an integral over all of \(\mathbb R^3\) is used. All pointwise
identities below require no support assumption. Retain \(c>0\), \(g>0\),
\(\lambda\in\mathbb R\setminus\{0\}\), \(X^0=cs\), \(X^i=x^i\), and the
metric \(\eta=\operatorname{diag}(-1,1,1,1)\).

For the Pauli matrices \(\sigma_a\), put \(T_a=-i\sigma_a/2\). Then \[
 [T_a,T_b]=\sum_d\varepsilon_{abd}T_d,\qquad
 -2\operatorname{tr}(T_aT_b)=\delta_{ab}.
 \tag{1}
\] Consequently the coefficient vector of \([v\cdot T,w\cdot T]\) is
\(v\times w\). We use this specified identification of the color basis
with the oriented Cartesian basis \((e_1,e_2,e_3)\) throughout.

Define a connection by its coefficient vectors: \[
 \mathbf A_0=0,\qquad
 \mathbf A_i=\lambda e_i\times u,\qquad
 A_\mu=\mathbf A_\mu\cdot T.
 \tag{2}
\] This is an injective linear map from the stated fluid fields to a
specified subspace of smooth connections. Its inverse on that image is
\[
 u=-\frac1{2\lambda}\sum_i e_i\times\mathbf A_i.
 \tag{3}
\] Indeed, \(\sum_i e_i\times(e_i\times u)=\sum_i(u_i e_i-u)=-2u\). This
is an inverse in the displayed representative and frame, not an
assertion that the same formula descends to all gauge-equivalence
classes.

There is also an exact correspondence with the earlier Cartan encoding
\(C_0=0,\ C_i=\lambda u_iT_3\). Recover
\(u_i=-2\operatorname{tr}(T_3C_i)/\lambda\), then use (2). Conversely
apply (3), then set \(C_i=\lambda u_iT_3\). Both compositions are the
identity on their respective specified images. The two encodings are
thus explicitly related without identifying their curvature or their
dynamics.

\subsection{2. Full curvature}\label{weak-coupling-part-18-full-curvature}

Use \[
 F_{\mu\nu}=\partial_\mu A_\nu-\partial_\nu A_\mu+[A_\mu,A_\nu],
 \quad
 \mathbf E_i=\mathbf F_{0i},\quad
 \mathbf B_i=\tfrac12\sum_{j,k}\varepsilon_{ijk}\mathbf F_{jk}.
 \tag{4}
\] The temporal term and the unreduced magnetic components are \[
 \mathbf E_i=\frac{\lambda}{c}e_i\times u_s,\qquad
 B_i^a=\lambda(\partial_a u_i-\delta_{ia}\operatorname{div}u)
                    +\lambda^2u_i u_a.
 \tag{5}
\] To verify the second formula, expand (4) and use \[
 \sum_k\varepsilon_{ijk}\varepsilon_{akb}
       =\delta_{ib}\delta_{ja}-\delta_{ia}\delta_{jb},
 \quad
 ((e_j\times u)\times(e_k\times u))_a
       =\sum_b\varepsilon_{jkb}u_bu_a,
 \quad
 \tfrac12\sum_{j,k}\varepsilon_{ijk}\varepsilon_{jkb}=\delta_{ib}.
 \tag{6}
\] For the original incompressible field this gives exactly \[
 \boxed{\,B_i^a=\lambda\partial_a u_i+\lambda^2u_i u_a\,}.
 \tag{7}
\] In particular the commutator is retained: \[
 [A_i,A_j]=\lambda^2\sum_k\varepsilon_{ijk}u_k\,(u\cdot T).
 \tag{8}
\] If \(u\ne0\) at a point, some \(u_k\ne0\), and (8) is nonzero for the
complementary oriented pair. This encoding does not take values in one
fixed commuting color line.

\subsection{3. Pointwise and integrated curvature
norms}\label{weak-coupling-part-18-pointwise-and-integrated-curvature-norms}

The coefficient norm is exactly the norm from (1). Expanding all terms
in (7) gives \[
 \begin{aligned}
 |\mathbf B|^2
 &=\lambda^2\sum_{i,a}(\partial_a u_i)^2
   +2\lambda^3\sum_{i,a}u_i u_a\partial_a u_i+\lambda^4|u|^4\\
 &=\lambda^2|\nabla u|^2+
   \lambda^3\operatorname{div}(|u|^2u)+\lambda^4|u|^4 .
 \end{aligned}
 \tag{9}
\] The last equality uses \(\operatorname{div}(|u|^2u)=
 2\sum_{i,a}u_i u_a\partial_a u_i+|u|^2\operatorname{div}u\). Compact
spatial support therefore implies \[
 \boxed{\ \int|\mathbf B|^2dx
  =\lambda^2\|\nabla u\|_2^2+\lambda^4\|u\|_4^4.\ }
 \tag{10}
\] All boundary terms vanish because their integrands are compactly
supported. Also \[
 \int|\operatorname{curl}u|^2dx
 =\int\sum_{i,a}(\partial_a u_i)^2dx
   -\int\sum_{i,a}\partial_a u_i\,\partial_i u_a\,dx
 =\|\nabla u\|_2^2,
 \tag{11}
\] because integration by parts turns the last sum into
\(\int(\operatorname{div}u)^2dx=0\).

There is a pointwise estimate which does not discard the derivative term
in (9). Regard \(B=(B_i^a)\) as a real \(3\)-by-\(3\) matrix. Its trace
is \(\lambda^2|u|^2\). Frobenius Cauchy--Schwarz gives \[
 (\lambda^2|u|^2)^2=|\langle B,I_3\rangle_F|^2
 \le 3|B|^2,\qquad
 \boxed{\,|B|^2\ge \lambda^4|u|^4/3\,}.
 \tag{12}
\] Equality holds precisely when \[
 \partial_a u_i=
 \lambda\bigl(|u|^2\delta_{ia}/3-u_i u_a\bigr).
 \tag{13}
\] The electric identity is \[
 \sum_i|\mathbf E_i|^2
 =\frac{\lambda^2}{c^2}\sum_i(|u_s|^2-(u_s)_i^2)
 =\frac{2\lambda^2}{c^2}|u_s|^2.
 \tag{14}
\] Thus, if the specified classical energy convention is
\(\mathcal E=(2g^2)^{-1}\int(|\mathbf E|^2+|\mathbf B|^2)dx\), its value
on (2) is exactly \[
 \mathcal E=
 \frac{\lambda^2}{g^2c^2}\|u_s\|_2^2+
 \frac{\lambda^2}{2g^2}\|\operatorname{curl}u\|_2^2+
 \frac{\lambda^4}{2g^2}\|u\|_4^4.
 \tag{15}
\] The factor \(2\) in (14) and the factor \(1/3\) in (12) depend on the
displayed dimensions and norm; neither is suppressed.

\subsection{4. Full current and its Gauss
component}\label{weak-coupling-part-18-full-current-and-its-gauss-component}

Define the induced current, without imposing that it vanish, by \[
 D^\mu F_{\mu\nu}=g^2j_\nu,\qquad
 D_\mu V=\partial_\mu V+[A_\mu,V].
 \tag{16}
\] Write \(\boldsymbol\omega=\operatorname{curl}u\) and
\(C=(u\cdot\nabla)u\). Direct substitution gives \[
 \boxed{\ g^2\mathbf j_0
 =-\frac{\lambda}{c}\operatorname{curl}u_s
  -\frac{\lambda^2}{c}u\times u_s.\ }
 \tag{17}
\] Indeed \(F_{i0}=-F_{0i}\), so the derivative part is
\(-(\lambda/c)\sum_i e_i\times\partial_i u_s\). The commutator part uses
the exact identity \(\sum_i(e_i\times u)\times(e_i\times v)=u\times v\).
Equation (17) is the lower-index component; \(j^0=-j_0\).

For each \(i\), the complete spatial expression is \[
 \begin{aligned}
 g^2\mathbf j_i={}&-\frac{\lambda}{c^2}e_i\times u_{ss}
                  -\lambda\nabla\omega_i\\
 &-\lambda^2\left[
       2\omega_i u+
       \sum_{j,k}\varepsilon_{ijk}u_k\partial_j u+
       e_i\times C\right]
       -\lambda^3|u|^2(e_i\times u).
 \end{aligned}
 \tag{18}
\] Here \(\nabla\omega_i\) is the color coefficient vector whose
component \(a\) is \(\partial_a\omega_i\). For a complete sign check,
first use \(F_{ji}=-\sum_k\varepsilon_{ijk}B_k\) to obtain \[
 g^2\mathbf j_i=
 -\frac{\lambda}{c^2}e_i\times u_{ss}
 -\sum_{j,k}\varepsilon_{ijk}
       [\,\partial_j\mathbf B_k+
               \lambda(e_j\times u)\times\mathbf B_k\,].
 \tag{19}
\] Then (7) and the vector triple-product identity give \[
 \begin{aligned}
 \partial_j\mathbf B_k+\lambda(e_j\times u)\times\mathbf B_k
 ={}&\lambda\nabla\partial_j u_k\\
 &+\lambda^2[\,2(\partial_j u_k)u+
                u_k\partial_j u-e_j C_k\,]\\
 &+\lambda^3[\,u_j u_k u-|u|^2u_k e_j\,].
 \end{aligned}
 \tag{20}
\] Contract (20) using \(\sum_{j,k}\varepsilon_{ijk}u_j u_k=0\) and
\(\sum_{j,k}\varepsilon_{ijk}e_jv_k=-e_i\times v\). This proves every
term of (18). In particular the coefficient \(2\) comes from two
different terms, the derivative of the quadratic magnetic part and the
commutator with its linear part.

Covariant conservation also holds with these same conventions: \[
 g^2D^\nu j_\nu
 =D^\nu D^\mu F_{\mu\nu}
 =\tfrac12[D^\nu,D^\mu]F_{\mu\nu}
 =\tfrac12[F^{\nu\mu},F_{\mu\nu}]=0.
 \tag{21}
\] The second equality follows by exchanging the two summed indices and
using antisymmetry. In the specified diagonal metric each term in the
last sum is the commutator of a scalar multiple of \(F_{\mu\nu}\) with
itself, and vanishes.

\subsection{5. Substitution of the original viscous fluid
equation}\label{weak-coupling-part-18-substitution-of-the-original-viscous-fluid-equation}

Suppose the same field satisfies the smooth forced incompressible
Navier--Stokes equation before its terminal time: \[
 u_s+(u\cdot\nabla)u-\nu\Delta u+\nabla p=f,\qquad \nu>0.
 \tag{22}
\] Define, with every original term retained, \[
 N=\nu\Delta u-(u\cdot\nabla)u-\nabla p+f.
 \tag{23}
\] Then \(u_s=N\) and direct differentiation gives \[
 u_{ss}=\nu\Delta N-(N\cdot\nabla)u-(u\cdot\nabla)N
                       -\nabla p_s+f_s.
 \tag{24}
\] Putting (23)--(24) into (17)--(18) is the exact fluid-to-current map;
it keeps viscosity, pressure, both differentiated transport terms,
forcing, color commutators and the nonzero Gauss component. No equation
in (22) sets (17) or (18) to zero.

The released positive-viscosity construction is an eligible input on
each smooth preterminal interval: see OpenAI, \emph{Finite time blowup
for Navier--Stokes},
\href{https://openai.com/index/navier-stokes-solution/}{paper and formal
source}. Equations (10), (14), and (17)--(24) specify what must be
evaluated on that actual field. They do not identify the fluid's smooth
external force with a smooth gauge current at a singular endpoint.

The current is derived from the connection, not removed. In particular
the non-Abelian quartic magnetic term in (10) is a proved feature of the
map; a source-free quantum spectrum is not a consequence of that
identity alone.

\subsection{6. An exact right inverse for non-Abelian curvature
samples}\label{weak-coupling-part-18-an-exact-right-inverse-for-non-abelian-curvature-samples}

The quartic term also gives a direct construction, independent of the
singular fluid profile. Let \(\mathsf E\) be any finite set of original
cubic-grid edges of physical spacing \(a>0\), and let \(p_e\) be their
distinct midpoints. Given arbitrary real weights \(h_e\ge0\), we now
construct a smooth compactly supported incompressible field \(U_h\) such
that the connection (2) satisfies \[
 \sum_i|\mathbf B_i[U_h](p_e)|^2=h_e\qquad(e\in\mathsf E).
 \tag{25}
\] This is a right inverse of the explicitly defined finite sampling
map, not a substitution for the released blowup solution.

Set \(\rho=a/8\). Distinct edge midpoints are separated by at least
\(a/\sqrt2\): in lattice units each has precisely one half-integer
coordinate; either their half-integer positions coincide and a nonzero
difference has size at least one, or two coordinates differ by at least
\(1/2\). Thus the closed radius-\(\rho\) balls are disjoint.

Define \(b(t)=e^{-1/t}\) for \(t>0\), and \(b(t)=0\) for \(t\le0\). All
derivatives of its positive branch tend to zero at zero, because each is
a polynomial in \(1/t\) times \(e^{-1/t}\); hence \(b\) is smooth. For
\(r_e=x-p_e\), put \[
 \chi_e(x)=
 \frac{b(1-|r_e|^2/\rho^2)}
 {b(1-|r_e|^2/\rho^2)+b(|r_e|^2/\rho^2-1/4)},\qquad
 v_e=\frac{h_e^{1/4}}{|\lambda|}e_3 .
 \tag{26}
\] The denominator is positive everywhere: both terms could vanish only
if \(|r_e|^2/\rho^2\ge1\) and \(|r_e|^2/\rho^2\le1/4\) simultaneously.
The cutoff equals one for \(|r_e|\le\rho/2\) and zero for
\(|r_e|\ge\rho\).

Set \[
 V_h(x)=\frac12\sum_{e\in\mathsf E}\chi_e(x)(v_e\times r_e),
 \qquad U_h=\operatorname{curl}V_h .
 \tag{27}
\] The sum is finite and smooth with compact support. Antisymmetry of
the curl and commutation of second derivatives give
\(\operatorname{div}U_h=0\). In the radius-\(\rho/2\) ball at \(p_e\)
only that summand remains, and its cutoff is one. Componentwise, \[
 \bigl(\operatorname{curl}(v\times r)\bigr)_i
 =\sum_{j,k,b,d}\varepsilon_{ijk}\varepsilon_{kbd}
          v_b\partial_j r_d
 =\sum_{j,k,b}\varepsilon_{ijk}\varepsilon_{kbj}v_b=2v_i.
 \tag{28}
\] Therefore \(U_h=v_e\), \(\nabla U_h=0\) throughout that inner ball.
By (7), \[
 B_i^a[U_h](p_e)=\lambda^2(v_e)_i(v_e)_a,\qquad
 \sum_{i,a}|B_i^a[U_h](p_e)|^2
       =\lambda^4|v_e|^4=h_e.
 \tag{29}
\] This proves (25). If \(h_e>0\), the sampled magnetic field is
entirely the retained commutator term, not a derivative term in a
commuting encoding. The fourth root is the stated set-map on nonnegative
weights; differentiability in \(h\) at a zero weight is not asserted.

For completeness these fields occur in explicit regular forced
Navier--Stokes flows for every prescribed \(\nu>0\). Let
\(\theta\in C_c^\infty((0,\infty);\mathbb R)\) equal one on a nonempty
time interval and put \[
 u(s,x)=\theta(s)U_h(x),\quad p(s,x)=0,\quad
 f(s,x)=\theta'(s)U_h(x)+
       \theta(s)^2(U_h\cdot\nabla)U_h(x)
       -\nu\theta(s)\Delta U_h(x).
 \tag{30}
\] Direct substitution proves (22), with zero initial velocity and
smooth compact spacetime forcing. The velocity and force are defined
through the time plateau and thereafter; no endpoint singularity is
introduced. On the plateau, at each sampled point, (17)--(18) give \[
 \mathbf j_0(p_e)=0,\qquad
 g^2\mathbf j_i(p_e)=
       -\lambda^3|v_e|^2(e_i\times v_e).
 \tag{31}
\] Thus the gauge source remains explicit even where the fluid is
locally constant. Formula (30) depends on the prescribed weights and the
grid spacing. Neither fixed forcing across regulators nor a source-free
field equation is claimed.

\end{document}
